% Integration — Pure Mathematics with Statistics Upper Sixth — Cours signé OnBuch+
% Official MINESEC syllabus. Compiler avec : tectonic PMS04-integration.tex
\newcommand{\DOCMATIERE}{Pure Mathematics with Statistics}
\newcommand{\DOCNIVEAU}{Upper Sixth}
\newcommand{\DOCMODULE}{Upper Sixth — Pure mathematics with statistics}
\newcommand{\DOCLECON}{4}
\newcommand{\DOCTITRE}{Integration}
% OnBuch+ — préambule « cours signés » ENRICHI (compilé avec tectonic / XeLaTeX)
% Système visuel « Pop » d'OnBuch+ : fond crème (jamais blanc), bordures encre
% épaisses, ombres « dures » décalées, titres Archivo Black en capitales.
% Version MATHÉMATIQUES : graphes (pgfplots/tikz), boîtes pédagogiques
% (définition, propriété, méthode, attention, le savais-tu, exercice résolu,
% exercice, TP, bilan), page de garde et signature OnBuch+ ancrée en bas.
%
% Macros attendues dans main.tex AVANT \input{preamble} :
%   \DOCMATIERE  \DOCNIVEAU  \DOCMODULE  \DOCLECON (numéro)  \DOCTITRE
\documentclass[11pt]{article}

\usepackage{fontspec}
\usepackage[english]{babel}
\usepackage[a4paper,margin=2cm,top=3.4cm,bottom=2.8cm,headheight=1.4cm,footskip=1.3cm]{geometry}
\usepackage{amsmath,amssymb,amsfonts}
\usepackage{mathtools} % \xmapsto, \coloneqq... souvent utilisés par les modèles
\usepackage{cancel} % simplifications barrées (bilans de forces)
% \abs{z} (module, valeur absolue) et \norm{u} (norme), utilisés par les modèles
\providecommand{\abs}{}\let\abs\relax\DeclarePairedDelimiter{\abs}{\lvert}{\rvert}
\providecommand{\norm}{}\let\norm\relax\DeclarePairedDelimiter{\norm}{\lVert}{\rVert}
\usepackage[table]{xcolor} % [table] : fournit aussi \rowcolors (alternance de lignes)
\usepackage{pagecolor}
\usepackage{tikz}
\usetikzlibrary{arrows.meta,positioning,calc,shapes.geometric,shapes.misc,shapes.arrows,decorations.pathmorphing,decorations.pathreplacing,decorations.markings,patterns,shadows,mindmap,trees,fit,backgrounds,matrix,intersections,angles,quotes}
\usepackage{pgfplots}
\usepackage[version=4]{mhchem} % équations nucléaires \ce{^{235}_{92}U}
\usepackage[european,siunitx]{circuitikz} % schémas électriques
\pgfplotsset{compat=1.18}
\usepgfplotslibrary{fillbetween,groupplots} % aires entre courbes (calcul intégral)
\usepackage{enumitem}
\usepackage{fancyhdr}
% [most] charge toutes les bibliothèques tcolorbox (listings, minted, raster,
% documentation...) et gonfle inutilement la mémoire TeX sur un document long
% (~300 boîtes) : on ne charge que ce qui est réellement utilisé.
\usepackage[skins,breakable]{tcolorbox}
\usepackage{graphicx}
\usepackage{colortbl}
\usepackage{booktabs}
\usepackage{tabularx}
\usepackage{diagbox} % case d'en-tête diagonale des tableaux à double entrée
% Colonnes extensibles centrée / gauche / droite (C, L, R), souvent utilisées par les modèles
\newcolumntype{C}{>{\centering\arraybackslash}X}
\newcolumntype{L}{>{\raggedright\arraybackslash}X}
\newcolumntype{R}{>{\raggedleft\arraybackslash}X}
\usepackage{array}
\usepackage{multicol}
\usepackage{siunitx}
% parse-numbers=false : accepte aussi \SI{2,5 \times 10^{-3}}{...}
\sisetup{locale=UK,output-decimal-marker={.},group-minimum-digits=4,parse-numbers=false}
\DeclareSIUnit\ohm{\ensuremath{\symup{\Omega}}} % U+2126 absent de la police
\DeclareSIUnit\division{div} % réglages d'oscilloscope (ms/div, V/div)
\DeclareSIUnit\tr{tr} % tours (vitesse de rotation, ex. \SI{1500}{\tr\per\minute})
\usepackage{qrcode}
% \degree : commande fréquemment utilisée par les modèles pour un angle ou une
% température en dehors de \SI{}{} (non fournie par les paquets chargés).
\providecommand{\degree}{^{\circ}}
% \qed (fin de démonstration), utilisé par les modèles sans amsthm
\providecommand{\qed}{\hfill$\square$}
% \barre : double barre des valeurs interdites dans un tableau de variations (tkz-tab)
\providecommand{\barre}{\ensuremath{\|}}
% environnement proof (amsthm non chargé) utilisé par les modèles
\ifdefined\proof\else\newenvironment{proof}[1][Proof]{\par\noindent\textit{#1.}\ }{\qed\par}\fi
\usepackage{lastpage}
\usepackage{hyperref}
\hypersetup{hidelinks}

% ============================================================
% Palette « Pop » OnBuch+ — mêmes hex que la classe P (plus_ui.dart)
% ============================================================
\definecolor{poporange}{HTML}{F59321}
\definecolor{popdark}{HTML}{E07A0C}
\definecolor{popcream}{HTML}{FFF6E8}
\definecolor{popink}{HTML}{1C1714}
\definecolor{poppurple}{HTML}{7A5AE0}
\definecolor{popgreen}{HTML}{2E9E6A}
\definecolor{poppink}{HTML}{E0457B}
\definecolor{popblue}{HTML}{2F7DE1}
\definecolor{popgold}{HTML}{C98A12}
\definecolor{popmuted}{HTML}{8A7F76}
\definecolor{popline}{HTML}{EADFD2}
\definecolor{popgray}{HTML}{8A7F76}
\colorlet{popgrey}{popgray}
% Alias tolérants (le modèle nomme parfois les couleurs différemment)
\colorlet{popwhite}{white}
\colorlet{popblack}{popink}
\colorlet{popred}{poppink}
\colorlet{popyellow}{popgold}
% Teintes claires pour fonds de tableaux / zones
\colorlet{popblueL}{popblue!10!white}
\colorlet{poporangeL}{poporange!14!white}
\colorlet{popgreenL}{popgreen!12!white}
\colorlet{poppurpleL}{poppurple!12!white}
\colorlet{poppinkL}{poppink!12!white}
\colorlet{popgoldL}{popgold!14!white}

\pagecolor{popcream}

% ============================================================
% Typographie
% ============================================================
\defaultfontfeatures{Ligatures=TeX}
\setmainfont{PlusJakartaSans-Medium.ttf}[Path=../../../fonts/,BoldFont=PlusJakartaSans-Bold.ttf]
% Maths en sans-serif (Fira Math) avec chiffres et lettres droites de Plus
% Jakarta, pour que les formules se fondent dans le texte.
\usepackage[math-style=ISO,bold-style=ISO]{unicode-math}
\setmathfont{FiraMath-Regular.otf}[Path=../../../fonts/]
\setmathfont{PlusJakartaSans-Medium.ttf}[Path=../../../fonts/,range={up/num,up/{Latin,latin},\mathup}]
% (icomma retiré : décimales avec point en anglais)
% Caractères mathématiques Unicode absents des polices de texte (Plus Jakarta,
% Archivo Black) : rendus par leur équivalent mathématique, en texte comme en
% formule — sinon ils s'affichent en carré vide (ex. « Arithmétique dans ℤ »).
\usepackage{newunicodechar}
\newunicodechar{ℝ}{\ensuremath{\mathbb{R}}}
\newunicodechar{ℤ}{\ensuremath{\mathbb{Z}}}
\newunicodechar{ℂ}{\ensuremath{\mathbb{C}}}
\newunicodechar{ℕ}{\ensuremath{\mathbb{N}}}
\newunicodechar{ℚ}{\ensuremath{\mathbb{Q}}}
\newunicodechar{𝔻}{\ensuremath{\mathbb{D}}}
\newunicodechar{α}{\ensuremath{\alpha}}
\newunicodechar{β}{\ensuremath{\beta}}
\newunicodechar{γ}{\ensuremath{\gamma}}
\newunicodechar{δ}{\ensuremath{\delta}}
\newunicodechar{ε}{\ensuremath{\varepsilon}}
\newunicodechar{θ}{\ensuremath{\theta}}
\newunicodechar{λ}{\ensuremath{\lambda}}
\newunicodechar{μ}{\ensuremath{\mu}}
\newunicodechar{σ}{\ensuremath{\sigma}}
\newunicodechar{τ}{\ensuremath{\tau}}
\newunicodechar{φ}{\ensuremath{\varphi}}
\newunicodechar{ψ}{\ensuremath{\psi}}
\newunicodechar{ω}{\ensuremath{\omega}}
\newunicodechar{Σ}{\ensuremath{\Sigma}}
% majuscules produites par \MakeUppercase dans les titres (α-aminés -> Α)
\newunicodechar{Α}{\ensuremath{\alpha}}
\newunicodechar{Β}{\ensuremath{\beta}}
\newunicodechar{Γ}{\ensuremath{\Gamma}}
\newunicodechar{Δ}{\ensuremath{\Delta}}
\newunicodechar{Ε}{\ensuremath{\varepsilon}}
\newunicodechar{Θ}{\ensuremath{\Theta}}
\newunicodechar{Λ}{\ensuremath{\Lambda}}
\newunicodechar{Μ}{\ensuremath{\mu}}
\newunicodechar{Τ}{\ensuremath{\tau}}
\newunicodechar{Φ}{\ensuremath{\Phi}}
\newunicodechar{Ψ}{\ensuremath{\Psi}}
\newunicodechar{Ω}{\ensuremath{\Omega}}
\newunicodechar{↦}{\ensuremath{\mapsto}}
\newunicodechar{∈}{\ensuremath{\in}}
\newunicodechar{∉}{\ensuremath{\notin}}
\newunicodechar{∧}{\ensuremath{\wedge}}
\newunicodechar{⇔}{\ensuremath{\Leftrightarrow}}
\newunicodechar{⟺}{\ensuremath{\Longleftrightarrow}}
\newunicodechar{⇒}{\ensuremath{\Rightarrow}}
\newunicodechar{⟹}{\ensuremath{\Longrightarrow}}
\newunicodechar{∘}{\ensuremath{\circ}}
\newunicodechar{∪}{\ensuremath{\cup}}
\newunicodechar{∩}{\ensuremath{\cap}}
\newunicodechar{≡}{\ensuremath{\equiv}}
\newunicodechar{⊂}{\ensuremath{\subset}}
\newunicodechar{⊥}{\ensuremath{\perp}}
\newunicodechar{∃}{\ensuremath{\exists}}
\newunicodechar{∀}{\ensuremath{\forall}}
\newunicodechar{∛}{\ensuremath{\sqrt[3]{\,}}}
\newunicodechar{∜}{\ensuremath{\sqrt[4]{\,}}}
\newunicodechar{⃗}{\ensuremath{{}^{\to}}}
\newunicodechar{ⁿ}{\ifmmode^{n}\else\textsuperscript{\ensuremath{n}}\fi}
\newunicodechar{ˣ}{\ifmmode^{x}\else\textsuperscript{\ensuremath{x}}\fi}
\newunicodechar{ᵇ}{\ifmmode^{b}\else\textsuperscript{\ensuremath{b}}\fi}
\newunicodechar{ʳ}{\ifmmode^{r}\else\textsuperscript{\ensuremath{r}}\fi}
\newunicodechar{ᵘ}{\ifmmode^{u}\else\textsuperscript{\ensuremath{u}}\fi}
\newunicodechar{ʸ}{\ifmmode^{y}\else\textsuperscript{\ensuremath{y}}\fi}
\newunicodechar{ᵃ}{\ifmmode^{a}\else\textsuperscript{\ensuremath{a}}\fi}
\newunicodechar{ᵉ}{\ifmmode^{e}\else\textsuperscript{\ensuremath{e}}\fi}
\newunicodechar{ᵏ}{\ifmmode^{k}\else\textsuperscript{\ensuremath{k}}\fi}
\newunicodechar{ᵖ}{\ifmmode^{p}\else\textsuperscript{\ensuremath{p}}\fi}
\newunicodechar{ᴺ}{\ifmmode^{N}\else\textsuperscript{\ensuremath{N}}\fi}
\newunicodechar{ᶜ}{\ifmmode^{c}\else\textsuperscript{\ensuremath{c}}\fi}
\newunicodechar{ʰ}{\ifmmode^{h}\else\textsuperscript{\ensuremath{h}}\fi}
\newunicodechar{ᵠ}{\ifmmode^{q}\else\textsuperscript{\ensuremath{q}}\fi}
\newunicodechar{ₙ}{\ifmmode_{n}\else\textsubscript{\ensuremath{n}}\fi}
\newunicodechar{ₐ}{\ifmmode_{a}\else\textsubscript{\ensuremath{a}}\fi}
\newunicodechar{ᵢ}{\ifmmode_{i}\else\textsubscript{\ensuremath{i}}\fi}
\newunicodechar{ᵣ}{\ifmmode_{r}\else\textsubscript{\ensuremath{r}}\fi}
\newunicodechar{ₖ}{\ifmmode_{k}\else\textsubscript{\ensuremath{k}}\fi}
\newunicodechar{ᵦ}{\ifmmode_{\beta}\else\textsubscript{\ensuremath{\beta}}\fi}
\newunicodechar{ₓ}{\ifmmode_{x}\else\textsubscript{\ensuremath{x}}\fi}
\newfontfamily\popdisplay{ArchivoBlack-Regular.ttf}[Path=../../../fonts/]
\newfontfamily\poplogo{SpaceGrotesk-Bold.ttf}[Path=../../../fonts/]
\newfontfamily\popsemi{PlusJakartaSans-SemiBold.ttf}[Path=../../../fonts/]
\newfontfamily\popxbold{PlusJakartaSans-ExtraBold.ttf}[Path=../../../fonts/]
\newfontfamily\popxboldls{PlusJakartaSans-ExtraBold.ttf}[Path=../../../fonts/,LetterSpace=3.0]

\setlist[enumerate]{leftmargin=1.7em,itemsep=5pt,topsep=4pt}
\setlist[itemize]{leftmargin=1.7em,itemsep=4pt,topsep=4pt,label={\color{poporange}\textbullet}}
\setlength{\parindent}{0pt}
\setlength{\parskip}{5pt}
\renewcommand{\arraystretch}{1.25}

% pgfplots : style Pop par défaut
\pgfplotsset{
  every axis/.append style={axis line style={popink,line width=1pt},tick style={popink},
    grid style={popline,line width=0.5pt},label style={font=\small},tick label style={font=\footnotesize}},
  popcurve/.style={poporange,line width=1.8pt,no markers,smooth},
}
% Même style disponible dans un \draw TikZ simple (hors pgfplots)
\tikzset{popcurve/.style={poporange,line width=1.8pt}}
\tikzset{popgrid/.style={popline,very thin}}
\colorlet{popcurve}{poporange} % le nom du style est parfois utilisé comme couleur

% ============================================================
% Pill / squiggle
% ============================================================
\newcommand{\poppill}[3]{%
  \tikz[baseline=-0.55ex]\node[draw=popink,line width=1.2pt,rounded corners=2.6pt,%
    fill=#2,inner xsep=6.5pt,inner ysep=3pt]{%
    \popxboldls\fontsize{7}{7}\selectfont\color{#3}\MakeUppercase{#1}};%
}
\newcommand{\popsquiggle}[1]{%
  \tikz[baseline=-0.4ex]{
    \draw[#1,line width=2.6pt,line cap=round]
      plot [smooth] coordinates {(0,0.09) (0.34,0.01) (0.68,0.21) (1.02,0.01) (1.36,0.10)};
  }%
}
% Mot-clé surligné (marqueur orange sous le texte)
\newcommand{\cle}[1]{{\popxbold\color{popdark}#1}}

% ============================================================
% En-tête / pied
% ============================================================
\pagestyle{fancy}
\fancyhf{}
\renewcommand{\headrulewidth}{0pt}
\renewcommand{\footrulewidth}{0pt}
\lhead{\raisebox{-0.15ex}{\poplogo\fontsize{13}{13}\selectfont\color{popink}OnBuch\color{poporange}+}}
\rhead{\poppill{\DOCMATIERE\ \textbullet\ \DOCNIVEAU\ \textbullet\ Chapter \DOCLECON}{white}{popink}}
\lfoot{\popxbold\fontsize{7.6}{7.6}\selectfont\color{popmuted}\MakeUppercase{Signed course OnBuch+}\ \textbullet\ {\popsemi\DOCTITRE}}
\rfoot{\tikz[baseline=-0.6ex]\node[rounded corners=8.5pt,fill=popink,minimum height=17pt,minimum width=17pt,inner xsep=5pt,inner sep=0]{\popxbold\fontsize{8}{8}\selectfont\color{popcream}\thepage\,/\,\pageref*{LastPage}};}

% ============================================================
% Titres
% ============================================================
\newcommand{\chaptitre}[2]{%
  \begin{tcolorbox}[enhanced,colback=poporange,colframe=popink,boxrule=2.6pt,arc=11pt,
      drop shadow=popink,left=15pt,right=15pt,top=12pt,bottom=11pt,before skip=3pt,after skip=18pt]
    \poppill{#2}{popink}{popcream}\par\vspace{9pt}
    {\raggedright\hyphenpenalty=10000\exhyphenpenalty=10000\popdisplay\fontsize{22}{24}\selectfont\color{popcream}\MakeUppercase{#1}\par}
  \end{tcolorbox}
}
% Section numérotée : pastille orange avec le numéro + titre Archivo
\newcounter{popsec}
\newcommand{\coursec}[1]{%
  \stepcounter{popsec}\setcounter{popsub}{0}%
  \par\vspace{16pt}\noindent
  \begin{minipage}{\linewidth}
  \tikz[baseline=-0.7ex]\node[draw=popink,line width=1.6pt,fill=poporange,rounded corners=4pt,minimum size=22pt,inner sep=0]{\popdisplay\fontsize{12}{12}\selectfont\color{popcream}\thepopsec};%
  \hspace{8pt}{\popdisplay\fontsize{15}{17}\selectfont\color{popink}\MakeUppercase{#1}}\par
  \vspace{4pt}\noindent{\color{poporange}\rule{36pt}{3.6pt}}
  \end{minipage}\par\nopagebreak\vspace{8pt}
}
\newcounter{popsub}
\newcommand{\courssub}[1]{%
  \stepcounter{popsub}%
  \par\vspace{9pt}\noindent{\popxbold\fontsize{11.5}{13}\selectfont\color{popink}{\color{poporange}\thepopsec.\thepopsub}\hspace{6pt}#1}\par\nopagebreak\vspace{4pt}
}

% ============================================================
% Boîtes pédagogiques — même recette : fond blanc, bordure épaisse colorée,
% ombre dure encre, badge plein. Titre optionnel [#1].
% ============================================================
\tcbset{popbase/.style={breakable,enhanced,colback=white,boxrule=2.3pt,arc=9pt,
  shadow={3pt}{-3pt}{0pt}{popink},left=14pt,right=14pt,top=11pt,bottom=11pt,before skip=11pt,after skip=15pt,
  fontupper=\popsemi\fontsize{10.6}{14.5}\selectfont\color{popink},
  fontlower=\popsemi\fontsize{10.6}{14.5}\selectfont\color{popink}}}
% Coche dessinée (le glyphe ✓ n'existe pas dans les polices du document)
\newcommand{\popcheck}[1]{\tikz[baseline=-0.5ex]\draw[#1,line width=1.6pt,line cap=round,line join=round](-0.13,0.0)--(-0.04,-0.09)--(0.13,0.1);}
\providecommand{\coche}{\popcheck{popgreen}}
\providecommand{\resultat}[1]{\par\smallskip\noindent\fcolorbox{popgreen}{popgreenL}{\parbox{\dimexpr\linewidth-2\fboxsep-2\fboxrule\relax}{\textbf{Result.} #1}}\par\smallskip}
\newcommand{\popboxhead}[3]{\poppill{#1}{#2}{white}\ifx\relax#3\relax\else\hspace{6pt}{\popxbold\fontsize{10.6}{12}\selectfont\color{popink}#3}\fi\par\vspace{7pt}}

\newtcolorbox{aretenir}[1][]{popbase,colframe=popgreen,before upper={\popboxhead{Key points}{popgreen}{#1}}}
\newtcolorbox{exemplebox}[1][]{popbase,colframe=poporange,before upper={\popboxhead{Exemple}{poporange}{#1}}}
\newtcolorbox{definition}[1][]{popbase,colframe=popblue,before upper={\popboxhead{Definition}{popblue}{#1}}}
\newtcolorbox{propriete}[1][]{popbase,colframe=poppurple,before upper={\popboxhead{Property}{poppurple}{#1}}}
\newtcolorbox{methode}[1][]{popbase,colframe=popgold,before upper={\popboxhead{Method}{popgold}{#1}}}
\newtcolorbox{attention}[1][]{popbase,colframe=poppink,before upper={\popboxhead{Warning}{poppink}{#1}}}
\newtcolorbox{experience}[1][]{popbase,colframe=popblue,colback=popblueL,before upper={\popboxhead{Experiment}{popblue}{#1}}}
% « Le savais-tu ? » : carte violette pleine (contexte, culture, Cameroun)
\newtcolorbox{savaistu}[1][]{popbase,colback=poppurple,colframe=popink,
  fontupper=\popsemi\fontsize{10.4}{14.2}\selectfont\color{white},
  before upper={\poppill{Did you know?}{white}{poppurple}\ifx\relax#1\relax\else\hspace{6pt}{\popxbold\color{white}#1}\fi\par\vspace{7pt}}}
% Exercice résolu : énoncé en haut, \tcblower puis la solution sur fond crème
\newcounter{popexo}\newcounter{popexr}
\newtcolorbox{exoresolu}[1][]{popbase,colframe=poporange,colbacklower=popcream,
  segmentation style={popink,line width=1.2pt,dashed},
  before upper={\stepcounter{popexr}\popboxhead{Worked example \thepopexr}{poporange}{#1}},
  before lower={\poppill{Solution}{popink}{popcream}\par\vspace{6pt}}}
% Exercice (non résolu en place) — la correction est dans la partie Corrigés
\newtcolorbox{exercice}[1][]{popbase,colframe=popink,
  before upper={\stepcounter{popexo}\popboxhead{Exercise \thepopexo}{popink}{#1}}}
% Corrigé
\newtcolorbox{corrige}[1][]{popbase,colframe=popgreen,colback=popgreenL,
  before upper={\popboxhead{Solution}{popgreen}{#1}}}
% Objectifs / prérequis / bilan
\newtcolorbox{objectifs}{popbase,colframe=popink,colback=poporangeL,
  before upper={\popboxhead{Chapter objectives}{popink}{}}}
\newtcolorbox{prerequis}{popbase,colframe=popink,colback=popblueL,
  before upper={\popboxhead{Prerequisites}{popblue}{}}}
\newtcolorbox{bilan}{popbase,colframe=popink,colback=white,boxrule=2.8pt,
  before upper={\popboxhead{Fiche bilan}{poporange}{L'essentiel à connaître pour l'examen}}}
% Figure encadrée avec légende
\newtcolorbox{popfigure}[1][]{enhanced,breakable=false,colback=white,colframe=popline,boxrule=1.4pt,arc=8pt,
  left=10pt,right=10pt,top=10pt,bottom=8pt,before skip=10pt,after skip=12pt,halign=center,
  lower separated=false,halign lower=center,
  fontlower=\popsemi\fontsize{9}{11.5}\selectfont\color{popmuted},
  #1}
\newcommand{\legende}[1]{\par\vspace{6pt}{\popsemi\fontsize{9}{11.5}\selectfont\color{popmuted}#1}}

% ============================================================
% Signature OnBuch+
% ============================================================
\newcommand{\poplogoinline}[1]{{\poplogo\fontsize{#1}{#1}\selectfont\color{popink}OnBuch\color{poporange}+}}
\newcommand{\signatureonbuch}{%
  \begin{tcolorbox}[enhanced,colback=popink,colframe=popink,boxrule=2.2pt,arc=10pt,
      drop shadow=poporange,left=14pt,right=14pt,top=10pt,bottom=10pt,before skip=0pt,after skip=0pt]
    \tikz[baseline=-0.7ex]\node[circle,fill=popgreen,draw=popcream,line width=1.3pt,minimum size=22pt,inner sep=0]{\popcheck{white}};%
    \hspace{9pt}\begin{minipage}[c]{0.62\linewidth}
      {\popxbold\fontsize{12}{13}\selectfont\color{popcream}Signed\ \textbullet\ {\poplogo OnBuch\color{poporange}+}}\par\vspace{2pt}
      {\raggedright\popsemi\fontsize{8}{10.5}\selectfont\color{popline}\MakeUppercase{OnBuch+ teaching team}\\\MakeUppercase{Follows the official MINESEC syllabus}\par}
    \end{minipage}\hfill
    {\poplogo\fontsize{22}{22}\selectfont\color{popcream}OnBuch\color{poporange}+}
  \end{tcolorbox}%
}

% ============================================================
% Page de garde
% #1 titre · #2 matière · #3 niveau · #4 module · #5 n° leçon · #6 durée ·
% #7 description / objectifs (texte ou itemize)
% ============================================================
\newcommand{\pagegarde}[7]{%
  \thispagestyle{empty}
  \newgeometry{a4paper,margin=1.7cm,top=1.6cm,bottom=1.4cm}%
  \begin{tikzpicture}[remember picture,overlay]
    \fill[poporange!18] ([xshift=-3.2cm,yshift=-3.6cm]current page.north east) circle (4.6cm);
    \fill[poppurple!14] ([xshift=2.4cm,yshift=7.6cm]current page.south west) circle (3.8cm);
  \end{tikzpicture}%
  {\poplogo\fontsize{30}{30}\selectfont\color{popink}OnBuch\color{poporange}+}\hfill
  \raisebox{0.6ex}{\poppill{Signed course \textbullet\ Premium}{popink}{popcream}}\par
  \vspace{1pt}\popsquiggle{poppurple}\par
  \vspace{8pt}
  \begin{tcolorbox}[enhanced,colback=poporange,colframe=popink,boxrule=2.8pt,arc=13pt,
      drop shadow=popink,left=18pt,right=18pt,top=12pt,bottom=13pt,after skip=10pt]
    \poppill{#2 \textbullet\ #3}{popink}{popcream}\hspace{5pt}\poppill{#4}{white}{popink}\par\vspace{8pt}
    {\popdisplay\fontsize{14}{14}\selectfont\color{popink}CHAPTER #5}\par\vspace{4pt}
    {\raggedright\hyphenpenalty=10000\exhyphenpenalty=10000\popdisplay\fontsize{25}{27}\selectfont\color{popcream}\MakeUppercase{#1}\par}
    \vspace{8pt}
    {\popxbold\fontsize{9.5}{9.5}\selectfont\color{popink}\MakeUppercase{Suggested time: #6}}
  \end{tcolorbox}
  \begin{tcolorbox}[enhanced,colback=white,colframe=popink,boxrule=2.2pt,arc=10pt,
      drop shadow=popink,left=16pt,right=16pt,top=10pt,bottom=10pt,after skip=8pt]
    \poppill{In this chapter}{poppurple}{white}\par\vspace{5pt}
    \popsemi\fontsize{9.3}{12.6}\selectfont\color{popink}#7
  \end{tcolorbox}
  \noindent\begin{minipage}[t]{0.487\linewidth}
    \begin{tcolorbox}[enhanced,colback=popgreen,colframe=popink,boxrule=2.2pt,arc=10pt,
        drop shadow=popink,left=11pt,right=11pt,top=10pt,bottom=10pt,height=3.1cm,valign=center]
      \tcbox[colback=white,colframe=popink,boxrule=1.3pt,arc=5pt,boxsep=0pt,nobeforeafter,
        left=3pt,right=3pt,top=3pt,bottom=3pt,valign=center]{\qrcode[height=1.9cm]{https://play.google.com/store/apps/details?id=com.luvvix.onbuch&hl=en}}%
      \hspace{8pt}\begin{minipage}[c]{0.5\linewidth}
        {\popdisplay\fontsize{11}{12.5}\selectfont\color{white}\MakeUppercase{Download OnBuch}}\par\vspace{4pt}
        {\raggedright\popsemi\fontsize{8.4}{11}\selectfont\color{white}Free \textbullet\ Google Play\\AI tutor, past papers, results\par}
      \end{minipage}
    \end{tcolorbox}
  \end{minipage}\hfill
  \begin{minipage}[t]{0.487\linewidth}
    \begin{tcolorbox}[enhanced,colback=poppurple,colframe=popink,boxrule=2.2pt,arc=10pt,
        drop shadow=popink,left=11pt,right=11pt,top=10pt,bottom=10pt,height=3.1cm,valign=center]
      \tikz[baseline=-0.7ex]\node[circle,fill=white,draw=popink,line width=1.3pt,minimum size=1.6cm,inner sep=0]{\popdisplay\fontsize{24}{24}\selectfont\color{poppurple}+};%
      \hspace{8pt}\begin{minipage}[c]{0.6\linewidth}
        {\popdisplay\fontsize{11}{12.5}\selectfont\color{white}\MakeUppercase{Go OnBuch+}}\par\vspace{4pt}
        {\raggedright\popsemi\fontsize{8.4}{11}\selectfont\color{white}All signed courses, unlimited AI tutor, PDF sheets — OnBuch+ tab in the app\par}
      \end{minipage}
    \end{tcolorbox}
  \end{minipage}
  \vspace{8pt}
  \popcontenu
  \vfill
  \signatureonbuch
  \restoregeometry
  \newpage
}

% Rangée « Ce cours contient » (page de garde)
\newcommand{\poptile}[3]{% #1 couleur #2 gros texte #3 légende
  \begin{tcolorbox}[enhanced,colback=#1,colframe=popink,boxrule=2pt,arc=9pt,shadow={2.5pt}{-2.5pt}{0pt}{popink},
      left=6pt,right=6pt,top=9pt,bottom=9pt,halign=center,valign=center,height=1.75cm,width=\linewidth,nobeforeafter]
    {\popdisplay\fontsize{12.5}{13}\selectfont\color{white}\MakeUppercase{#2}}\par\vspace{3pt}
    {\centering\popsemi\fontsize{7.8}{9.5}\selectfont\color{white}#3\par}
  \end{tcolorbox}}
\newcommand{\popcontenu}{%
  \par\noindent{\popxboldls\fontsize{7.5}{7.5}\selectfont\color{popmuted}THIS SIGNED COURSE CONTAINS}\par\vspace{7pt}
  \noindent\begin{minipage}[t]{0.235\linewidth}\poptile{popblue}{Course}{complete, illustrated, step by step}\end{minipage}\hfill
  \begin{minipage}[t]{0.235\linewidth}\poptile{poporange}{Methods}{and worked examples}\end{minipage}\hfill
  \begin{minipage}[t]{0.235\linewidth}\poptile{poppink}{Exercises}{exam style + solutions}\end{minipage}\hfill
  \begin{minipage}[t]{0.235\linewidth}\poptile{popgreen}{Summary sheet}{the essentials to remember}\end{minipage}\par}

% Sommaire
\newtcolorbox{sommaire}{enhanced,colback=white,colframe=popink,boxrule=2.2pt,arc=10pt,
  drop shadow=popink,left=18pt,right=18pt,top=13pt,bottom=13pt,before skip=6pt,after skip=20pt,
  before upper={\poppill{Contents}{poppurple}{white}\par\vspace{9pt}\popsemi\fontsize{10.6}{14.5}\selectfont\color{popink}}}

% Signature de fin de cours — poussée tout en bas de la dernière page
\newcommand{\findecours}{%
  \par\vfill\nopagebreak
  \begin{center}\popsquiggle{poporange}\end{center}\vspace{4pt}
  \signatureonbuch
}
\AtBeginDocument{\def\llbracket{\mathopen{[\mkern-3.5mu[}}\def\rrbracket{\mathclose{]\mkern-3.5mu]}}} % intervalles d'entiers (glyphe ⟦ absent de la police)
% Glyphes absents de Fira Math : redéfinis à partir de glyphes disponibles
\AtBeginDocument{%
  \def\setminus{\mathbin{\backslash}}\let\smallsetminus\setminus
  \def\vdots{\mathord{\vbox{\baselineskip3.2pt\lineskiplimit0pt\kern4pt\hbox{.}\hbox{.}\hbox{.}}}}%
  \def\ddots{\mathinner{\mkern1mu\raise6.5pt\hbox{.}\mkern2mu\raise3.6pt\hbox{.}\mkern2mu\raise0.7pt\hbox{.}\mkern1mu}}%
  \def\triangle{\mathord{\scalebox{0.9}{$\symup{\Delta}$}}}%
  \def\nabla{\mathord{\raisebox{\depth}{\scalebox{1}[-1]{$\symup{\Delta}$}}}}%
  \def\checkmark{\popcheck{popdark}}%
  \def\star{\mathbin{\ast}}\def\bigstar{\mathord{\ast}}%
  \def\diamond{\mathbin{\scalebox{0.6}{$\Diamond$}}}%
  \def\bigcup{\mathop{\mathchoice{\raisebox{-0.3em}{\scalebox{1.7}{$\cup$}}}{\scalebox{1.25}{$\cup$}}{\cup}{\cup}}\displaylimits}%
  \def\bigcap{\mathop{\mathchoice{\raisebox{-0.3em}{\scalebox{1.7}{$\cap$}}}{\scalebox{1.25}{$\cap$}}{\cap}{\cap}}\displaylimits}%
  \def\lesseqqgtr{\mathrel{\lessgtr}}\def\bigoplus{\mathop{\oplus}\displaylimits}%
}
\providecommand{\ssi}{if and only if} % abréviation fréquente
\AtBeginDocument{\providecommand{\wideparen}{\overparen}} % arc de cercle

% Fonctions usuelles que les modèles utilisent sans que amsmath les fournisse
\providecommand{\arsinh}{\operatorname{arsinh}}\providecommand{\arcosh}{\operatorname{arcosh}}
\providecommand{\artanh}{\operatorname{artanh}}\providecommand{\arsech}{\operatorname{arsech}}
\providecommand{\arcsch}{\operatorname{arcsch}}\providecommand{\arcoth}{\operatorname{arcoth}}
\providecommand{\arcsinh}{\operatorname{arsinh}}\providecommand{\arccosh}{\operatorname{arcosh}}
\providecommand{\arctanh}{\operatorname{artanh}}\providecommand{\sech}{\operatorname{sech}}
\providecommand{\csch}{\operatorname{csch}}\providecommand{\arcsec}{\operatorname{arcsec}}
\providecommand{\arccsc}{\operatorname{arccsc}}\providecommand{\arccot}{\operatorname{arccot}}
\providecommand{\sgn}{\operatorname{sgn}}\providecommand{\lcm}{\operatorname{lcm}}
\providecommand{\Var}{\operatorname{Var}}\providecommand{\Cov}{\operatorname{Cov}}

%% FIGFIX-BEGIN v5 — correctifs globaux des figures (docs/onbuch-generation/tools/figfix)
\usepackage{adjustbox}\usepackage{etoolbox}\tcbuselibrary{raster}
% toute figure TikZ/pgfplots plus large que la ligne est réduite à la largeur disponible
\BeforeBeginEnvironment{tikzpicture}{\begin{adjustbox}{max width=\linewidth}}
\AfterEndEnvironment{tikzpicture}{\end{adjustbox}}
% légendes pgfplots lisibles par-dessus les courbes
\pgfplotsset{every axis/.append style={legend style={fill=white,fill opacity=0.92,text opacity=1,draw=black!15,inner sep=3pt}}}
% étiquettes d'arêtes (syntaxe edge["..."]) avec fond blanc
\tikzset{every edge quotes/.append style={fill=white,inner sep=1.2pt}}
%% FIGFIX-END
\begin{document}

\pagegarde{Integration}{Pure Mathematics with Statistics}{Upper Sixth}{Upper Sixth — Pure mathematics with statistics}{4}{9 periods}{This chapter presents the definite integral as the limit of Riemann sums and develops its great geometric and physical applications: areas between curves and axes, volumes of solids of revolution, centroids and centres of mass. It also introduces the trapezium rule for numerical approximation and the mean value theorem for integrals, all in the spirit of the GCE Advanced Level examination.
\vspace{4pt}
\begin{multicols}{2}\small
\begin{itemize}
\item By the end of this chapter I can express a definite integral as the limit of a Riemann sum and interpret it as a signed area.
\item By the end of this chapter I can use definite integrals to calculate the area between a curve and the x-axis, including regions below the axis.
\item By the end of this chapter I can calculate the area between a curve and the y-axis, and the area between two curves.
\item By the end of this chapter I can calculate volumes of solids of revolution about the x-axis and about the y-axis.
\item By the end of this chapter I can use integration to find the centroid of a plane region and the centre of mass of a lamina or solid.
\item By the end of this chapter I can approximate a definite integral with the trapezium rule and estimate the error involved.
\end{itemize}\end{multicols}}

\begin{sommaire}
\begin{multicols}{2}
\begin{enumerate}
\item The Definite Integral as a Limit of Riemann Sums
\item Areas Bounded by Curves and the Coordinate Axes
\item Volumes of Solids of Revolution
\item Centroids and Centres of Mass by Integration
\item Numerical Integration: the Trapezium Rule
\item Mean Value Theorem for Integrals and Synthesis
\item Integration activity
\item Methods and worked examples
\item Exercises
\item Solutions to the exercises
\item Summary sheet
\end{enumerate}
\end{multicols}
\end{sommaire}

\begin{objectifs}
\begin{itemize}
\item By the end of this chapter I can express a definite integral as the limit of a Riemann sum and interpret it as a signed area.
\item By the end of this chapter I can use definite integrals to calculate the area between a curve and the x-axis, including regions below the axis.
\item By the end of this chapter I can calculate the area between a curve and the y-axis, and the area between two curves.
\item By the end of this chapter I can calculate volumes of solids of revolution about the x-axis and about the y-axis.
\item By the end of this chapter I can use integration to find the centroid of a plane region and the centre of mass of a lamina or solid.
\item By the end of this chapter I can approximate a definite integral with the trapezium rule and estimate the error involved.
\item By the end of this chapter I can state and apply the mean value theorem for integrals and compute the mean value of a function on an interval.
\item By the end of this chapter I can model and solve authentic GCE-style problems combining several of these applications.
\end{itemize}
\end{objectifs}

\begin{prerequis}
\begin{itemize}
\item Indefinite integration: standard integrals, integration by substitution and by parts (Lower Sixth / PMS03).
\item The Fundamental Theorem of Calculus: evaluating a definite integral with a primitive.
\item Curve sketching of polynomials, rational, trigonometric, exponential and logarithmic functions.
\item Solving equations to find intersection points of curves.
\item Sigma notation and limits of sequences and sums (Lower Sixth).
\end{itemize}
\end{prerequis}

\newpage

\coursec{The Definite Integral as a Limit of Riemann Sums}

During the rainy season, engineers at the Lagdo dam must estimate the volume of water stored behind the curved wall of the reservoir, while a cocoa cooperative in Kumba wants to know the exact area of an irregularly shaped plot bounded by a stream. Neither the area nor the volume can be found with the simple formulas of geometry, because the boundaries are \emph{curves}, not straight lines. Yet by slicing the region into very thin strips, summing their contributions, and letting the strips become infinitely thin, integration gives the \emph{exact} answer. In this chapter you will learn how the definite integral, born as a limit of sums, becomes a powerful measuring tool for areas, volumes, centres of mass and average values — tools used daily by Cameroonian engineers, surveyors and economists.

The central question of this section is therefore: \textbf{how can we define and compute the area of a region bounded by a curve?} The answer, due to Bernhard Riemann, is one of the most beautiful ideas in all of mathematics.

\courssub{The Area Problem and the Idea of Exhaustion}

You already know how to find the area of a rectangle ($\text{length}\times\text{width}$), a triangle, a trapezium, and even a circle ($\pi r^2$, a formula you were asked to accept on trust). But consider the region under the parabola $y=x^2$, above the $x$-axis, between $x=0$ and $x=1$. No elementary formula applies.

The strategy — used since Archimedes and called the \cle{method of exhaustion} — is simple and powerful:
\begin{enumerate}
\item \textbf{Slice} the region into $n$ thin vertical strips.
\item \textbf{Approximate} each strip by a rectangle whose area we can compute.
\item \textbf{Sum} the areas of the rectangles to get an approximation of the total area.
\item \textbf{Refine}: let $n$ grow without bound so that the strips become infinitely thin; the approximations approach a limit, which we \emph{define} to be the exact area.
\end{enumerate}

The figure below shows this idea for $y=x^2$ on $[0,1]$: with $n=4$ rectangles the approximation is crude, but with $n=10$ the rectangles already hug the curve closely.

\begin{popfigure}
\begin{tikzpicture}[scale=3.2]
% n = 4 right rectangles
\foreach \i in {1,2,3,4}{
  \pgfmathsetmacro{\x}{\i/4}
  \pgfmathsetmacro{\y}{\x*\x}
  \draw[fill=popblueL, draw=popblue] ({\x-0.25},0) rectangle (\x,\y);
}
\draw[->] (-0.1,0) -- (1.25,0) node[right] {$x$};
\draw[->] (0,-0.1) -- (0,1.2) node[above] {$y$};
\draw[domain=0:1.1, smooth, thick, popdark] plot (\x,{\x*\x}) node[right] {$y=x^2$};
\foreach \i in {0,1,2,3,4}{\draw ({\i/4},0.02) -- ({\i/4},-0.02) node[below] {\tiny $\i/4$};}
\node at (0.5,-0.35) {$n=4$ (right endpoints)};
\end{tikzpicture}\hfill
\begin{tikzpicture}[scale=3.2]
% n = 10 right rectangles
\foreach \i in {1,...,10}{
  \pgfmathsetmacro{\x}{\i/10}
  \pgfmathsetmacro{\y}{\x*\x}
  \draw[fill=popblueL, draw=popblue] ({\x-0.1},0) rectangle (\x,\y);
}
\draw[->] (-0.1,0) -- (1.25,0) node[right] {$x$};
\draw[->] (0,-0.1) -- (0,1.2) node[above] {$y$};
\draw[domain=0:1.1, smooth, thick, popdark] plot (\x,{\x*\x}) node[right] {$y=x^2$};
\node at (0.5,-0.35) {$n=10$ (right endpoints)};
\end{tikzpicture}
\legende{Right-endpoint rectangles approximating the area under $y=x^2$ on $[0,1]$. As $n$ increases, the staircase fills the region ever more closely.}
\end{popfigure}

\courssub{Partitions, Lower Sums and Upper Sums}

\begin{definition}[Partition of an interval]
A \cle{partition} of the interval $[a,b]$ into $n$ sub-intervals is a finite set of points
\[ a = x_0 < x_1 < x_2 < \dots < x_{n-1} < x_n = b. \]
The $i$-th sub-interval is $[x_{i-1},x_i]$, of width $\Delta x_i = x_i - x_{i-1}$. The \cle{mesh} of the partition is $\max_i \Delta x_i$. The partition is \cle{regular} (or uniform) when all sub-intervals have the same width $\Delta x = \dfrac{b-a}{n}$, so that $x_i = a + i\,\dfrac{b-a}{n}$.
\end{definition}

Suppose $f$ is continuous and positive on $[a,b]$. On each sub-interval $[x_{i-1},x_i]$, since $f$ is continuous it attains a minimum value $m_i$ and a maximum value $M_i$. We build two families of rectangles:

\begin{itemize}
\item the \cle{inscribed} rectangles, of height $m_i$ (they lie inside the region under the curve);
\item the \cle{circumscribed} rectangles, of height $M_i$ (they cover the region).
\end{itemize}

\begin{definition}[Lower and upper sums]
The \cle{lower sum} and the \cle{upper sum} of $f$ over the partition are
\[ L_n = \sum_{i=1}^{n} m_i\,\Delta x_i, \qquad U_n = \sum_{i=1}^{n} M_i\,\Delta x_i. \]
If $A$ denotes the (unknown) exact area under the curve, then for every partition
\[ L_n \;\le\; A \;\le\; U_n. \]
\end{definition}

For an \emph{increasing} function on a regular partition, the minimum on $[x_{i-1},x_i]$ is $f(x_{i-1})$ (left endpoint) and the maximum is $f(x_i)$ (right endpoint). Hence
\[ L_n = \frac{b-a}{n}\sum_{i=0}^{n-1} f(x_i), \qquad U_n = \frac{b-a}{n}\sum_{i=1}^{n} f(x_i). \]

\begin{popfigure}
\begin{tikzpicture}
\begin{axis}[width=6.6cm, height=5.4cm, xmin=-0.05, xmax=1.15, ymin=0, ymax=1.15,
 axis lines=middle, xlabel={$x$}, ylabel={$y$}, title={Lower sum $L_4$},
 xtick={0,0.25,0.5,0.75,1}, xticklabels={$0$,$\frac14$,$\frac12$,$\frac34$,$1$},
 tick label style={font=\tiny}, title style={font=\small}, samples=60]
\addplot[domain=0:1.08, thick, color=popdark] {x^2};
\foreach \i in {0,1,2,3}{
  \pgfmathsetmacro{\xa}{\i/4}
  \pgfmathsetmacro{\xb}{(\i+1)/4}
  \pgfmathsetmacro{\ya}{\xa*\xa}
  \addplot[draw=popgreen, fill=popgreenL, forget plot] coordinates {(\xa,0) (\xa,\ya) (\xb,\ya) (\xb,0)} \closedcycle;
}
\end{axis}
\end{tikzpicture}\hfill
\begin{tikzpicture}
\begin{axis}[width=6.6cm, height=5.4cm, xmin=-0.05, xmax=1.15, ymin=0, ymax=1.15,
 axis lines=middle, xlabel={$x$}, ylabel={$y$}, title={Upper sum $U_4$},
 xtick={0,0.25,0.5,0.75,1}, xticklabels={$0$,$\frac14$,$\frac12$,$\frac34$,$1$},
 tick label style={font=\tiny}, title style={font=\small}, samples=60]
\addplot[domain=0:1.08, thick, color=popdark] {x^2};
\foreach \i in {0,1,2,3}{
  \pgfmathsetmacro{\xa}{\i/4}
  \pgfmathsetmacro{\xb}{(\i+1)/4}
  \pgfmathsetmacro{\yb}{\xb*\xb}
  \addplot[draw=poporange, fill=poporange!20, forget plot] coordinates {(\xa,0) (\xa,\yb) (\xb,\yb) (\xb,0)} \closedcycle;
}
\end{axis}
\end{tikzpicture}
\legende{Lower sum (inscribed rectangles, left) and upper sum (circumscribed rectangles, right) for $f(x)=x^2$ on $[0,1]$ with $n=4$. The exact area is trapped between the two.}
\end{popfigure}

\begin{exemplebox}[{Bounding the area under $y=x^2$ on $[0,1]$}]
With $n=4$, $\Delta x = \tfrac14$ and $f$ increasing:
\begin{align*}
L_4 &= \tfrac14\left[f(0)+f\!\left(\tfrac14\right)+f\!\left(\tfrac12\right)+f\!\left(\tfrac34\right)\right]
= \tfrac14\left(0+\tfrac{1}{16}+\tfrac{4}{16}+\tfrac{9}{16}\right) = \tfrac{14}{64} = 0.21875,\\
U_4 &= \tfrac14\left[f\!\left(\tfrac14\right)+f\!\left(\tfrac12\right)+f\!\left(\tfrac34\right)+f(1)\right]
= \tfrac14\left(\tfrac{1}{16}+\tfrac{4}{16}+\tfrac{9}{16}+\tfrac{16}{16}\right) = \tfrac{30}{64} = 0.46875.
\end{align*}
So $0.21875 \le A \le 0.46875$. The exact value, as we shall prove below, is $A=\tfrac13 \approx 0.3333$: indeed trapped between the two bounds.
\end{exemplebox}

\courssub{Riemann Sums with Arbitrary Sample Points}

Lower and upper sums use the extreme values of $f$ on each strip, which may be hard to find. Riemann's insight was that \emph{any} choice of height works in the limit.

\begin{definition}[Riemann sum]
Let $f$ be defined on $[a,b]$ and let $a=x_0<x_1<\dots<x_n=b$ be a partition. In each sub-interval $[x_{i-1},x_i]$ choose any \cle{sample point} $c_i$. The sum
\[ S_n = \sum_{i=1}^{n} f(c_i)\,\Delta x_i \]
is called a \cle{Riemann sum} of $f$ over the partition. For a regular partition, $\Delta x = \dfrac{b-a}{n}$ and the classical choices are:
\begin{itemize}
\item \textbf{left-endpoint sum}: $c_i = x_{i-1}$, \quad $S_n^{L} = \dfrac{b-a}{n}\displaystyle\sum_{i=0}^{n-1} f\!\left(a+i\tfrac{b-a}{n}\right)$;
\item \textbf{right-endpoint sum}: $c_i = x_i$, \quad $S_n^{R} = \dfrac{b-a}{n}\displaystyle\sum_{i=1}^{n} f\!\left(a+i\tfrac{b-a}{n}\right)$;
\item \textbf{midpoint sum}: $c_i = \dfrac{x_{i-1}+x_i}{2}$, \quad $S_n^{M} = \dfrac{b-a}{n}\displaystyle\sum_{i=1}^{n} f\!\left(a+\left(i-\tfrac12\right)\tfrac{b-a}{n}\right)$.
\end{itemize}
\end{definition}

Since $m_i \le f(c_i) \le M_i$ on each sub-interval, every Riemann sum satisfies $L_n \le S_n \le U_n$. If the lower and upper sums squeeze towards the same number as the mesh tends to $0$, then \emph{every} Riemann sum is forced to the same limit.

\begin{definition}[The definite integral as a limit of Riemann sums]
Let $f$ be defined on $[a,b]$. If the Riemann sums $S_n = \sum_{i=1}^n f(c_i)\,\Delta x_i$ approach one and the same finite limit as the mesh of the partition tends to $0$, \emph{whatever} the choice of partitions and sample points, then $f$ is said to be \cle{integrable} on $[a,b]$, and this limit is the \cle{definite integral} of $f$ from $a$ to $b$:
\[ \int_a^b f(x)\,dx \;=\; \lim_{\max\Delta x_i \to 0}\ \sum_{i=1}^{n} f(c_i)\,\Delta x_i
\;=\; \lim_{n\to\infty}\ \frac{b-a}{n}\sum_{i=1}^{n} f\!\left(a+i\tfrac{b-a}{n}\right) \]
(the second equality holding for regular partitions). We call $a$ and $b$ the \cle{limits of integration}, $f$ the \cle{integrand}, and $x$ a \cle{dummy variable}: $\int_a^b f(x)\,dx = \int_a^b f(t)\,dt$.
\end{definition}

\begin{propriete}[Integrability of continuous functions — admitted]
Every function \textbf{continuous} on $[a,b]$ is integrable on $[a,b]$: its lower and upper sums converge to a common limit. The same holds for bounded functions with a finite number of jump discontinuities. (The proof belongs to university analysis and is not examinable; you may use the result freely.)
\end{propriete}

\courssub{First Principles: Computing Integrals as Limits of Sums}

We shall need the classical summation formulas (proved by induction in the chapter on sequences):
\[ \sum_{i=1}^{n} i = \frac{n(n+1)}{2}, \qquad
\sum_{i=1}^{n} i^2 = \frac{n(n+1)(2n+1)}{6}. \]

\begin{methode}[Evaluating an integral from first principles]
\begin{enumerate}
\item Write the regular partition of $[a,b]$: $\Delta x = \dfrac{b-a}{n}$, $x_i = a + i\Delta x$.
\item Choose sample points (right endpoints are usually easiest) and form the Riemann sum $S_n = \Delta x \sum_{i=1}^n f(x_i)$.
\item Evaluate the sum in closed form using known summation formulas.
\item Take the limit as $n\to\infty$: this limit \emph{is} $\int_a^b f(x)\,dx$.
\end{enumerate}
\end{methode}

\begin{exoresolu}[Computing $\displaystyle\int_0^1 x\,dx$ from the limit definition]
Compute $\displaystyle\int_0^1 x\,dx$ directly as a limit of Riemann sums, and check against the geometric area.
\tcblower
\textbf{Solution.} Here $a=0$, $b=1$, so $\Delta x = \dfrac1n$ and $x_i = \dfrac{i}{n}$. Using right endpoints with $f(x)=x$:
\[ S_n = \sum_{i=1}^{n} f(x_i)\,\Delta x = \frac1n \sum_{i=1}^n \frac{i}{n} = \frac{1}{n^2}\sum_{i=1}^n i = \frac{1}{n^2}\cdot\frac{n(n+1)}{2} = \frac{n+1}{2n} = \frac12\left(1+\frac1n\right). \]
Therefore
\[ \int_0^1 x\,dx = \lim_{n\to\infty} \frac12\left(1+\frac1n\right) = \frac12. \]
\textbf{Check.} The region is a triangle of base $1$ and height $1$, of area $\tfrac12\times 1\times 1 = \tfrac12$. \checkmark
\end{exoresolu}

\begin{exoresolu}[Computing $\displaystyle\int_0^1 x^2\,dx$ from the limit definition]
Prove, using only the limit definition, that $\displaystyle\int_0^1 x^2\,dx = \frac13$.
\tcblower
\textbf{Solution.} With $\Delta x = \dfrac1n$, $x_i = \dfrac{i}{n}$ and right endpoints:
\[ S_n = \frac1n\sum_{i=1}^n \left(\frac{i}{n}\right)^2 = \frac{1}{n^3}\sum_{i=1}^n i^2 = \frac{1}{n^3}\cdot\frac{n(n+1)(2n+1)}{6} = \frac{(n+1)(2n+1)}{6n^2}. \]
Expanding, $S_n = \dfrac{2n^2+3n+1}{6n^2} = \dfrac13 + \dfrac{1}{2n} + \dfrac{1}{6n^2}$. Hence
\[ \int_0^1 x^2\,dx = \lim_{n\to\infty}\left(\frac13 + \frac{1}{2n} + \frac{1}{6n^2}\right) = \frac13. \]
This confirms the bound $0.21875 \le \tfrac13 \le 0.46875$ found earlier, and agrees with $\left[\dfrac{x^3}{3}\right]_0^1 = \dfrac13$ from antiderivatives.
\end{exoresolu}

\courssub{Signed Area: What the Integral Really Measures}

So far we assumed $f \ge 0$, so the integral equals a genuine geometric area. But the limit definition makes sense for \emph{any} continuous $f$, positive or negative. When $f(c_i) < 0$, the product $f(c_i)\Delta x_i$ is negative: the strip contributes \emph{negatively} to the sum.

\begin{aretenir}[The definite integral is a signed area]
$\displaystyle\int_a^b f(x)\,dx$ is the \cle{algebraic (signed) area} between the curve, the $x$-axis and the lines $x=a$, $x=b$:
\begin{itemize}
\item regions \textbf{above} the $x$-axis contribute $+$ their area;
\item regions \textbf{below} the $x$-axis contribute $-$ their area.
\end{itemize}
The \emph{geometric} (total) area is obtained by integrating $|f|$, i.e.\ splitting at the zeros of $f$ and taking absolute values of the pieces.
\end{aretenir}

\begin{popfigure}
\begin{tikzpicture}
\begin{axis}[width=11.5cm, height=5.6cm, xmin=-0.2, xmax=3.4, ymin=-1.6, ymax=1.8,
 axis lines=middle, xlabel={$x$}, ylabel={$y$}, samples=80,
 xtick={1,2,3}, ytick={-1,1}, tick label style={font=\small}]
\addplot[domain=0:3.2, thick, color=popdark] {2*sin(deg(1.7*x)) - 0.4};
\addplot[draw=none, fill=popgreenL, domain=0:2.1] {max(2*sin(deg(1.7*x)) - 0.4,0)} \closedcycle;
\addplot[draw=none, fill=poppink!35, domain=2.1:3.2] {min(2*sin(deg(1.7*x)) - 0.4,0)} \closedcycle;
\node at (axis cs:1.0,0.75) {$+A_1$};
\node at (axis cs:2.75,-0.75) {$-A_2$};
\end{axis}
\end{tikzpicture}
\legende{Signed area: $\displaystyle\int_a^b f(x)\,dx = A_1 - A_2$. The green region counts positively, the pink region negatively.}
\end{popfigure}

\begin{exemplebox}[An integral that is zero, an area that is not]
$\displaystyle\int_0^{2\pi} \sin x\,dx = \big[-\cos x\big]_0^{2\pi} = -1-(-1) = 0$, because the hump above the axis on $[0,\pi]$ (area $2$) exactly cancels the hump below on $[\pi,2\pi]$ (area $2$). The \emph{total geometric area} is $2+2 = 4$, computed as $\displaystyle\int_0^{\pi}\sin x\,dx + \left|\int_{\pi}^{2\pi}\sin x\,dx\right| = 2+2$.
\end{exemplebox}

\begin{attention}[Integral $\neq$ area]
A definite integral can be \textbf{negative} or even \textbf{zero} while the geometric area between the curve and the axis is strictly positive. At the GCE, when a question asks for an \emph{area}, you must locate the zeros of $f$ on the interval, split the integral there, and add the \emph{absolute values} of the pieces. Answering with a bare $\int_a^b f(x)\,dx$ when $f$ changes sign is a classic way to lose marks.
\end{attention}

\courssub{Properties Deduced from the Limit Definition}

Each property below follows by applying the corresponding operation to the Riemann sums and then passing to the limit.

\begin{propriete}[Linearity and Chasles relation — proofs examinable]
Let $f,g$ be continuous on $[a,b]$ and $\lambda,\mu\in\mathbb{R}$. Then:
\begin{enumerate}
\item \textbf{Linearity:} $\displaystyle\int_a^b \big(\lambda f(x)+\mu g(x)\big)\,dx = \lambda\int_a^b f(x)\,dx + \mu\int_a^b g(x)\,dx$.
\item \textbf{Additivity (Chasles):} for any $c\in[a,b]$, $\displaystyle\int_a^b f(x)\,dx = \int_a^c f(x)\,dx + \int_c^b f(x)\,dx$.
\item \textbf{Conventions:} $\displaystyle\int_a^a f(x)\,dx = 0$ and $\displaystyle\int_b^a f(x)\,dx = -\int_a^b f(x)\,dx$.
\end{enumerate}
\end{propriete}

\textbf{Justification from Riemann sums.} For linearity, a Riemann sum of $\lambda f + \mu g$ is
\[ \sum_{i=1}^n \big(\lambda f(c_i)+\mu g(c_i)\big)\Delta x_i = \lambda\sum_{i=1}^n f(c_i)\Delta x_i + \mu\sum_{i=1}^n g(c_i)\Delta x_i, \]
and taking the limit as $n\to\infty$ on both sides gives the result. For Chasles, refine partitions of $[a,b]$ so that $c$ is always a partition point; then each Riemann sum over $[a,b]$ splits into a sum over $[a,c]$ plus a sum over $[c,b]$, and the limit preserves the equality. \hfill$\square$

\begin{propriete}[Comparison and bounding properties]
Let $f,g$ be continuous on $[a,b]$ with $a<b$.
\begin{enumerate}
\item \textbf{Positivity:} if $f(x)\ge 0$ on $[a,b]$, then $\displaystyle\int_a^b f(x)\,dx \ge 0$.
\item \textbf{Comparison:} if $f(x)\le g(x)$ on $[a,b]$, then $\displaystyle\int_a^b f(x)\,dx \le \int_a^b g(x)\,dx$.
\item \textbf{Bounding:} if $m\le f(x)\le M$ on $[a,b]$, then
\[ m(b-a) \;\le\; \int_a^b f(x)\,dx \;\le\; M(b-a). \]
\item \textbf{Absolute value:} $\displaystyle\left|\int_a^b f(x)\,dx\right| \le \int_a^b |f(x)|\,dx$.
\end{enumerate}
Each statement is immediate on Riemann sums (e.g.\ $f(c_i)\ge 0$ forces $S_n\ge 0$, and a limit of non-negative numbers is non-negative) and passes to the limit.
\end{propriete}

\begin{exemplebox}[Bounding an integral we cannot compute exactly]
On $[0,1]$, $0 \le \dfrac{x^2}{1+x^3} \le x^2$ (since $1+x^3\ge 1$). By comparison,
\[ 0 \;\le\; \int_0^1 \frac{x^2}{1+x^3}\,dx \;\le\; \int_0^1 x^2\,dx = \frac13. \]
(The exact value is $\tfrac13\ln 2 \approx 0.231$, indeed within the bounds.)
\end{exemplebox}

\courssub{The Link with the Fundamental Theorem of Calculus}

Computing limits of Riemann sums by hand is instructive but quickly becomes impossible (try $f(x)=\sin x$!). The \cle{Fundamental Theorem of Calculus}, studied in the chapter on primitives, rescues us.

\begin{aretenir}[Fundamental Theorem of Calculus — recalled]
If $f$ is continuous on $[a,b]$ and $F$ is \emph{any} primitive (antiderivative) of $f$, i.e.\ $F'=f$, then
\[ \int_a^b f(x)\,dx = F(b) - F(a) = \big[F(x)\big]_a^b. \]
Thus the mysterious limit of sums is computed exactly by antidifferentiation: the two faces of calculus — differentiation and integration — are inverse operations.
\end{aretenir}

For instance, since $F(x)=\dfrac{x^3}{3}$ satisfies $F'(x)=x^2$, the theorem gives instantly $\displaystyle\int_0^1 x^2\,dx = \dfrac{1^3}{3}-\dfrac{0^3}{3} = \dfrac13$ — exactly the limit we laboriously obtained from the sum formula. The Riemann-sum viewpoint remains essential: it is what \emph{defines} the integral, it justifies every application (areas, volumes, centroids, averages) by a slicing-and-summing argument, and it underlies the numerical methods we shall meet later in this chapter.

\begin{exercice}[From first principles]
\begin{enumerate}
\item Using right-endpoint Riemann sums and the formula $\sum_{i=1}^n i = \dfrac{n(n+1)}{2}$, show that $\displaystyle\int_0^2 x\,dx = 2$. Check with the area of a triangle.
\item Using $\sum_{i=1}^n i^2 = \dfrac{n(n+1)(2n+1)}{6}$, compute $\displaystyle\int_0^2 x^2\,dx$ from the limit definition and verify with the Fundamental Theorem.
\item Show that for $f(x)=x$ on $[0,1]$, the left-endpoint sum is $S_n^{L}=\dfrac{n-1}{2n}$ and the right-endpoint sum is $S_n^{R}=\dfrac{n+1}{2n}$. Conclude that both tend to $\dfrac12$ and that $S_n^{R} - S_n^{L} = \dfrac1n \to 0$.
\end{enumerate}
\end{exercice}

\begin{corrige}[Exercise 1]
\textbf{1.} $\Delta x = \dfrac2n$, $x_i = \dfrac{2i}{n}$. Then $S_n = \dfrac2n\sum_{i=1}^n \dfrac{2i}{n} = \dfrac{4}{n^2}\cdot\dfrac{n(n+1)}{2} = \dfrac{2(n+1)}{n} \to 2$. Geometric check: triangle of base $2$, height $2$, area $\tfrac12\times2\times2=2$. \checkmark

\textbf{2.} $S_n = \dfrac2n\sum_{i=1}^n \left(\dfrac{2i}{n}\right)^2 = \dfrac{8}{n^3}\cdot\dfrac{n(n+1)(2n+1)}{6} = \dfrac{4(n+1)(2n+1)}{3n^2} \to \dfrac{4\times 2}{3} = \dfrac83$. Check: $\left[\dfrac{x^3}{3}\right]_0^2 = \dfrac83$. \checkmark

\textbf{3.} $S_n^{L} = \dfrac1n\sum_{i=0}^{n-1}\dfrac{i}{n} = \dfrac{1}{n^2}\cdot\dfrac{(n-1)n}{2} = \dfrac{n-1}{2n}$; $S_n^{R} = \dfrac1n\sum_{i=1}^{n}\dfrac{i}{n} = \dfrac{n+1}{2n}$. Both tend to $\tfrac12$, and $S_n^{R} - S_n^{L} = \dfrac{2}{2n} = \dfrac1n \to 0$: the left and right sums squeeze the area, confirming integrability and the value $\int_0^1 x\,dx = \tfrac12$.
\end{corrige}

\coursec{Areas Bounded by Curves and the Coordinate Axes}

In the previous section we built the definite integral from the ground up: we approximated the region under a curve by thin rectangles, wrote the sum $\sum_{i=1}^{n} f(x_i^*)\,\Delta x$, and defined $\int_a^b f(x)\,dx$ as the limit of these Riemann sums. That construction was not an accident. The integral was \emph{born} as an area. In this section we harvest the fruit of that work: we learn to compute, quickly and rigorously, the areas of regions bounded by curves and by the coordinate axes, including regions below the $x$-axis, regions trapped between two curves, and regions described most naturally by $x$ as a function of $y$.

\begin{experience}[Discovering the area between two curves]
Take two continuous curves $y=f(x)$ and $y=g(x)$ on $[a,b]$ with $f(x) \ge g(x)$. Partition $[a,b]$ into $n$ subintervals of width $\Delta x$. On the $i$-th subinterval, construct a rectangle whose height is the vertical distance between the curves: $f(x_i^*) - g(x_i^*)$. The area of this rectangle is $\bigl(f(x_i^*)-g(x_i^*)\bigr)\Delta x$. Summing over all rectangles gives a Riemann sum $\sum_{i=1}^n \bigl(f(x_i^*)-g(x_i^*)\bigr)\Delta x$. As $n\to\infty$, this sum tends to $\int_a^b \bigl(f(x)-g(x)\bigr)\,dx$. This geometric argument justifies the formula for the area between two curves without any new theory.
\end{experience}

\courssub{Area Between a Curve and the $x$-axis}

\textbf{The positive case.}\par
Suppose a trader in Bamenda records the rate at which water flows into a tank; the total volume collected is the area under the rate--time graph. More generally, let $f$ be continuous on $[a,b]$ with $f(x) \geq 0$ for all $x \in [a,b]$. Each rectangle in a Riemann sum has height $f(x_i^*) \geq 0$, so every partial sum is a genuine sum of rectangle areas, and the limit is exactly the area of the region under the curve.

\begin{propriete}[Area between a curve and the $x$-axis]
Let $f$ be continuous on $[a,b]$.
\begin{itemize}
\item If $f(x) \geq 0$ on $[a,b]$, the area of the region bounded by $y=f(x)$, the $x$-axis, and the lines $x=a$, $x=b$ is
\[ A = \int_a^b f(x)\,dx. \]
\item If $f(x) \leq 0$ on $[a,b]$, the integral is negative and the area is
\[ A = -\int_a^b f(x)\,dx = \int_a^b \bigl(-f(x)\bigr)\,dx. \]
\end{itemize}
Area is always a \cle{non-negative} number, measured in square units.
\end{propriete}

\begin{exemplebox}[A parabolic region]
Find the area enclosed by $y = 4 - x^2$ and the $x$-axis.

The curve meets the $x$-axis where $4-x^2=0$, i.e.\ $x = \pm 2$, and $4-x^2 \geq 0$ on $[-2,2]$. Hence
\[ A = \int_{-2}^{2} (4-x^2)\,dx = \left[4x - \frac{x^3}{3}\right]_{-2}^{2} = \left(8-\tfrac{8}{3}\right) - \left(-8+\tfrac{8}{3}\right) = \frac{32}{3}. \]
The area is $\dfrac{32}{3}$ square units.
\end{exemplebox}

\begin{popfigure}
\begin{center}
\begin{tikzpicture}
\begin{axis}[width=11cm, height=7cm, axis lines=middle, xlabel=$x$, ylabel=$y$,
xmin=-3, xmax=3, ymin=-1, ymax=5, xtick={-2,2}, ytick={4},
legend style={at={(0.5,1.02)}, anchor=south, draw=none}]
\addplot[popblue, thick, domain=-2.6:2.6, samples=100] {4-x^2};
\addlegendentry{$y=4-x^2$}
\addplot[fill=popblueL, draw=none, domain=-2:2, samples=100] {4-x^2} \closedcycle;
\node[popdark] at (axis cs:0,1.6) {$A=\dfrac{32}{3}$};
\end{axis}
\end{tikzpicture}
\end{center}
\legende{Area under $y=4-x^2$ above the $x$-axis on $[-2,2]$.}
\end{popfigure}

\courssub{Regions Below the Axis: Total Area versus Net Integral}

When $f$ changes sign on $[a,b]$, the integral $\int_a^b f(x)\,dx$ computes a \cle{net} (signed) quantity: areas above the axis count positively, areas below count negatively, and they may partly or wholly cancel. The \emph{total} area, which is what examination questions almost always ask for, is $\int_a^b |f(x)|\,dx$.

\begin{methode}[Total area when the curve crosses the axis]
\begin{enumerate}
\item Solve $f(x)=0$ on $[a,b]$ to locate the zeros $x_1 < x_2 < \cdots$.
\item Determine the sign of $f$ on each subinterval (test a point in each).
\item Integrate piece by piece, taking the \textbf{opposite} of the integral on subintervals where $f<0$.
\item Add the positive pieces: $A = \displaystyle\sum_k \left|\int_{x_{k-1}}^{x_k} f(x)\,dx\right|$.
\end{enumerate}
\end{methode}

\begin{attention}[The classic trap]
Computing $\int_a^b f(x)\,dx$ in one piece when $f$ crosses the axis gives the \emph{net} value, not the area. For $f(x)=x^3-x$ on $[-1,1]$, the single integral is $0$, but the region clearly has positive area! You must split at the zeros.
\end{attention}

\begin{exoresolu}[Total area under a cubic]
Find the total area of the region bounded by $y = x^3 - x$ and the $x$-axis on $[-1,1]$.
\tcblower
\textbf{Zeros:} $x^3 - x = x(x-1)(x+1) = 0$ gives $x = -1, 0, 1$.\\
\textbf{Signs:} for $x \in (-1,0)$, e.g.\ $x=-\tfrac12$: $f(-\tfrac12) = -\tfrac18 + \tfrac12 = \tfrac38 > 0$. For $x \in (0,1)$, e.g.\ $x=\tfrac12$: $f(\tfrac12) = \tfrac18 - \tfrac12 = -\tfrac38 < 0$.\\
\textbf{Split the integral:}
\begin{align*}
\int_{-1}^{0}(x^3-x)\,dx &= \left[\frac{x^4}{4}-\frac{x^2}{2}\right]_{-1}^{0} = 0 - \left(\tfrac14 - \tfrac12\right) = \frac14,\\
\int_{0}^{1}(x^3-x)\,dx &= \left[\frac{x^4}{4}-\frac{x^2}{2}\right]_{0}^{1} = \frac14 - \frac12 = -\frac14.
\end{align*}
\textbf{Total area:} $A = \left|\tfrac14\right| + \left|-\tfrac14\right| = \dfrac12$ square units. (The net integral is $\tfrac14 - \tfrac14 = 0$, which is \emph{not} the area.)
\end{exoresolu}

\begin{popfigure}
\begin{center}
\begin{tikzpicture}
\begin{axis}[width=11cm, height=7cm, axis lines=middle, xlabel=$x$, ylabel=$y$,
xmin=-1.6, xmax=1.6, ymin=-0.8, ymax=0.8, xtick={-1,0,1}, ytick=\empty,
legend style={at={(0.5,1.02)}, anchor=south, draw=none}]
\addplot[poporange, thick, domain=-1.35:1.35, samples=150] {x^3-x};
\addlegendentry{$y=x^3-x$}
\addplot[fill=popgreenL, draw=none, domain=-1:0, samples=100] {x^3-x} \closedcycle;
\addplot[fill=popblueL, draw=none, domain=0:1, samples=100] {x^3-x} \closedcycle;
\node[popdark] at (axis cs:-0.55,0.18) {$+\tfrac14$};
\node[popdark] at (axis cs:0.55,-0.18) {$-\tfrac14$};
\end{axis}
\end{tikzpicture}
\end{center}
\legende{The two lobes of $y=x^3-x$: the integral over $[-1,1]$ is $0$, but the total area is $\tfrac12$.}
\end{popfigure}

\courssub{Area Between Two Curves}

Let $f$ and $g$ be continuous on $[a,b]$ with $f(x) \geq g(x)$ there. Think of vertical strips of width $dx$: each strip runs from the lower curve $g(x)$ up to the upper curve $f(x)$, so its height is $f(x)-g(x)$ and its area is $\bigl(f(x)-g(x)\bigr)dx$. Summing and passing to the limit:

\begin{propriete}[Area between two curves]
If $f(x) \geq g(x)$ on $[a,b]$, the area of the region between the curves is
\[ A = \int_a^b \bigl(f(x) - g(x)\bigr)\,dx. \]
If the curves cross inside $[a,b]$, split at the intersection points and integrate $|f(x)-g(x)|$ on each piece. The formula works even if one or both curves dip below the $x$-axis, because the \emph{difference} of heights is what matters.
\end{propriete}

\begin{methode}[Area between two curves — exam procedure]
\begin{enumerate}
\item \textbf{Intersect:} solve $f(x)=g(x)$ to find the boundary $x$-values.
\item \textbf{Decide which is on top:} test one point between consecutive intersections.
\item \textbf{Integrate the difference} (top minus bottom) on each interval.
\item \textbf{Add} the positive contributions.
\end{enumerate}
\end{methode}

\begin{exoresolu}[Between a parabola and a line]
Find the area of the region enclosed by $y = x^2$ and $y = x$.
\tcblower
\textbf{Intersection:} $x^2 = x \Rightarrow x(x-1)=0 \Rightarrow x=0$ or $x=1$, giving the points $(0,0)$ and $(1,1)$.\\
\textbf{Which is on top?} At $x=\tfrac12$: $x = \tfrac12 > x^2 = \tfrac14$, so the line is above on $[0,1]$.\\
\textbf{Integrate:}
\[ A = \int_0^1 (x - x^2)\,dx = \left[\frac{x^2}{2} - \frac{x^3}{3}\right]_0^1 = \frac12 - \frac13 = \frac16. \]
The area is $\dfrac16$ square units.
\end{exoresolu}

\begin{popfigure}
\begin{center}
\begin{tikzpicture}[scale=4]
\fill[popblueL] plot[domain=0:1, samples=60] (\x,{\x}) -- plot[domain=1:0, samples=60] (\x,{\x*\x}) -- cycle;
\draw[->] (-0.15,0) -- (1.25,0) node[right] {$x$};
\draw[->] (0,-0.15) -- (0,1.25) node[above] {$y$};
\draw[popblue, thick, domain=-0.1:1.15, samples=60] plot (\x,{\x*\x}) node[right] {$y=x^2$};
\draw[poporange, thick, domain=-0.1:1.15] plot (\x,{\x}) node[right, yshift=2pt] {$y=x$};
\fill[popdark] (0,0) circle (0.015) node[below left] {$(0,0)$};
\fill[popdark] (1,1) circle (0.015) node[above right] {$(1,1)$};
\node[popdark] at (0.55,0.4) {$A=\tfrac16$};
\end{tikzpicture}
\end{center}
\legende{Region between $y=x$ (top) and $y=x^2$ (bottom), with intersections $(0,0)$ and $(1,1)$.}
\end{popfigure}

\courssub{Area Bounded by the $y$-axis: Integrating with Respect to $y$}

Some regions are described far more naturally by expressing $x$ in terms of $y$. Consider the curve $x = y^2$ (a parabola opening to the right) together with the $y$-axis. Solving for $y$ would give two branches $y = \pm\sqrt{x}$, which is clumsy. Instead, use \cle{horizontal strips}: a strip at height $y$, of thickness $dy$, runs from the $y$-axis ($x=0$) to the curve ($x = g(y)$), so its area is $g(y)\,dy$.

\begin{propriete}[Area between a curve and the $y$-axis]
If $x = g(y)$ is continuous with $g(y) \geq 0$ for $c \leq y \leq d$, the area bounded by the curve, the $y$-axis, and the lines $y=c$, $y=d$ is
\[ A = \int_c^d g(y)\,dy. \]
More generally, the area between $x = g_2(y)$ (right) and $x = g_1(y)$ (left) is $\displaystyle\int_c^d \bigl(g_2(y)-g_1(y)\bigr)\,dy$.
\end{propriete}

\begin{methode}[Integrating with respect to $y$]
\begin{enumerate}
\item Rewrite the boundary as $x$ in terms of $y$.
\item Find the $y$-coordinates of the intersection points (these become the limits $c$ and $d$).
\item Identify the right-hand and left-hand boundaries.
\item Compute $A = \int_c^d (\text{right} - \text{left})\,dy$.
\end{enumerate}
\end{methode}

\begin{exoresolu}[Region bounded by $x=y^2$ and the $y$-axis]
Find the area of the region bounded by $x = y^2$, the $y$-axis, and the lines $y=0$ and $y=2$.
\tcblower
Here $g(y) = y^2 \geq 0$ on $[0,2]$, and the left boundary is the $y$-axis ($x=0$). Hence
\[ A = \int_0^2 y^2\,dy = \left[\frac{y^3}{3}\right]_0^2 = \frac{8}{3}. \]
The area is $\dfrac{8}{3}$ square units. \emph{Check with vertical strips:} the region is $\{(x,y): 0\le x\le 4,\ \sqrt{x}\le y\le 2\}$, so $A = \int_0^4 (2-\sqrt{x})\,dx = \left[2x - \tfrac{2}{3}x^{3/2}\right]_0^4 = 8 - \tfrac{16}{3} = \tfrac83$. Same answer — but the $y$-integration was one line.
\end{exoresolu}

\begin{popfigure}
\begin{center}
\begin{tikzpicture}[scale=1.1]
\fill[popgreenL] plot[domain=0:2, samples=60] ({\x*\x},\x) -- (0,2) -- (0,0) -- cycle;
\draw[->] (-0.4,0) -- (4.8,0) node[right] {$x$};
\draw[->] (0,-0.4) -- (0,2.6) node[above] {$y$};
\draw[popblue, thick, domain=0:2.25, samples=80] plot ({\x*\x},\x) node[right] {$x=y^2$};
\draw[popdark, dashed] (0,2) node[left] {$2$} -- (4,2);
\fill[poporange] (2.25,1.5) rectangle (4,1.62);
\node[popdark] at (3.1,1.85) {strip: $y^2\,dy$};
\node[popdark] at (1.3,0.6) {$A=\tfrac83$};
\end{tikzpicture}
\end{center}
\legende{Region bounded by $x=y^2$ and the $y$-axis: a horizontal strip has length $y^2$ and thickness $dy$.}
\end{popfigure}

\textbf{Choosing the variable of integration.}\par
As the check above shows, the same region can often be attacked with vertical strips ($dx$) or horizontal strips ($dy$). Choose the variable that (i) avoids splitting the region into pieces, and (ii) avoids awkward inverse functions. If the boundary is given as $x$ in terms of $y$, integrate with respect to $y$; if one formula $y=f(x)$ describes the whole upper boundary, vertical strips are natural.

\begin{savaistu}[Historical note: From exhaustion to Riemann]
The method of exhaustion used by Archimedes (3rd century BC) to find the area of a parabolic segment is the direct ancestor of Riemann sums. Archimedes inscribed triangles in the segment, summed their areas, and took the limit — essentially computing $\int_0^1 x^2\,dx = \frac13$ without algebraic notation. The modern $\int$ symbol, introduced by Leibniz in 1675, is an elongated `S' for \emph{summa} (sum), reminding us that every definite integral is fundamentally a limit of sums.
\end{savaistu}

\courssub{Areas Involving Trigonometric, Exponential and Logarithmic Curves}

GCE questions frequently combine area calculations with the standard integrals of $\sin$, $\cos$, $e^x$ and $\ln x$. The technique is unchanged: sketch, locate intersections and zeros, split if necessary, integrate.

\begin{exoresolu}[Area under a sine arch and an exponential region]
\begin{enumerate}
\item[(a)] Find the area of one arch of $y = \sin x$.
\item[(b)] Find the area bounded by $y = e^x$, the $x$-axis, the $y$-axis and the line $x = 1$.
\end{enumerate}
\tcblower
\textbf{(a)} One arch corresponds to $0 \leq x \leq \pi$, where $\sin x \geq 0$:
\[ A = \int_0^{\pi} \sin x\,dx = \bigl[-\cos x\bigr]_0^{\pi} = -(-1) - (-1) = 2 \text{ square units}. \]
\textbf{(b)} $e^x > 0$ everywhere, so
\[ A = \int_0^1 e^x\,dx = \bigl[e^x\bigr]_0^1 = e - 1 \approx 1.718 \text{ square units}. \]
\end{exoresolu}

\begin{exoresolu}[A logarithmic region — integrate with respect to $y$]
Find the area of the region bounded by $y = \ln x$, the $x$-axis and the line $x = e$.
\tcblower
\textbf{Vertical strips:} $\ln x \geq 0$ on $[1,e]$, and $\ln x = 0$ at $x=1$:
\[ A = \int_1^e \ln x\,dx = \bigl[x\ln x - x\bigr]_1^e = (e\cdot 1 - e) - (0 - 1) = 1 \text{ square unit}. \]
\textbf{Horizontal strips (alternative):} write $x = e^y$; the region lies between $x = e^y$ (left) and $x = e$ (right) for $0 \leq y \leq 1$:
\[ A = \int_0^1 \left(e - e^y\right)dy = \left[ey - e^y\right]_0^1 = (e - e) - (0 - 1) = 1. \]
Both routes agree, as they must.
\end{exoresolu}

\courssub{Examination Technique}

\begin{methode}[Sketch, shade, justify]
\begin{enumerate}
\item \textbf{Sketch} both curves (or curve and axis) on the same axes, marking all intersection points found algebraically.
\item \textbf{Shade} the exact region whose area is required — this fixes the limits and the order (top minus bottom, or right minus left).
\item \textbf{Justify the limits:} state that the limits are the $x$-coordinates (or $y$-coordinates) of the intersections or of the given boundary lines.
\item \textbf{Check the sign} of your answer: an area must be positive; if you obtain a negative number, you integrated bottom minus top or forgot to split at a zero.
\item \textbf{State the units:} ``square units'' when no scale is given.
\end{enumerate}
\end{methode}

\begin{attention}[Frequent errors at the GCE]
\begin{itemize}
\item Integrating across a zero of $f$ without splitting — you get the net signed value.
\item Using the $x$-coordinates of intersections as limits when integrating with respect to $y$ — the limits must be $y$-values.
\item Subtracting in the wrong order and reporting a negative area.
\item Forgetting that $\int \ln x\,dx = x\ln x - x$ (integration by parts result), not $\tfrac{1}{x}$.
\end{itemize}
\end{attention}

\begin{exercice}[Consolidation]
\begin{enumerate}
\item Find the total area bounded by $y = x^2 - 4$ and the $x$-axis.
\item Find the area enclosed by $y = x^2$ and $y = 2x$.
\item Find the area bounded by $y = \cos x$, the $x$-axis, and the lines $x=0$, $x=\pi$.
\item Find the area bounded by $x = 4 - y^2$ and the $y$-axis.
\item Find the area enclosed by $y = e^x$, $y = e^{-x}$ and the line $x = 1$.
\end{enumerate}
\end{exercice}

\begin{corrige}[Solutions to Consolidation Exercises]
\textbf{1.} Zeros at $x=\pm2$; $x^2-4 \le 0$ on $[-2,2]$, so $A = -\int_{-2}^{2}(x^2-4)\,dx$. Compute the integral:
\[ \int_{-2}^{2}(x^2-4)\,dx = \left[\frac{x^3}{3}-4x\right]_{-2}^{2} = \left(\frac{8}{3}-8\right) - \left(-\frac{8}{3}+8\right) = \frac{16}{3} - 16 = -\frac{32}{3}. \]
Hence $A = \frac{32}{3}$ square units.\\
\textbf{2.} Intersections: $x^2=2x \Rightarrow x=0,2$. Line on top: $A = \int_0^2(2x-x^2)dx = \left[x^2-\frac{x^3}{3}\right]_0^2 = 4-\frac{8}{3} = \frac{4}{3}$ square units.\\
\textbf{3.} $\cos x \ge 0$ on $[0,\frac{\pi}{2}]$, $\cos x \le 0$ on $[\frac{\pi}{2},\pi]$. Split at $\frac{\pi}{2}$:
\[ A = \int_0^{\pi/2}\cos x\,dx - \int_{\pi/2}^{\pi}\cos x\,dx = \bigl[\sin x\bigr]_0^{\pi/2} - \bigl[\sin x\bigr]_{\pi/2}^{\pi} = 1 - (-1) = 2 \text{ square units}. \]
\textbf{4.} $x=4-y^2 \ge 0$ for $-2 \le y \le 2$: $A = \int_{-2}^{2}(4-y^2)dy = \left[4y-\frac{y^3}{3}\right]_{-2}^{2} = \left(8-\frac{8}{3}\right) - \left(-8+\frac{8}{3}\right) = \frac{32}{3}$ square units.\\
\textbf{5.} On $[0,1]$, $e^x \ge e^{-x}$; curves meet at $x=0$. $A = \int_0^1 (e^x - e^{-x})dx = \left[e^x + e^{-x}\right]_0^1 = (e + e^{-1}) - (1+1) = e + e^{-1} - 2 \approx 1.086$ square units.
\end{corrige}

\begin{aretenir}[Key points of this section]
\begin{itemize}
\item $f \ge 0$ on $[a,b]$: $A = \int_a^b f\,dx$; if $f \le 0$, $A = -\int_a^b f\,dx$.
\item Total area $= \int_a^b |f|\,dx$: split at every zero of $f$.
\item Between two curves: $A = \int_a^b (\text{top} - \text{bottom})\,dx$, splitting where the curves cross.
\item With respect to $y$: $A = \int_c^d (\text{right} - \text{left})\,dy$, with $y$-limits.
\item Always sketch, shade, and check that the final area is positive.
\end{itemize}
\end{aretenir}

\coursec{Volumes of Solids of Revolution}

In the previous section we used the definite integral to add up infinitely many thin strips of \emph{area} and so compute the area of a plane region. We now push the same idea one dimension further: we take a plane region, \cle{rotate} it about an axis, and ask for the \emph{volume} of the solid generated. The definite integral will again do the adding for us — this time of thin \emph{discs} instead of thin strips. This is one of the most frequently examined applications of integration at the GCE Advanced Level, and it rewards careful, methodical work.

\courssub{The Idea of the Disc Method}

\begin{experience}[From a region to a solid]
Take a piece of cardboard cut in the shape of the region under the curve $y=\sqrt{x}$ between $x=0$ and $x=4$. Pin it along the $x$-axis with a knitting needle and spin the needle rapidly. The cardboard sweeps out a solid that looks like a trumpet bell or a drinking gourd. Every point of the region traces a circle, and the whole region traces a solid of \cle{revolution}.
\end{experience}

How can we compute the volume of such a solid? We imitate exactly what we did for areas in Lesson 52. Slice the solid by planes \emph{perpendicular to the axis of rotation}. Each slice is (very nearly) a thin circular \cle{disc} — a cylinder of small thickness $\delta x$. The volume of a cylinder is $\pi r^{2}h$, so:

\begin{itemize}
\item the radius of the disc at position $x$ is the distance from the curve to the axis, namely $r = y = f(x)$;
\item the thickness of the disc is a small increment $\delta x$;
\item the volume of one disc is therefore approximately $\pi [f(x)]^{2}\,\delta x$.
\end{itemize}

Adding the volumes of all the discs from $x=a$ to $x=b$ and taking the limit as $\delta x \to 0$ (a Riemann sum, exactly as in Lesson 52) gives the exact volume:
\[ V = \lim_{\delta x \to 0} \sum_{x=a}^{b} \pi y^{2}\,\delta x = \pi\int_{a}^{b} y^{2}\,dx. \]

\begin{popfigure}
\begin{center}
\begin{tikzpicture}[scale=1.05]
% axis
\draw[->, popmuted] (-0.4,0) -- (5.2,0) node[right] {$x$};
\draw[->, popmuted] (0,-0.6) -- (0,2.6) node[above] {$y$};
% curve
\draw[very thick, popblue, domain=0:4.6, samples=80] plot (\x,{sqrt(max(0,\x))}) node[above right, black] {$y=\sqrt{x}$};
% back half of solid (dashed suggestion of 3D)
\draw[popblue, dashed, domain=0.05:4.6, samples=60] plot (\x,{-sqrt(max(0,\x))});
% ellipses at ends
\draw[popblue] (4.6,0) ellipse (0.14 and 2.144);
\draw[popblue] (0.05,0) ellipse (0.03 and 0.22);
% one disc at x=3
\fill[poporange, opacity=0.35] (3,0) ellipse (0.13 and 1.732);
\draw[poporange, thick] (3,0) ellipse (0.13 and 1.732);
\draw[poporange, thick] (3,0) -- (3,1.732) node[midway, right, popdark] {$y$};
\draw[<->, popdark] (3,-0.5) -- (3.35,-0.5) node[midway, below] {$\delta x$};
\draw[poporange, thick, opacity=0.6] (3.35,0) ellipse (0.13 and 1.83);
\fill[poporange, opacity=0.2] (3.35,0) ellipse (0.13 and 1.83);
\draw[popmuted] (3,1.732) -- (3.35,1.83);
\draw[popmuted] (3,-1.732) -- (3.35,-1.83);
% labels
\node[below, popmuted] at (0,0) {$O$};
\node[below, popmuted] at (4.6,-0.15) {$b$};
\end{tikzpicture}
\end{center}
\legende{The region under $y=\sqrt{x}$ rotated about the $x$-axis. A typical slice is a disc of radius $y$ and thickness $\delta x$, of volume $\pi y^{2}\,\delta x$.}
\end{popfigure}

\courssub{Rotation About the $x$-Axis (Lesson 56)}

\begin{propriete}[Volume of revolution about the $x$-axis]
The volume of the solid generated when the region under the curve $y=f(x)$, above the $x$-axis, between $x=a$ and $x=b$, is rotated through one complete turn ($2\pi$ radians) about the $x$-axis is
\[ V=\pi\int_{a}^{b} y^{2}\,dx=\pi\int_{a}^{b}\bigl[f(x)\bigr]^{2}\,dx. \]
The derivation from the Riemann sum of disc volumes is examinable: $V=\lim\limits_{\delta x\to 0}\sum \pi y^{2}\,\delta x=\pi\displaystyle\int_{a}^{b}y^{2}\,dx$.
\end{propriete}

\begin{attention}[Square the radius!]
The integrand is $\bigl[f(x)\bigr]^{2}$, \textbf{not} $f(x)$. Forgetting to square is the single most common error in this topic. Remember: a disc of radius $y$ has \emph{area} $\pi y^{2}$, and it is that area which is integrated along the axis. Also note the factor $\pi$ multiplies the \emph{whole} integral.
\end{attention}

\begin{methode}[The disc method about the $x$-axis, step by step]
\begin{enumerate}
\item \textbf{Sketch} the region and the axis of rotation; identify the limits $a$ and $b$ (these are $x$-values).
\item \textbf{Identify the radius}: the perpendicular distance from the curve to the axis, $r=y=f(x)$.
\item \textbf{Square} the radius and write down $V=\pi\displaystyle\int_{a}^{b}[f(x)]^{2}\,dx$.
\item \textbf{Integrate} $[f(x)]^{2}$ with respect to $x$ (expand brackets first if necessary).
\item \textbf{Evaluate} between the limits and multiply by $\pi$; give the exact answer in terms of $\pi$, then a decimal approximation if required. State units (e.g.\ $\mathrm{units}^{3}$).
\end{enumerate}
\end{methode}

\begin{exoresolu}[Volume of a cone — standard derivation]
The line segment $y=\dfrac{r}{h}x$ between $x=0$ and $x=h$ is rotated about the $x$-axis. Show that the volume of the cone generated is $V=\dfrac{1}{3}\pi r^{2}h$.
\tcblower
\textbf{Solution.} The radius of the disc at position $x$ is $y=\dfrac{r}{h}x$. Hence
\begin{align*}
V&=\pi\int_{0}^{h}\left(\frac{r}{h}x\right)^{2}dx
=\pi\,\frac{r^{2}}{h^{2}}\int_{0}^{h}x^{2}\,dx
=\pi\,\frac{r^{2}}{h^{2}}\left[\frac{x^{3}}{3}\right]_{0}^{h}\\
&=\pi\,\frac{r^{2}}{h^{2}}\cdot\frac{h^{3}}{3}
=\frac{1}{3}\pi r^{2}h,
\end{align*}
which is the familiar formula for the volume of a cone of base radius $r$ and height $h$. This confirms that the disc method reproduces a result we already trust.
\end{exoresolu}

\begin{exoresolu}[Volume of a sphere — standard derivation]
The semicircle $y=\sqrt{a^{2}-x^{2}}$ ($-a\le x\le a$) is rotated about the $x$-axis. Show that the sphere generated has volume $V=\dfrac{4}{3}\pi a^{3}$.
\tcblower
\textbf{Solution.} Here $y^{2}=a^{2}-x^{2}$, so
\begin{align*}
V&=\pi\int_{-a}^{a}\left(a^{2}-x^{2}\right)dx
=\pi\left[a^{2}x-\frac{x^{3}}{3}\right]_{-a}^{a}\\
&=\pi\left[\left(a^{3}-\frac{a^{3}}{3}\right)-\left(-a^{3}+\frac{a^{3}}{3}\right)\right]
=\pi\left(\frac{2a^{3}}{3}+\frac{2a^{3}}{3}\right)
=\frac{4}{3}\pi a^{3}.
\end{align*}
\end{exoresolu}

\begin{popfigure}
\begin{center}
\begin{tikzpicture}[scale=0.9]
\draw[->, popmuted] (-3.4,0) -- (3.6,0) node[right] {$x$};
\draw[->, popmuted] (0,-2.2) -- (0,2.6) node[above] {$y$};
% semicircle
\draw[very thick, popblue, domain=-2:2, samples=100] plot (\x,{sqrt(max(0,4-\x*\x))});
\draw[popblue, dashed, domain=-2:2, samples=100] plot (\x,{-sqrt(max(0,4-\x*\x))});
% sphere outline
\draw[poppurple, thick] (0,0) circle (2);
\draw[poppurple] (0,0) ellipse (2 and 0.55);
\draw[poppurple, dashed] (2,0) arc (0:180:2 and 0.55);
% radius
\draw[->, popdark, thick] (0,0) -- (1.414,1.414) node[midway, above left] {$a$};
\node[below left, popmuted] at (0,0) {$O$};
\node[below, popmuted] at (2,-0.05) {$a$};
\node[below, popmuted] at (-2,-0.05) {$-a$};
\node[popblue] at (2.6,1.5) {$y=\sqrt{a^{2}-x^{2}}$};
\end{tikzpicture}
\end{center}
\legende{Rotating the semicircle $y=\sqrt{a^{2}-x^{2}}$ about the $x$-axis generates a sphere of radius $a$.}
\end{popfigure}

\begin{exoresolu}[A paraboloid]
Find the volume generated when the region under $y=\sqrt{x}$ from $x=0$ to $x=4$ is rotated completely about the $x$-axis.
\tcblower
\textbf{Solution.} Since $y^{2}=x$,
\[ V=\pi\int_{0}^{4}x\,dx=\pi\left[\frac{x^{2}}{2}\right]_{0}^{4}=\pi\left(\frac{16}{2}-0\right)=8\pi. \]
The volume is $8\pi$ cubic units $\approx 25.13$ cubic units.
\end{exoresolu}

\courssub{Rotation About the $y$-Axis (Lesson 57)}

If instead the region is rotated about the $y$-axis, the discs are stacked \emph{vertically}. The radius of a disc is now a horizontal distance, $x$, and the thickness is $\delta y$. Interchanging the roles of $x$ and $y$ in the previous argument:

\begin{propriete}[Volume of revolution about the $y$-axis]
The volume generated when the region bounded by the curve $x=g(y)$, the $y$-axis, and the lines $y=c$ and $y=d$ is rotated through one complete turn about the $y$-axis is
\[ V=\pi\int_{c}^{d} x^{2}\,dy=\pi\int_{c}^{d}\bigl[g(y)\bigr]^{2}\,dy. \]
\end{propriete}

\begin{attention}[Integrate with respect to $y$]
When rotating about the $y$-axis the integral must be with respect to $y$. If the curve is given as $y=f(x)$, you must first \textbf{rewrite $x$ in terms of $y$} (i.e.\ find $x=g(y)$), and the limits $c,d$ are \textbf{$y$-values}, not $x$-values. Mixing $x$ and $y$ in the same integral is a classic source of lost marks.
\end{attention}

\begin{methode}[The disc method about the $y$-axis, step by step]
\begin{enumerate}
\item \textbf{Sketch} the region and the axis of rotation; identify the limits $c$ and $d$ (these are $y$-values).
\item \textbf{Express $x$ in terms of $y$}: rearrange $y=f(x)$ to get $x=g(y)$. The radius is $r=x=g(y)$.
\item \textbf{Square} the radius and write down $V=\pi\displaystyle\int_{c}^{d}[g(y)]^{2}\,dy$.
\item \textbf{Integrate} with respect to $y$.
\item \textbf{Evaluate} between the $y$-limits and multiply by $\pi$; give the exact answer in terms of $\pi$.
\end{enumerate}
\end{methode}

\begin{exoresolu}[Rotation about the $y$-axis]
The region bounded by the curve $y=x^{2}$, the $y$-axis and the line $y=4$ is rotated about the $y$-axis. Find the volume generated.
\tcblower
\textbf{Solution.} From $y=x^{2}$ we get $x^{2}=y$ (we need $x^{2}$ directly — no square root is even necessary). The limits are $y=0$ (where the curve meets the $y$-axis) to $y=4$. Hence
\[ V=\pi\int_{0}^{4}x^{2}\,dy=\pi\int_{0}^{4}y\,dy=\pi\left[\frac{y^{2}}{2}\right]_{0}^{4}=8\pi. \]
The volume is $8\pi$ cubic units. Note how the same region rotated about the two different axes can give quite different solids.
\end{exoresolu}

\courssub{Hollow Solids: the Washer Method (Extension for Regions Between Two Curves)}

GCE questions often ask for the volume generated by a region bounded by \emph{two} curves. Suppose the region rotated lies \emph{between two curves}, $y=f(x)$ (outer) and $y=g(x)$ (inner), with $f(x)\ge g(x)\ge 0$ on $[a,b]$. When this region spins about the $x$-axis, each slice is no longer a solid disc but a \cle{washer} (an annulus): a disc of outer radius $R=f(x)$ with a hole of inner radius $r=g(x)$. Its cross-sectional area is $\pi R^{2}-\pi r^{2}$, so:

\begin{methode}[Washer method for the region between two curves (about $x$-axis)]
\begin{enumerate}
\item Sketch both curves and find their intersections to fix the limits $a$ and $b$.
\item Identify the \textbf{outer radius} $R$ (curve farther from the axis) and the \textbf{inner radius} $r$ (curve nearer the axis).
\item Write $\displaystyle V=\pi\int_{a}^{b}\left(R^{2}-r^{2}\right)dx=\pi\int_{a}^{b}\Bigl([f(x)]^{2}-[g(x)]^{2}\Bigr)dx$.
\item Integrate and evaluate; keep the answer exact in terms of $\pi$.
\end{enumerate}
Equivalently: volume of the outer solid \emph{minus} volume of the inner solid.
\end{methode}

\begin{attention}[$R^{2}-r^{2}$ is not $(R-r)^{2}$]
A very frequent mistake is to integrate $[f(x)-g(x)]^{2}$. This is wrong: $\pi(R^{2}-r^{2})\neq\pi(R-r)^{2}$. Subtract the \emph{volumes} (or the squared radii), never square the difference of the radii.
\end{attention}

\begin{popfigure}
\begin{center}
\begin{tikzpicture}[scale=2.4]
\draw[->, popmuted] (-0.3,0) -- (1.5,0) node[right] {$x$};
\draw[->, popmuted] (0,-0.3) -- (0,1.4) node[above] {$y$};
% curves
\draw[very thick, popblue, domain=0:1.15, samples=60] plot (\x,{\x}) node[right, black] {$y=x$};
\draw[very thick, popgreen, domain=0:1.1, samples=60] plot (\x,{\x*\x}) node[right, black] {$y=x^{2}$};
% shaded region
\fill[popblueL, opacity=0.6, domain=0:1, samples=60] plot (\x,{\x}) -- plot[domain=1:0, samples=60] (\x,{\x*\x}) -- cycle;
% washer at x=0.6 (drawn as annulus suggestion)
\fill[poporange, opacity=0.45] (0.6,0) ellipse (0.05 and 0.6);
\draw[poporange, thick] (0.6,0) ellipse (0.05 and 0.6);
\draw[poporange, thick] (0.6,0) ellipse (0.05 and 0.36);
\fill[white] (0.6,0) ellipse (0.05 and 0.36);
\draw[poporange, thick] (0.6,0.36) -- (0.6,0.6);
\node[popdark, scale=0.85] at (0.82,0.62) {$R=x$};
\node[popdark, scale=0.85] at (0.85,0.28) {$r=x^{2}$};
\draw[popmuted, dashed] (0.6,0) -- (0.6,0.36);
\node[below, popmuted] at (1,0) {$1$};
\node[below left, popmuted] at (0,0) {$O$};
\end{tikzpicture}
\end{center}
\legende{The region between $y=x$ and $y=x^{2}$ on $[0,1]$. Rotated about the $x$-axis, a slice gives a washer of outer radius $R=x$ and inner radius $r=x^{2}$.}
\end{popfigure}

\begin{exoresolu}[Washer method in action]
The region enclosed by $y=x$ and $y=x^{2}$ is rotated completely about the $x$-axis. Find the volume generated.
\tcblower
\textbf{Solution.} The curves meet where $x=x^{2}$, i.e.\ $x=0$ and $x=1$. On $[0,1]$, $x\ge x^{2}$, so the outer radius is $R=x$ and the inner radius is $r=x^{2}$. Then
\begin{align*}
V&=\pi\int_{0}^{1}\Bigl(x^{2}-\bigl(x^{2}\bigr)^{2}\Bigr)dx
=\pi\int_{0}^{1}\bigl(x^{2}-x^{4}\bigr)\,dx\\
&=\pi\left[\frac{x^{3}}{3}-\frac{x^{5}}{5}\right]_{0}^{1}
=\pi\left(\frac{1}{3}-\frac{1}{5}\right)
=\frac{2\pi}{15}.
\end{align*}
The volume is $\dfrac{2\pi}{15}$ cubic units $\approx 0.419$ cubic units.
\end{exoresolu}

\courssub{Regions Bounded by Axes and Given Ordinates: Mixed Problems}

GCE questions frequently describe the region in words: \emph{``the region bounded by the curve $y=\ldots$, the $x$-axis and the lines $x=a$ and $x=b$''}, or \emph{``bounded by the curve, the $y$-axis and the line $y=d$''}. The essential skill is translating the description into correct limits and a correct radius \emph{before} integrating.

\begin{exoresolu}[Choosing the correct axis and limits]
The region $R$ is bounded by the curve $y=\dfrac{1}{x}$, the $x$-axis and the lines $x=1$ and $x=2$.
\begin{enumerate}
\item[(a)] Find the volume when $R$ is rotated about the $x$-axis.
\item[(b)] The region $S$ bounded by $y=\dfrac{1}{x}$, the $y$-axis, $y=\dfrac{1}{2}$ and $y=1$ is rotated about the $y$-axis. Find its volume.
\end{enumerate}
\tcblower
\textbf{Solution.} (a) About the $x$-axis, radius $y=\dfrac{1}{x}$:
\[ V=\pi\int_{1}^{2}\frac{1}{x^{2}}\,dx=\pi\left[-\frac{1}{x}\right]_{1}^{2}=\pi\left(-\frac{1}{2}+1\right)=\frac{\pi}{2}. \]
(b) About the $y$-axis we need $x$ in terms of $y$: from $y=\dfrac{1}{x}$, $x=\dfrac{1}{y}$. The limits are $y=\dfrac{1}{2}$ to $y=1$:
\[ V=\pi\int_{1/2}^{1}\frac{1}{y^{2}}\,dy=\pi\left[-\frac{1}{y}\right]_{1/2}^{1}=\pi\left(-1+2\right)=\pi. \]
Notice how the same curve produces $\dfrac{\pi}{2}$ about one axis and $\pi$ about the other: the axis of rotation changes everything.
\end{exoresolu}

\begin{aretenir}[Key points on volumes of revolution]
\begin{itemize}
\item About the $x$-axis: $\displaystyle V=\pi\int_{a}^{b}y^{2}\,dx$; about the $y$-axis: $\displaystyle V=\pi\int_{c}^{d}x^{2}\,dy$.
\item The radius is the \emph{perpendicular distance from the curve to the axis of rotation}; always \textbf{square} it.
\item For rotation about the $y$-axis, rewrite $x$ as a function of $y$ and use $y$-limits.
\item Between two curves (washer method): $V=\pi\displaystyle\int_{a}^{b}\bigl(R^{2}-r^{2}\bigr)\,dx$, never $\pi\int(R-r)^{2}\,dx$.
\item Leave answers \textbf{exact in terms of $\pi$} unless a decimal is requested; volumes carry cubic units (e.g.\ $\mathrm{cm}^{3}$ if $x,y$ are in cm).
\end{itemize}
\end{aretenir}

\begin{exercice}[Practice]
\begin{enumerate}
\item The region under $y=x^{2}+1$ from $x=0$ to $x=2$ is rotated about the $x$-axis. Find the volume generated.
\item The region bounded by $y=2\sqrt{x}$, the $x$-axis and the line $x=3$ is rotated about the $x$-axis. Find the volume.
\item The region bounded by $y=x^{3}$, the $y$-axis and the line $y=8$ is rotated about the $y$-axis. Find the volume.
\item The region enclosed by $y=x^{2}$ and $y=2x$ is rotated about the $x$-axis. Show that the volume is $\dfrac{64\pi}{15}$.
\end{enumerate}
\end{exercice}

\begin{corrige}[Exercise 1]
\begin{enumerate}
\item $V=\pi\displaystyle\int_{0}^{2}(x^{2}+1)^{2}\,dx=\pi\int_{0}^{2}(x^{4}+2x^{2}+1)\,dx=\pi\left[\dfrac{x^{5}}{5}+\dfrac{2x^{3}}{3}+x\right]_{0}^{2}=\pi\left(\dfrac{32}{5}+\dfrac{16}{3}+2\right)=\dfrac{206\pi}{15}$.
\item $V=\pi\displaystyle\int_{0}^{3}4x\,dx=\pi\left[2x^{2}\right]_{0}^{3}=18\pi$.
\item $x^{2}=y^{2/3}$; $V=\pi\displaystyle\int_{0}^{8}y^{2/3}\,dy=\pi\left[\dfrac{3}{5}y^{5/3}\right]_{0}^{8}=\pi\cdot\dfrac{3}{5}\cdot 32=\dfrac{96\pi}{5}$.
\item Intersections: $x^{2}=2x\Rightarrow x=0,2$; on $[0,2]$, $2x\ge x^{2}$. $V=\pi\displaystyle\int_{0}^{2}\left(4x^{2}-x^{4}\right)dx=\pi\left[\dfrac{4x^{3}}{3}-\dfrac{x^{5}}{5}\right]_{0}^{2}=\pi\left(\dfrac{32}{3}-\dfrac{32}{5}\right)=\dfrac{64\pi}{15}$.
\end{enumerate}
\end{corrige}

\begin{savaistu}[Cooling towers and gourds]
The graceful curved shape of a cooling tower at a power station is a solid (in fact a shell) of revolution: a hyperbola rotated about its axis, chosen because the shape resists wind load efficiently. Closer to home, the calabashes used in many Cameroonian households are natural solids of revolution — and a potter shaping a pot on a wheel is literally performing a rotation about a vertical axis, the same mathematics as $V=\pi\displaystyle\int x^{2}\,dy$.
\end{savaistu}

\coursec{Centroids and Centres of Mass by Integration}

In the previous section we used integration to add up infinitely many thin discs and so obtain the \emph{volume} of a solid of revolution. The same philosophy — slice, compute the contribution of one slice, integrate — answers another very practical question: \emph{where does an object balance?} A carpenter cutting a metal plate in Douala, an engineer designing a water tank in Bamenda, a potter shaping a vessel in Bafoussam: each needs to know the \cle{centre of mass} of the object. In this section we derive, rigorously, the integral formulas that locate it.

\courssub{From Point Masses to Continuous Bodies}

\textbf{Reminder: discrete systems.}\par
Consider point masses $m_1, m_2, \dots, m_n$ placed on a light rod along the $x$-axis at positions $x_1, x_2, \dots, x_n$. The \emph{moment} of the mass $m_i$ about the origin (more precisely about the axis through $O$ perpendicular to the rod) is $m_i x_i$: mass times signed distance. The rod balances on a knife-edge placed at
\[ \bar{x} = \frac{\sum_{i=1}^{n} m_i x_i}{\sum_{i=1}^{n} m_i}, \]
a \emph{weighted average} of the positions. For masses located at $(x_i, y_i)$ in the plane, the centre of mass coordinates are
\[ \bar{x} = \frac{\sum m_i x_i}{\sum m_i}, \qquad \bar{y} = \frac{\sum m_i y_i}{\sum m_i}. \]

\begin{definition}[Moment and centroid of a plane region]
Let $\mathcal{R}$ be a plane region of uniform density $\rho$ (mass per unit area). The \cle{moment} of $\mathcal{R}$ about an axis is the limit, as the partition becomes infinitely fine, of the sum of (mass of each element) $\times$ (signed distance of the element's own centroid from the axis). The \cle{centroid} of $\mathcal{R}$ is the point $(\bar{x}, \bar{y})$ at which the region would balance if it were cut from a thin plate of uniform density. For a uniform lamina, the density $\rho$ cancels out of the formulas, so the centroid depends only on geometry.
\end{definition}

\textbf{From discrete to continuous: region under a curve.}\par
Take the region bounded by $y = f(x) \ge 0$, the $x$-axis, and the lines $x = a$, $x = b$. Slice it into thin \emph{vertical} strips of width $\delta x$. A strip at position $x$ has height $f(x)$, so its area is approximately $f(x)\,\delta x$ and its mass is $\rho f(x)\,\delta x$. Two moments are needed.

\begin{itemize}
\item \textbf{Moment about the $y$-axis.} Every point of the strip is at distance $x$ from the $y$-axis, so the strip's moment is approximately $x \cdot \rho f(x)\,\delta x$. Summing and passing to the limit:
\[ M_y = \rho \int_a^b x\,f(x)\,dx. \]
\item \textbf{Moment about the $x$-axis.} The strip is a thin rectangle; its \emph{own} centroid is halfway up, at height $\tfrac{1}{2}f(x)$. Hence the strip's moment about the $x$-axis is approximately $\tfrac{1}{2}f(x) \cdot \rho f(x)\,\delta x = \tfrac{1}{2}\rho [f(x)]^2\,\delta x$, giving
\[ M_x = \frac{\rho}{2}\int_a^b [f(x)]^2\,dx. \]
\end{itemize}

The total mass is $M = \rho A$ where $A = \displaystyle\int_a^b f(x)\,dx$ is the area. The centroid coordinates are
\[ \bar{x} = \frac{M_y}{M} = \frac{1}{A}\int_a^b x\,f(x)\,dx, \qquad
\bar{y} = \frac{M_x}{M} = \frac{1}{2A}\int_a^b [f(x)]^2\,dx. \]

\begin{attention}[The factor $\tfrac{1}{2}$ in $\bar{y}$]
A very common error is to write $M_x = \int_a^b y\,f(x)\,dx$ or to forget the $\tfrac12$. Remember: for the moment about the $x$-axis, the relevant distance is the height of the \emph{strip's own centroid}, which is $\tfrac{1}{2}f(x)$, not $f(x)$. That is exactly where the factor $\tfrac12$ comes from.
\end{attention}

\begin{popfigure}
\begin{center}
\begin{tikzpicture}[scale=1.1]
\draw[->] (-0.4,0) -- (4.6,0) node[right] {$x$};
\draw[->] (0,-0.4) -- (0,3.2) node[above] {$y$};
\draw[domain=0:4, smooth, thick, popblue] plot (\x, {0.25*\x*(4-\x)+0.4});
\fill[popblueL, opacity=0.6] (0,0.4) -- plot[domain=0:4] (\x, {0.25*\x*(4-\x)+0.4}) -- (4,0) -- (0,0) -- cycle;
\draw[domain=0:4, smooth, thick, popblue] plot (\x, {0.25*\x*(4-\x)+0.4});
\fill[poporange, opacity=0.75] (2.4,0) rectangle (2.75,1.84);
\draw[popdark] (2.4,0) rectangle (2.75,1.84);
\fill[popink] (2.575,0.92) circle (0.05);
\node[popink, right] at (2.75,0.92) {\scriptsize $(x,\ \tfrac12 f(x))$};
\draw[<->, popmuted] (2.4,-0.25) -- (2.75,-0.25) node[midway, below] {\scriptsize $\delta x$};
\node[popdark] at (1.1,2.6) {\scriptsize $y = f(x)$};
\end{tikzpicture}
\end{center}
\legende{A vertical strip of height $f(x)$ under a curve; its own centroid is at $\left(x, \tfrac12 f(x)\right)$.}
\end{popfigure}

\begin{propriete}[Centroid of the region under a curve]
For the region bounded by $y = f(x) \ge 0$, the $x$-axis and the lines $x = a$, $x = b$, with area $A = \displaystyle\int_a^b f(x)\,dx$:
\[ \bar{x} = \frac{1}{A}\int_a^b x\,f(x)\,dx, \qquad \bar{y} = \frac{1}{2A}\int_a^b [f(x)]^2\,dx. \]
\end{propriete}

\begin{propriete}[Centroid of a region between two curves]
If the region lies between $y = f(x)$ (top) and $y = g(x)$ (bottom), with $f(x) \ge g(x)$ for $a \le x \le b$, the strip runs from $g(x)$ to $f(x)$. Its height is $f(x)-g(x)$, its own centroid is at height $\tfrac{1}{2}[f(x)+g(x)]$, and its area is $[f(x)-g(x)]\,\delta x$. The moment about the $x$-axis is therefore $\tfrac{1}{2}[f(x)+g(x)][f(x)-g(x)]\,\delta x = \tfrac{1}{2}\bigl([f(x)]^2 - [g(x)]^2\bigr)\,\delta x$. With area $A = \displaystyle\int_a^b [f(x)-g(x)]\,dx$:
\[ \bar{x} = \frac{1}{A}\int_a^b x\,[f(x) - g(x)]\,dx, \qquad
\bar{y} = \frac{1}{2A}\int_a^b \left([f(x)]^2 - [g(x)]^2\right)dx. \]
\end{propriete}

\courssub{Centroids of Solids of Revolution}

When a plane region is rotated about an axis, the resulting solid has its centre of mass on that axis (by symmetry). We derive the formula for rotation about the $x$-axis; the $y$-axis case is analogous.

Consider the region under $y = f(x) \ge 0$, $a \le x \le b$, rotated about the $x$-axis. Slice the solid into thin discs of thickness $\delta x$ at position $x$. A disc has radius $f(x)$, volume $\pi [f(x)]^2\,\delta x$, and mass $\rho \pi [f(x)]^2\,\delta x$ (uniform density $\rho$). Its own centroid is at its centre, i.e. at $(x,0,0)$. The moment of this disc about the $yz$-plane (the plane $x=0$) is $x \cdot \rho \pi [f(x)]^2\,\delta x$. Summing and passing to the limit:
\[ M_{yz} = \rho \pi \int_a^b x\,[f(x)]^2\,dx. \]
The total mass is $M = \rho V = \rho \pi \displaystyle\int_a^b [f(x)]^2\,dx$. Hence the $x$-coordinate of the centre of mass is
\[ \bar{x} = \frac{M_{yz}}{M} = \frac{\displaystyle\int_a^b x\,[f(x)]^2\,dx}{\displaystyle\int_a^b [f(x)]^2\,dx}. \]
By symmetry, $\bar{y} = \bar{z} = 0$.

\begin{propriete}[Centre of mass of a solid of revolution about the $x$-axis]
For the solid generated by rotating the region under $y = f(x) \ge 0$, $a \le x \le b$, about the $x$-axis (uniform density):
\[ \bar{x} = \frac{\displaystyle\int_a^b x\,y^2\,dx}{\displaystyle\int_a^b y^2\,dx}, \qquad \bar{y} = 0. \]
For rotation about the $y$-axis (region bounded by $x = g(y)$, $c \le y \le d$):
\[ \bar{y} = \frac{\displaystyle\int_c^d y\,x^2\,dy}{\displaystyle\int_c^d x^2\,dy}, \qquad \bar{x} = 0. \]
\end{propriete}

\begin{methode}[Exploiting symmetry]
A centroid (or centre of mass) lies on \emph{every} axis of symmetry of the region or solid. Therefore:
\begin{enumerate}
\item Before computing anything, look for axes of symmetry.
\item Each axis of symmetry gives one coordinate of the centroid \emph{for free} (e.g. $\bar{x} = 0$ for a region symmetric about the $y$-axis).
\item Only the remaining coordinate(s) need integration.
\end{enumerate}
This halves the work for semicircles, cones, parabolas symmetric about an axis, etc.
\end{methode}

\courssub{Worked Examples: Regions Under Curves}

\begin{exoresolu}[Centroid under a parabola]
Find the centroid of the region bounded by $y = 4 - x^2$, the $x$-axis and the $y$-axis in the first quadrant (so $0 \le x \le 2$).
\tcblower
\textbf{Area:}
\[ A = \int_0^2 (4 - x^2)\,dx = \left[4x - \frac{x^3}{3}\right]_0^2 = 8 - \frac{8}{3} = \frac{16}{3}. \]
\textbf{Moment about the $y$-axis:}
\[ \int_0^2 x(4 - x^2)\,dx = \left[2x^2 - \frac{x^4}{4}\right]_0^2 = 8 - 4 = 4
\quad\Rightarrow\quad \bar{x} = \frac{4}{16/3} = \frac{3}{4}. \]
\textbf{Moment about the $x$-axis:}
\[ \frac12\int_0^2 (4-x^2)^2\,dx = \frac12\int_0^2 (16 - 8x^2 + x^4)\,dx
= \frac12\left[16x - \frac{8x^3}{3} + \frac{x^5}{5}\right]_0^2 = \frac12\left(32 - \frac{64}{3} + \frac{32}{5}\right) = \frac{128}{15}, \]
so $\bar{y} = \dfrac{128/15}{16/3} = \dfrac{8}{5}$. The centroid is $\left(\tfrac34, \tfrac85\right)$.
\end{exoresolu}

\begin{exoresolu}[Centroid of a region between two curves]
Find the centroid of the region bounded by $y = x^2$ and $y = x$ in the first quadrant.
\tcblower
The curves intersect at $x=0$ and $x=1$. Top curve: $f(x)=x$, bottom: $g(x)=x^2$.
\textbf{Area:}
\[ A = \int_0^1 (x - x^2)\,dx = \left[\frac{x^2}{2} - \frac{x^3}{3}\right]_0^1 = \frac{1}{2} - \frac{1}{3} = \frac{1}{6}. \]
\textbf{Moment about the $y$-axis:}
\[ \int_0^1 x(x - x^2)\,dx = \int_0^1 (x^2 - x^3)\,dx = \left[\frac{x^3}{3} - \frac{x^4}{4}\right]_0^1 = \frac{1}{3} - \frac{1}{4} = \frac{1}{12}
\quad\Rightarrow\quad \bar{x} = \frac{1/12}{1/6} = \frac{1}{2}. \]
\textbf{Moment about the $x$-axis:}
\[ \frac12\int_0^1 (x^2 - x^4)\,dx = \frac12\left[\frac{x^3}{3} - \frac{x^5}{5}\right]_0^1 = \frac12\left(\frac{1}{3} - \frac{1}{5}\right) = \frac{1}{15}
\quad\Rightarrow\quad \bar{y} = \frac{1/15}{1/6} = \frac{2}{5}. \]
The centroid is $\left(\tfrac12, \tfrac25\right)$. Note that $\bar{x}=\tfrac12$ could have been deduced from symmetry about the line $x=\tfrac12$ (the region is symmetric under $x \mapsto 1-x$).
\end{exoresolu}

\courssub{Standard Results Derived by Integration}

\begin{exoresolu}[Centroid of a semicircular lamina]
Show that the centroid of a uniform semicircular lamina of radius $r$ lies on its axis of symmetry at distance $\dfrac{4r}{3\pi}$ from the diameter.
\tcblower
Place the diameter on the $x$-axis, centre at the origin: the lamina is the region under $y = \sqrt{r^2 - x^2}$, $-r \le x \le r$. By symmetry $\bar{x} = 0$. The area is $A = \tfrac12\pi r^2$.
\[ M_x = \frac12\int_{-r}^{r}(r^2 - x^2)\,dx = \frac12\left[r^2 x - \frac{x^3}{3}\right]_{-r}^{r} = \frac12\cdot\frac{4r^3}{3} = \frac{2r^3}{3}. \]
Hence $\bar{y} = \dfrac{2r^3/3}{\pi r^2/2} = \dfrac{4r}{3\pi} \approx 0.424\,r$.
\end{exoresolu}

\begin{popfigure}
\begin{center}
\begin{tikzpicture}[scale=1.3]
\fill[popblueL] (2,0) arc (0:180:2) -- cycle;
\draw[thick, popblue] (2,0) arc (0:180:2) -- cycle;
\draw[->] (-2.5,0) -- (2.6,0) node[right] {$x$};
\draw[->] (0,-0.3) -- (0,2.5) node[above] {$y$};
\draw[dashed, popmuted] (0,0) -- (0,2.3);
\fill[popink] (0,0.849) circle (0.055);
\node[popink, right] at (0.1,0.849) {\scriptsize $G\left(0,\ \tfrac{4r}{3\pi}\right)$};
\draw[<->, popmuted] (-0.25,0) -- (-0.25,0.849) node[midway, left] {\scriptsize $\tfrac{4r}{3\pi}$};
\node[popdark] at (1.5,0.35) {\scriptsize $r$};
\draw[popdark] (0,0) -- (1.414,1.414);
\end{tikzpicture}
\end{center}
\legende{Semicircular lamina: the centroid lies on the axis of symmetry at height $\tfrac{4r}{3\pi}$ above the diameter.}
\end{popfigure}

\begin{exoresolu}[Centre of mass of a solid cone]
A uniform solid right circular cone has height $h$ and base radius $R$. Show that its centre of mass is on its axis, at distance $\dfrac{h}{4}$ from the base.
\tcblower
Generate the cone by rotating the line $y = \dfrac{R}{h}x$ about the $x$-axis for $0 \le x \le h$ (apex at the origin, base at $x = h$). By symmetry the centre of mass is on the $x$-axis.
\[ V = \pi\int_0^h \frac{R^2}{h^2}x^2\,dx = \frac{\pi R^2 h}{3}, \qquad
\pi\int_0^h x\,y^2\,dx = \pi\frac{R^2}{h^2}\int_0^h x^3\,dx = \pi\frac{R^2}{h^2}\cdot\frac{h^4}{4} = \frac{\pi R^2 h^2}{4}. \]
Therefore $\bar{x} = \dfrac{\pi R^2 h^2/4}{\pi R^2 h/3} = \dfrac{3h}{4}$ from the apex, i.e. $h - \tfrac{3h}{4} = \dfrac{h}{4}$ \textbf{from the base}.
\end{exoresolu}

\begin{popfigure}
\begin{center}
\begin{tikzpicture}[scale=1.0]
\fill[popblueL] (0,0) -- (4,1.4) arc (90:270:0.35 and 1.4) -- cycle;
\draw[thick, popblue] (0,0) -- (4,1.4);
\draw[thick, popblue] (0,0) -- (4,-1.4);
\draw[thick, popblue] (4,0) ellipse (0.35 and 1.4);
\draw[dashed, popmuted] (0,0) -- (4.8,0);
\fill[popink] (3,0) circle (0.06);
\node[popink, above] at (3,0.08) {\scriptsize $G$};
\draw[<->, popmuted] (3,-1.7) -- (4,-1.7) node[midway, below] {\scriptsize $\tfrac{h}{4}$};
\draw[<->, popmuted] (0,-1.7) -- (4,-1.7) node[midway, below] {\scriptsize $h$};
\node[popdark] at (0,-0.35) {\scriptsize apex};
\node[popdark, right] at (4.15,1.1) {\scriptsize base};
\end{tikzpicture}
\end{center}
\legende{Solid cone: the centre of mass is on the axis, at $\tfrac{h}{4}$ from the base ($\tfrac{3h}{4}$ from the apex).}
\end{popfigure}

\begin{exoresolu}[Centre of mass of a solid hemisphere]
Show that the centre of mass of a uniform solid hemisphere of radius $r$ is at distance $\tfrac{3r}{8}$ from its plane face.
\tcblower
Generate the hemisphere by rotating the quarter-circle $y = \sqrt{r^2 - x^2}$, $0 \le x \le r$, about the $x$-axis. The plane face is at $x=0$. By symmetry the centre of mass is on the $x$-axis.
\[ V = \pi\int_0^r (r^2 - x^2)\,dx = \pi\left[r^2 x - \frac{x^3}{3}\right]_0^r = \frac{2}{3}\pi r^3. \]
Moment about the plane face ($x=0$):
\[ \pi\int_0^r x(r^2 - x^2)\,dx = \pi\left[\frac{r^2 x^2}{2} - \frac{x^4}{4}\right]_0^r = \pi\left(\frac{r^4}{2} - \frac{r^4}{4}\right) = \frac{\pi r^4}{4}. \]
Hence $\bar{x} = \dfrac{\pi r^4/4}{2\pi r^3/3} = \dfrac{3r}{8}$ from the plane face.
\end{exoresolu}

\begin{aretenir}[Standard centroid and centre of mass results]
\begin{itemize}
\item \textbf{Triangular lamina:} centroid at intersection of medians, one third of the height above the base: $\bar{y} = \tfrac{h}{3}$.
\item \textbf{Semicircular lamina} (radius $r$): $\bar{y} = \dfrac{4r}{3\pi}$ above the diameter.
\item \textbf{Solid hemisphere} (radius $r$): $\bar{x} = \dfrac{3r}{8}$ from the flat face.
\item \textbf{Solid cone} (height $h$): $\bar{x} = \dfrac{h}{4}$ from the base (or $\tfrac{3h}{4}$ from the apex).
\end{itemize}
\end{aretenir}

\courssub{Applications and Examination Strategy}

The centre of mass governs \emph{stability}: an object topples when the vertical line through its centre of mass falls outside its base. This is why a full water tanker is more stable when the tank's centre of mass is low, why the Leaning Tower of Pisa still stands (its centre of mass still projects inside the base), and why arches and drinking vessels are designed with low balancing points. Locating $G$ is also the first step in computing moments of inertia in mechanics.

\begin{methode}[Exam technique for centroid questions]
\begin{enumerate}
\item \textbf{Sketch} the region or solid; indicate the axis of rotation if applicable.
\item \textbf{Identify} symmetry: state immediately any coordinate fixed by symmetry.
\item \textbf{Quote} the formula you will use ($\bar{x}$, $\bar{y}$ or the solid-of-revolution formula).
\item \textbf{Compute separately} the area $A$ (or volume $V$) and each moment integral, keeping exact values (fractions, $\pi$, surds).
\item \textbf{Divide} at the very end, and simplify.
\item \textbf{Check} plausibility: the centroid must lie inside (or on the boundary side of) the region, and any symmetry coordinate must match.
\end{enumerate}
\end{methode}

\begin{attention}[Common mistakes]
\begin{itemize}
\item Forgetting the factor $\tfrac12$ in $\bar{y}$ for a lamina.
\item Using $y$ instead of $y^2$ in the solid-of-revolution formula $\bar{x} = \int xy^2\,dx \big/ \int y^2\,dx$.
\item Mixing up ``from the apex'' and ``from the base'' for a cone: state clearly which reference point you use.
\item Ignoring symmetry and computing an integral that is obviously zero or obviously fixed.
\item Forgetting that the formulas for $\bar{x}$ and $\bar{y}$ of a lamina assume the density is constant and cancels; if density varies, the mass element is $\rho(x,y)\,dA$ and the formulas change.
\end{itemize}
\end{attention}

\begin{exercice}[Practice]
\begin{enumerate}
\item Find the centroid of the region bounded by $y = \sqrt{x}$, the $x$-axis and the line $x = 4$.
\item The region under $y = \sin x$, $0 \le x \le \pi$, is rotated about the $x$-axis. Find the centre of mass of the resulting solid.
\item Show that the centre of mass of a uniform solid hemisphere of radius $r$ is at distance $\tfrac{3r}{8}$ from its plane face.
\item Find the centroid of the region bounded by $y = x^2$ and $y = 2x$.
\item A uniform solid is generated by rotating the region bounded by $y = x^2$, $x = 1$ and the $x$-axis about the $y$-axis. Find the centre of mass of this solid.
\end{enumerate}
\end{exercice}

\begin{corrige}[Exercise 1]
$A = \int_0^4 \sqrt{x}\,dx = \left[\tfrac23 x^{3/2}\right]_0^4 = \tfrac{16}{3}$.
$\int_0^4 x\sqrt{x}\,dx = \left[\tfrac25 x^{5/2}\right]_0^4 = \tfrac{64}{5}$, so $\bar{x} = \tfrac{64/5}{16/3} = \tfrac{12}{5}$.
$\tfrac12\int_0^4 x\,dx = \tfrac12\cdot 8 = 4$, so $\bar{y} = \tfrac{4}{16/3} = \tfrac34$. Centroid: $\left(\tfrac{12}{5}, \tfrac34\right)$.
\end{corrige}

\begin{corrige}[Exercise 2]
By symmetry about $x = \tfrac{\pi}{2}$, $\bar{x} = \tfrac{\pi}{2}$. Verification: $\int_0^\pi \sin^2 x\,dx = \tfrac{\pi}{2}$ and $\int_0^\pi x\sin^2 x\,dx = \tfrac{\pi^2}{4}$ (integrate by parts after writing $\sin^2 x = \tfrac{1-\cos 2x}{2}$), so $\bar{x} = \tfrac{\pi^2/4}{\pi/2} = \tfrac{\pi}{2}$, and $\bar{y} = 0$.
\end{corrige}

\begin{corrige}[Exercise 3]
Rotate $y = \sqrt{r^2 - x^2}$, $0 \le x \le r$, about the $x$-axis. $V = \pi\int_0^r (r^2 - x^2)\,dx = \tfrac23\pi r^3$. Moment: $\pi\int_0^r x(r^2 - x^2)\,dx = \pi\left[\tfrac{r^2x^2}{2} - \tfrac{x^4}{4}\right]_0^r = \tfrac{\pi r^4}{4}$. Hence $\bar{x} = \dfrac{\pi r^4/4}{2\pi r^3/3} = \dfrac{3r}{8}$.
\end{corrige}

\begin{corrige}[Exercise 4]
Curves intersect at $x=0$ and $x=2$. Top: $f(x)=2x$, bottom: $g(x)=x^2$.
$A = \int_0^2 (2x - x^2)\,dx = \left[x^2 - \frac{x^3}{3}\right]_0^2 = 4 - \frac{8}{3} = \frac{4}{3}$.
$\int_0^2 x(2x - x^2)\,dx = \int_0^2 (2x^2 - x^3)\,dx = \left[\frac{2x^3}{3} - \frac{x^4}{4}\right]_0^2 = \frac{16}{3} - 4 = \frac{4}{3}$, so $\bar{x} = 1$.
$\frac12\int_0^2 (4x^2 - x^4)\,dx = \frac12\left[\frac{4x^3}{3} - \frac{x^5}{5}\right]_0^2 = \frac12\left(\frac{32}{3} - \frac{32}{5}\right) = \frac{32}{15}$, so $\bar{y} = \frac{32/15}{4/3} = \frac{8}{5}$. Centroid: $\left(1, \frac{8}{5}\right)$.
\end{corrige}

\begin{corrige}[Exercise 5]
Region: $y=x^2$, $x=1$, $x$-axis. Rotate about $y$-axis. Use cylindrical shells or express $x=\sqrt{y}$.
Volume by shells: $V = 2\pi\int_0^1 x\cdot x^2\,dx = 2\pi\left[\frac{x^4}{4}\right]_0^1 = \frac{\pi}{2}$.
Moment about $y=0$ (i.e. $\bar{y}$): each shell at $x$ has height $x^2$, centroid at $y=\frac{x^2}{2}$, mass $2\pi x\cdot x^2\,\delta x$. Moment = $\int 2\pi x\cdot x^2\cdot \frac{x^2}{2}\,dx = \pi\int_0^1 x^5\,dx = \frac{\pi}{6}$.
Thus $\bar{y} = \frac{\pi/6}{\pi/2} = \frac{1}{3}$. By symmetry $\bar{x}=0$. Centre of mass: $\left(0, \frac{1}{3}\right)$.
\end{corrige}

\begin{savaistu}[Pappus and his theorems]
The Greek mathematician Pappus of Alexandria (4th century) proved two beautiful theorems:
\begin{enumerate}
\item The area of a surface of revolution equals the length of the generating curve times the distance travelled by its centroid.
\item The volume of a solid of revolution equals the area of the generating region times the distance travelled by its centroid: $V = 2\pi\bar{y}\,A$ (for rotation about the $x$-axis).
\end{enumerate}
Check it on the cone: rotating a right triangle of area $\tfrac12 Rh$ whose centroid is at height $\tfrac{R}{3}$ gives $V = 2\pi\cdot\tfrac{R}{3}\cdot\tfrac{Rh}{2} = \tfrac13\pi R^2 h$ — exactly the formula you know. Centroids and volumes are two faces of the same coin.
\end{savaistu}

\coursec{Numerical Integration: the Trapezium Rule}

In the previous section we used integration to locate centroids and centres of mass, and in every case we were able to find a primitive of the integrand and apply the Fundamental Theorem of Calculus. But life -- and the GCE examination -- is not always so kind. In this section we answer a very practical question: \emph{what can we do when no primitive is available, or when the function is only known through a table of measured values?} The answer is \cle{numerical integration}, and the first tool we study is the \cle{trapezium rule} (also called the trapezoidal rule).

\courssub{Why Numerical Integration?}

The Fundamental Theorem tells us that $\int_a^b f(x)\,dx = F(b) - F(a)$ \emph{provided} we can find a primitive $F$ of $f$. Two common situations defeat this strategy.

\textbf{1. Functions without elementary primitives.} Some perfectly well-behaved functions have primitives that cannot be expressed in terms of the functions we know (polynomials, exponentials, logarithms, trigonometric functions). Famous examples include
\[
\int e^{x^2}\,dx, \qquad \int \frac{\sin x}{x}\,dx, \qquad \int \sqrt{1+x^3}\,dx .
\]
These integrals are not ``impossible'' -- the area under the curve certainly exists -- but there is no closed formula for the primitive. The integral $\int_0^1 e^{x^2}\,dx$, for instance, appears constantly in probability (the normal distribution), yet it can only be evaluated \emph{approximately}.

\textbf{2. Data given only as a table.} In experimental work we often do not know a formula for $f$ at all. A meteorological station in Bamenda records rainfall intensity every hour; an engineer measures the cross-sectional width of the Mungo river every $5$ metres; a speedometer reading is taken every $10$ seconds. The quantity we want (total rainfall, cross-sectional area, distance travelled) is an integral, but the integrand exists only as a list of numbers.

In both cases the strategy is the same: replace the exact area by a sum of simple areas that we \emph{can} compute from the available values.

\begin{savaistu}[Historical note]
The trapezium rule was known to ancient Babylonian astronomers (c. 350 BCE) who used it to compute the distance travelled by Jupiter from velocity measurements. In Cameroon, surveyors still use it today to estimate land areas from irregular boundary measurements.
\end{savaistu}

\courssub{From Curves to Chords: One Trapezium}

Recall how the definite integral was born (Section 1): we sliced $[a,b]$ into thin strips and approximated each strip by a \emph{rectangle} (Riemann sums). The trapezium rule improves on this by approximating each strip by a \emph{trapezium}: instead of replacing the curve by a horizontal segment, we replace it by the \cle{chord} joining the two endpoints of the arc.

Consider one strip of width $h$, bounded by the ordinates $y_0 = f(x_0)$ and $y_1 = f(x_1)$, with $x_1 - x_0 = h$. The chord from $(x_0, y_0)$ to $(x_1, y_1)$ forms, together with the $x$-axis and the two ordinates, a trapezium whose parallel sides have lengths $y_0$ and $y_1$ and whose height (distance between them) is $h$. The area of a trapezium is
\[
A = \frac{1}{2}(\text{sum of parallel sides}) \times (\text{distance between them}) = \frac{h}{2}(y_0 + y_1).
\]

\begin{propriete}[Area of one trapezium]
If the arc of $y = f(x)$ over $[x_0, x_1]$ is replaced by its chord, then
\[
\int_{x_0}^{x_1} f(x)\,dx \;\approx\; \frac{h}{2}\bigl(y_0 + y_1\bigr), \qquad h = x_1 - x_0 .
\]
\end{propriete}

\courssub{The Trapezium Rule with $n$ Strips}

Now divide $[a, b]$ into $n$ strips of \emph{equal} width
\[
h = \frac{b-a}{n},
\]
with subdivision points $x_0 = a,\ x_1 = a+h,\ \dots,\ x_n = b$, and ordinates $y_r = f(x_r)$ for $r = 0, 1, \dots, n$. Applying the one-strip formula to each strip and adding,
\[
\int_a^b f(x)\,dx \approx \frac{h}{2}(y_0+y_1) + \frac{h}{2}(y_1+y_2) + \dots + \frac{h}{2}(y_{n-1}+y_n).
\]
Every \emph{interior} ordinate $y_1, \dots, y_{n-1}$ appears in exactly two neighbouring trapezia, hence is counted twice, while the two \emph{end} ordinates $y_0$ and $y_n$ appear only once. Factoring $\frac{h}{2}$ gives the rule.

\begin{propriete}[The trapezium rule]
For $n$ strips of equal width $h = \dfrac{b-a}{n}$,
\[
\boxed{\;\int_a^b f(x)\,dx \;\approx\; \frac{h}{2}\Bigl[\,y_0 + 2\bigl(y_1 + y_2 + \dots + y_{n-1}\bigr) + y_n\,\Bigr]\;}
\]
where $y_r = f(a + rh)$. In words: \emph{half the width, times (first + last + twice the sum of the middle ordinates)}. (The derivation above, by summing the areas of $n$ trapezia, is examinable.)
\end{propriete}

\begin{popfigure}
\begin{center}
\begin{tikzpicture}
\begin{axis}[
    width=12cm, height=7cm,
    axis lines=middle,
    xlabel={$x$}, ylabel={$y$},
    xmin=-0.3, xmax=4.6, ymin=-0.4, ymax=3.4,
    xtick={0,1,2,3,4},
    xticklabels={$x_0$,$x_1$,$x_2$,$x_3$,$x_4$},
    ytick=\empty,
    clip=false]
\addplot[domain=0:4, samples=100, thick, popblue] {0.25*(x-1.2)^2 + 0.9};
\addplot[domain=0:1, draw=poporange, thick, fill=poporange, fill opacity=0.15] {1.04 - 0.15*x} \closedcycle;
\addplot[domain=1:2, draw=poporange, thick, fill=poporange, fill opacity=0.15] {0.905 + 0.155*(x-1)} \closedcycle;
\addplot[domain=2:3, draw=poporange, thick, fill=poporange, fill opacity=0.15] {1.06 + 0.65*(x-2)} \closedcycle;
\addplot[domain=3:4, draw=poporange, thick, fill=poporange, fill opacity=0.15] {1.71 + 1.15*(x-3)} \closedcycle;
\addplot[domain=0:4, samples=100, thick, popblue] {0.25*(x-1.2)^2 + 0.9};
\draw[dashed, popmuted] (axis cs:0,0) -- (axis cs:0,1.04) node[pos=0.6, left, popdark] {$y_0$};
\draw[dashed, popmuted] (axis cs:1,0) -- (axis cs:1,0.905) node[pos=0.65, left, popdark] {$y_1$};
\draw[dashed, popmuted] (axis cs:2,0) -- (axis cs:2,1.06) node[pos=0.65, left, popdark] {$y_2$};
\draw[dashed, popmuted] (axis cs:3,0) -- (axis cs:3,1.71) node[pos=0.65, left, popdark] {$y_3$};
\draw[dashed, popmuted] (axis cs:4,0) -- (axis cs:4,2.86) node[pos=0.65, right, popdark] {$y_4$};
\addplot[only marks, mark=*, mark size=1.5pt, popdark] coordinates {(0,1.04) (1,0.905) (2,1.06) (3,1.71) (4,2.86)};
\end{axis}
\end{tikzpicture}
\legende{The trapezium rule with $n = 4$ strips: the curve (blue) is replaced by chords (orange), and the area is approximated by the sum of the four trapezia.}
\end{center}
\end{popfigure}

\begin{methode}[Applying the trapezium rule]
\begin{enumerate}
\item \textbf{Compute the strip width} $h = \dfrac{b-a}{n}$, where $n$ is the number of \emph{strips} (intervals).
\item \textbf{Build a table of ordinates}: list $x_0, x_1, \dots, x_n$ (there are $n+1$ of them) and compute each $y_r = f(x_r)$, keeping at least 4 decimal places.
\item \textbf{Apply the formula}:
\[
\int_a^b f(x)\,dx \approx \frac{h}{2}\Bigl[\underbrace{y_0 + y_n}_{\text{ends}} + 2\underbrace{(y_1 + \dots + y_{n-1})}_{\text{middle ones}}\Bigr].
\]
\item \textbf{Round only at the end}, to the accuracy requested (usually 3 or 4 significant figures).
\end{enumerate}
\end{methode}

\begin{attention}[Strips versus ordinates]
``Use the trapezium rule with \textbf{4 intervals}'' and ``with \textbf{5 ordinates}'' mean exactly the same thing: $n = 4$ strips require $n + 1 = 5$ ordinates $y_0, y_1, y_2, y_3, y_4$. Confusing the two is one of the most frequent errors at the GCE: candidates who read ``4 intervals'' and compute only 4 ordinates lose a whole strip and get a wrong $h$. Always check: \emph{number of ordinates} $=$ \emph{number of intervals} $+ 1$.
\end{attention}

\begin{exoresolu}[Estimating $\displaystyle\int_0^1 e^{x^2}\,dx$]
Use the trapezium rule with $4$ intervals to estimate $\displaystyle I = \int_0^1 e^{x^2}\,dx$, giving your answer to $4$ decimal places.
\tcblower
\textbf{Solution.} Here $a = 0$, $b = 1$, $n = 4$, so $h = \dfrac{1-0}{4} = 0.25$. The $5$ ordinates are computed in a table:
\begin{center}
\begin{tabular}{|c|c|c|}
\hline
\rowcolor{popblueL} $r$ & $x_r$ & $y_r = e^{x_r^2}$ \\ \hline
$0$ & $0.00$ & $e^{0} = 1.000000$ \\ \hline
$1$ & $0.25$ & $e^{0.0625} = 1.064494$ \\ \hline
$2$ & $0.50$ & $e^{0.25} = 1.284025$ \\ \hline
$3$ & $0.75$ & $e^{0.5625} = 1.755055$ \\ \hline
$4$ & $1.00$ & $e^{1} = 2.718282$ \\ \hline
\end{tabular}
\end{center}
Ends: $y_0 + y_4 = 1.000000 + 2.718282 = 3.718282$.\\
Middle: $y_1 + y_2 + y_3 = 1.064494 + 1.284025 + 1.755055 = 4.103574$.\\[2pt]
Therefore
\[
I \approx \frac{0.25}{2}\bigl[3.718282 + 2(4.103574)\bigr] = 0.125 \times 11.925430 = 1.490679.
\]
Hence $\displaystyle\int_0^1 e^{x^2}\,dx \approx 1.4907$ (4 d.p.).
\end{exoresolu}

\begin{attention}[Do not round the ordinates too early]
If in the example above you round each ordinate to 2 decimal places ($1.00,\ 1.06,\ 1.28,\ 1.76,\ 2.72$), you obtain $I \approx 1.4900$ instead of $1.4907$: the rounding errors accumulate inside the bracket and are then multiplied by $\frac{h}{2}$. \textbf{Rule of thumb:} keep at least 4 decimal places (or 2 more than requested) in all intermediate values, and round only the final answer.
\end{attention}

\courssub{Working with Tabulated Data}

When the function is unknown and only measurements are given, the same formula applies directly to the data -- this is precisely why the trapezium rule is so useful in applications.

\begin{exemplebox}[Cross-sectional area of a river]
To estimate the flow of a river near Maroua, an engineer measures the depth $y$ (in metres) of the channel at points $2$ metres apart across its $10$-metre width:
\begin{center}
\begin{tabular}{|l|c|c|c|c|c|c|}
\hline
\rowcolor{popblueL} Distance $x$ (m) & $0$ & $2$ & $4$ & $6$ & $8$ & $10$ \\ \hline
Depth $y$ (m) & $0.0$ & $1.2$ & $1.8$ & $1.5$ & $0.9$ & $0.0$ \\ \hline
\end{tabular}
\end{center}
The cross-sectional area is $A = \int_0^{10} y\,dx$. With $h = 2$ and $n = 5$ strips:
\[
A \approx \frac{2}{2}\bigl[(0.0 + 0.0) + 2(1.2 + 1.8 + 1.5 + 0.9)\bigr] = 1 \times 2(5.4) = 10.8\ \mathrm{m^2}.
\]
No formula for $y(x)$ was needed -- the table alone suffices.
\end{exemplebox}

\courssub{Accuracy: Increasing $n$, Over- and Under-estimation}

\textbf{Effect of $n$.} Each trapezium replaces a small arc by a chord. The thinner the strips, the closer each chord lies to its arc, so the approximation improves as $n$ increases. It can be shown that the error of the trapezium rule behaves roughly like $\dfrac{(b-a)^3}{12n^2}\,\max|f''|$: doubling the number of strips divides the error by about $4$. For $\int_0^1 e^{x^2}\,dx$, whose exact value is $1.46265\ldots$, the trapezium rule gives $1.4907$ with $n = 4$ and $1.4697$ with $n = 8$ -- clearly converging to the true value.

\textbf{Over- or under-estimation.} Look at the concavity of the graph:
\begin{itemize}
\item If $f$ is \textbf{concave down} ($f'' < 0$, shaped like $\frown$), each chord lies \emph{below} the arc, so the trapezium rule \emph{under-estimates} the integral.
\item If $f$ is \textbf{concave up} ($f'' > 0$, shaped like $\smile$), each chord lies \emph{above} the arc, so the rule \emph{over-estimates} the integral.
\end{itemize}
Since $\frac{d^2}{dx^2}\bigl(e^{x^2}\bigr) = (2 + 4x^2)e^{x^2} > 0$, the curve in our worked example is concave up, and indeed $1.4907 > 1.4627$: an over-estimate, as predicted.

\begin{popfigure}
\begin{center}
\begin{tikzpicture}[scale=2.2]
\draw[->] (-0.2,0) -- (3.4,0) node[right] {$x$};
\draw[->] (0,-0.2) -- (0,2.2) node[above] {$y$};
\draw[very thick, popblue, domain=0.4:2.8, samples=80] plot (\x, {0.6 + 1.1*sin(deg(1.05*\x))});
\draw[thick, poporange] (0.4, {0.6 + 1.1*sin(deg(1.05*0.4))}) -- (2.8, {0.6 + 1.1*sin(deg(1.05*2.8))});
\draw[dashed, popmuted] (0.4,0) -- (0.4, {0.6 + 1.1*sin(deg(1.05*0.4))});
\draw[dashed, popmuted] (2.8,0) -- (2.8, {0.6 + 1.1*sin(deg(1.05*2.8))});
\node[popdark] at (1.6, 1.95) {curve (concave down)};
\node[poporange] at (1.6, 1.12) {chord};
\node[align=center, popdark] at (1.6, 0.45) {trapezium area\\ $<$ true area};
\end{tikzpicture}
\legende{For a concave-down arc, the chord lies below the curve: the trapezium under-estimates the area. For a concave-up arc the opposite holds.}
\end{center}
\end{popfigure}

\begin{attention}[Exam tip: comment on your estimate]
GCE questions frequently add: ``\emph{State, with a reason, whether your estimate is an over-estimate or an under-estimate.}'' Do not compute the exact integral to answer! Simply look at the graph (or the sign of $f''$): concave up $\Rightarrow$ over-estimate; concave down $\Rightarrow$ under-estimate. Sketching the curve and one chord in the margin is an excellent way to see it.
\end{attention}

\courssub{Comparison with the Exact Value and Error Estimation}

When the integral \emph{can} be computed exactly, we can measure the quality of the approximation.

\begin{exoresolu}[Trapezium rule versus exact integration]
(a) Use the trapezium rule with $4$ intervals to estimate $\displaystyle I = \int_0^2 x^2\,dx$.\\
(b) Find the exact value of $I$ and hence the percentage error of the estimate.\\
(c) Explain the sign of the error.
\tcblower
\textbf{Solution.} (a) $h = \dfrac{2-0}{4} = 0.5$. Ordinates: $y_0 = 0$, $y_1 = 0.25$, $y_2 = 1$, $y_3 = 2.25$, $y_4 = 4$.
\[
I \approx \frac{0.5}{2}\bigl[(0 + 4) + 2(0.25 + 1 + 2.25)\bigr] = 0.25 \times \bigl[4 + 7\bigr] = 2.75.
\]
(b) Exact value: $I = \left[\dfrac{x^3}{3}\right]_0^2 = \dfrac{8}{3} = 2.6667$ (4 d.p.). Hence
\[
\text{error} = 2.75 - 2.6667 = 0.0833, \qquad \text{percentage error} = \frac{0.0833}{2.6667}\times 100\% \approx 3.1\%.
\]
(c) $f(x) = x^2$ has $f''(x) = 2 > 0$: the curve is concave up, every chord lies above the arc, so the trapezium rule must \emph{over-estimate} -- consistent with $2.75 > 2.6667$.
\end{exoresolu}

\courssub{Extra: a Glimpse of the Midpoint Rule}

The trapezium rule is not the only possibility. In the \cle{midpoint rule}, each strip is approximated by the rectangle whose height is the value of $f$ at the \emph{midpoint} of the strip:
\[
\int_a^b f(x)\,dx \approx h\Bigl[f\Bigl(x_0 + \tfrac{h}{2}\Bigr) + f\Bigl(x_0 + \tfrac{3h}{2}\Bigr) + \dots + f\Bigl(x_0 + \tfrac{(2n-1)h}{2}\Bigr)\Bigr].
\]
For smooth curves its error is typically about half that of the trapezium rule and of opposite sign, so the two estimates bracket the true value -- a handy way to check a numerical answer. At GCE level, however, only the trapezium rule is examinable; the midpoint rule is mentioned here for culture and as a checking device.

\begin{aretenir}[Key points]
\begin{itemize}
\item Numerical integration is needed when no elementary primitive exists (e.g.\ $e^{x^2}$, $\frac{\sin x}{x}$) or when $f$ is known only through a table of values.
\item Trapezium rule with $n$ strips of width $h = \frac{b-a}{n}$:
\[
\int_a^b f(x)\,dx \approx \frac{h}{2}\bigl[y_0 + 2(y_1 + \dots + y_{n-1}) + y_n\bigr].
\]
\item $n$ intervals $\Leftrightarrow$ $n+1$ ordinates; keep at least 4 decimal places in intermediate values.
\item Increasing $n$ improves accuracy (error roughly proportional to $1/n^2$).
\item Concave up $\Rightarrow$ over-estimate; concave down $\Rightarrow$ under-estimate.
\end{itemize}
\end{aretenir}

\begin{exercice}[]
The speed $v$ (in m\,s$^{-1}$) of a car leaving Douala is recorded every $5$ seconds:
\begin{center}
\begin{tabular}{|l|c|c|c|c|c|}
\hline
\rowcolor{popblueL} $t$ (s) & $0$ & $5$ & $10$ & $15$ & $20$ \\ \hline
$v$ (m\,s$^{-1}$) & $0$ & $6.5$ & $11.0$ & $13.5$ & $15.0$ \\ \hline
\end{tabular}
\end{center}
Use the trapezium rule with $4$ intervals to estimate the distance travelled during these $20$ seconds.
\end{exercice}

\begin{corrige}[Exercise 1]
Distance $= \int_0^{20} v\,dt$, with $h = 5$:
\[
\frac{5}{2}\bigl[(0 + 15.0) + 2(6.5 + 11.0 + 13.5)\bigr] = 2.5 \times \bigl[15 + 62\bigr] = 2.5 \times 77 = 192.5\ \mathrm{m}.
\]
\end{corrige}

\begin{exercice}[]
Use the trapezium rule with $5$ ordinates to estimate $\displaystyle\int_1^2 \frac{\sin x}{x}\,dx$, working to $4$ decimal places. Given that $\frac{d^2}{dx^2}\left(\frac{\sin x}{x}\right) < 0$ on $[1,2]$, state whether your answer is an over- or under-estimate.
\end{exercice}

\begin{corrige}[Exercise 2]
$5$ ordinates means $n = 4$ strips, $h = 0.25$. With $f(x) = \frac{\sin x}{x}$ (radians):
\begin{center}
\begin{tabular}{|c|c|c|c|c|c|}
\hline
\rowcolor{popblueL} $x$ & $1.00$ & $1.25$ & $1.50$ & $1.75$ & $2.00$ \\ \hline
$f(x)$ & $0.8415$ & $0.7592$ & $0.6650$ & $0.5623$ & $0.4546$ \\ \hline
\end{tabular}
\end{center}
\[
\int_1^2 \frac{\sin x}{x}\,dx \approx \frac{0.25}{2}\bigl[(0.8415 + 0.4546) + 2(0.7592 + 0.6650 + 0.5623)\bigr] = 0.125 \times 5.2691 = 0.6586.
\]
Since $f'' < 0$, the curve is concave down, chords lie below the arcs: the estimate $0.6586$ is an \textbf{under-estimate}.
\end{corrige}

\coursec{Mean Value Theorem for Integrals and Synthesis}

In the previous section we learned how to \emph{estimate} a definite integral numerically with the trapezium rule when an exact primitive is out of reach. We now close the chapter by asking a different, very natural question: if $\displaystyle\int_a^b f(x)\,dx$ measures an ``accumulated quantity'', what is the \emph{average} of the function over $[a,b]$? This single idea — the \cle{mean value} of a function — leads to one of the most elegant theorems of the integral calculus, gives us a powerful way to bound integrals, and appears regularly in GCE Advanced Level synthesis questions. We finish with a revision map of the whole chapter and an examination strategy.

\courssub{The Mean (Average) Value of a Function}

\textbf{From a discrete average to a continuous average.}\par
If a student scores marks $x_1, x_2, \dots, x_n$ in $n$ tests, the average mark is $\bar{x} = \dfrac{1}{n}\displaystyle\sum_{i=1}^{n} x_i$. Now suppose a quantity varies \emph{continuously}: the temperature in Yaoundé at every instant of a day, the speed of a bus at every second of its journey from Douala to Bafoussam. We cannot average infinitely many values by adding them, but we can sample: divide $[a,b]$ into $n$ equal strips of width $\Delta x = \dfrac{b-a}{n}$ and average the sample values $f(x_i)$:
\[
\text{average} \approx \frac{1}{n}\sum_{i=1}^{n} f(x_i) = \frac{1}{b-a}\sum_{i=1}^{n} f(x_i)\,\Delta x ,
\]
where we have used $\dfrac{1}{n} = \dfrac{\Delta x}{b-a}$. As $n \to \infty$ the Riemann sum tends to the definite integral, and we recognise the limit as $\dfrac{1}{b-a}\displaystyle\int_a^b f(x)\,dx$. This motivates the definition.

\begin{definition}[Mean value of a function on an interval]
Let $f$ be continuous on $[a,b]$ with $a<b$. The \cle{mean value} (or \cle{average value}) of $f$ on $[a,b]$ is
\[
\bar{f} = \frac{1}{b-a}\int_a^b f(x)\,dx .
\]
\end{definition}

\textbf{Geometric interpretation.}\par
Multiplying the definition by $(b-a)$ gives
\[
\bar{f}\,(b-a) = \int_a^b f(x)\,dx .
\]
The right-hand side is the (signed) area under the curve; the left-hand side is the area of a \emph{rectangle} of base $(b-a)$ and height $\bar{f}$. So the mean value is exactly the \cle{height of the rectangle} on the same base whose area equals the area under the curve.

\begin{popfigure}
\begin{center}
\begin{tikzpicture}
\begin{axis}[
  width=12cm, height=7cm,
  axis lines=middle,
  xlabel={$x$}, ylabel={$y$},
  xmin=0, xmax=6.4, ymin=0, ymax=4.2,
  xtick={1,3.05,5}, xticklabels={$a$,$c$,$b$},
  ytick=\empty,
  clip=false
]
\addplot[draw=poporange, fill=poporange, fill opacity=0.15] coordinates {(1,0) (5,0) (5,2.85) (1,2.85)} -- cycle;
\addplot[domain=1:5, samples=100, draw=none, fill=popblue, fill opacity=0.12] {2.2 + 0.9*sin(deg(x-1)) + 0.15*(x-1)} \closedcycle;
\addplot[domain=1:5, samples=100, thick, popblue] {2.2 + 0.9*sin(deg(x-1)) + 0.15*(x-1)};
\addplot[domain=1:5, samples=2, dashed, thick, poporange] {2.85};
\addplot[only marks, mark=*, popdark] coordinates {(3.05,2.85)};
\draw[dashed, popmuted] (axis cs:3.05,0) -- (axis cs:3.05,2.85);
\node[poporange, anchor=south west] at (axis cs:4.1,2.9) {$y=\bar{f}$};
\node[popblue] at (axis cs:4.6,3.9) {$y=f(x)$};
\node[popdark, anchor=west] at (axis cs:3.1,2.55) {$(c,\bar{f})$};
\end{axis}
\end{tikzpicture}
\legende{The rectangle of height $\bar{f}$ and base $(b-a)$ has the same area as the region under the curve. The curve crosses the line $y=\bar{f}$ at $x=c$.}
\end{center}
\end{popfigure}

\begin{attention}[The mean value is not $\dfrac{f(a)+f(b)}{2}$]
A very common error is to average only the two endpoints: $\bar{f} \neq \dfrac{f(a)+f(b)}{2}$ in general. The formula $\dfrac{f(a)+f(b)}{2}$ is exact \emph{only when $f$ is linear} on $[a,b]$ (it is then the trapezium rule with one strip, which is exact for straight lines). For example, for $f(x)=x^2$ on $[0,2]$: $\dfrac{f(0)+f(2)}{2} = 2$, but the true mean value is $\dfrac{1}{2}\displaystyle\int_0^2 x^2\,dx = \dfrac{1}{2}\cdot\dfrac{8}{3} = \dfrac{4}{3}$.
\end{attention}

\courssub{The Mean Value Theorem for Integrals}

The figure above suggests something remarkable: the horizontal line $y=\bar{f}$ must actually \emph{meet} the curve somewhere inside the interval. This is always true for continuous functions.

\begin{propriete}[Mean value theorem for integrals]
If $f$ is continuous on $[a,b]$, then there exists at least one number $c \in [a,b]$ such that
\[
f(c) = \frac{1}{b-a}\int_a^b f(x)\,dx, \qquad\text{equivalently}\qquad f(c)\,(b-a) = \int_a^b f(x)\,dx .
\]
\emph{Proof (instructive, not required for the examination).} Since $f$ is continuous on the closed interval $[a,b]$, it attains a minimum $m$ and a maximum $M$ there. Integrating the inequality $m \le f(x) \le M$ over $[a,b]$ gives
\[
m(b-a) \le \int_a^b f(x)\,dx \le M(b-a), \qquad\text{so}\qquad m \le \frac{1}{b-a}\int_a^b f(x)\,dx \le M .
\]
Thus the mean value $\bar{f}$ lies between the least and greatest values of $f$. By the \emph{intermediate value theorem}, the continuous function $f$ takes every value between $m$ and $M$; in particular there exists $c\in[a,b]$ with $f(c)=\bar{f}$. $\square$
\end{propriete}

\begin{attention}[Existence, not uniqueness]
The theorem guarantees \emph{at least one} value $c$; there may be several. It also says nothing about \emph{where} $c$ is — that must be found by solving an equation. Note that continuity is essential: a function with a jump on $[a,b]$ may never take its mean value.
\end{attention}

\begin{methode}[Finding the value(s) of $c$ explicitly]
\begin{enumerate}
\item Compute the integral $\displaystyle I = \int_a^b f(x)\,dx$ exactly.
\item Divide by $(b-a)$ to obtain the mean value $\bar{f} = \dfrac{I}{b-a}$.
\item Solve the equation $f(c) = \bar{f}$.
\item Keep only the solutions lying in $[a,b]$.
\end{enumerate}
\end{methode}

\begin{exoresolu}[Finding $c$ for a polynomial]
Find the mean value of $f(x) = x^2 + 1$ on $[0,3]$, and determine the value(s) of $c \in [0,3]$ guaranteed by the mean value theorem for integrals.
\tcblower
\textbf{Solution.}
\textbf{Step 1:} $\displaystyle \int_0^3 (x^2+1)\,dx = \left[\frac{x^3}{3}+x\right]_0^3 = (9+3) - 0 = 12$.

\textbf{Step 2:} $\bar{f} = \dfrac{12}{3-0} = 4$.

\textbf{Step 3:} Solve $f(c)=4$: $c^2+1=4 \Rightarrow c^2=3 \Rightarrow c=\pm\sqrt{3}$.

\textbf{Step 4:} Only $c=\sqrt{3}\approx 1.73$ lies in $[0,3]$. Hence $c=\sqrt{3}$, and indeed $f(\sqrt{3})=3+1=4=\bar{f}$.
\end{exoresolu}

\begin{exoresolu}[Finding $c$ for a trigonometric function]
Find the mean value of $f(x)=\sin x$ on $\left[0,\pi\right]$ and the corresponding value(s) of $c$.
\tcblower
\textbf{Solution.}
$\displaystyle\int_0^{\pi} \sin x\,dx = \big[-\cos x\big]_0^{\pi} = (-\cos\pi) - (-\cos 0) = 1+1 = 2$, so $\bar{f} = \dfrac{2}{\pi} \approx 0.637$.

Solve $\sin c = \dfrac{2}{\pi}$ on $[0,\pi]$: since $\dfrac{2}{\pi}<1$ and $\sin$ rises from $0$ to $1$ then falls back to $0$ on this interval, there are \emph{two} solutions,
\[
c_1 = \arcsin\frac{2}{\pi} \approx 0.69, \qquad c_2 = \pi - \arcsin\frac{2}{\pi} \approx 2.45 .
\]
Both lie in the interval, illustrating that $c$ need not be unique.
\end{exoresolu}

\courssub{Bounding an Integral}

The proof of the mean value theorem contained a result that is extremely useful in its own right, especially when an integral cannot be evaluated exactly.

\begin{propriete}[Bounds for an integral]
If $f$ is continuous on $[a,b]$ and $m \le f(x) \le M$ for all $x\in[a,b]$, then
\[
m(b-a) \le \int_a^b f(x)\,dx \le M(b-a) .
\]
Taking $m$ and $M$ to be the minimum and maximum of $f$ on $[a,b]$ gives the sharpest such bounds. (Proof examinable in one line: integrate the inequality $m\le f(x)\le M$ term by term over $[a,b]$, using the fact that integration preserves order.)
\end{propriete}

\begin{popfigure}
\begin{center}
\begin{tikzpicture}[scale=1.15]
\draw[->] (0,0) -- (6.6,0) node[right] {$x$};
\draw[->] (0,0) -- (0,4.4) node[above] {$y$};
\draw[thick, popblue, domain=1:5.5, samples=80] plot (\x, {2.3 + 0.8*sin(deg(2*(\x-1))) + 0.1*(\x-1)});
\draw[dashed, thick, popgreen] (0.6,1.6) -- (5.9,1.6) node[right, popgreen] {$y=m$};
\draw[dashed, thick, poporange] (0.6,3.6) -- (5.9,3.6) node[right, poporange] {$y=M$};
\draw[popmuted, dashed] (1,0) node[below] {$a$} -- (1,1.6);
\draw[popmuted, dashed] (5.5,0) node[below] {$b$} -- (5.5,1.6);
\draw[<->, popdark] (6.15,0) -- (6.15,1.6) node[midway, right] {\scriptsize $m(b-a)$};
\draw[<->, popdark] (6.35,0) -- (6.35,3.6) node[midway, right] {\scriptsize $M(b-a)$};
\node[popblue] at (3.2,4.0) {$y=f(x)$};
\end{tikzpicture}
\legende{The area under the curve is trapped between the areas of the rectangles of heights $m$ and $M$: $m(b-a)\le\displaystyle\int_a^b f(x)\,dx \le M(b-a)$.}
\end{center}
\end{popfigure}

\begin{exemplebox}[Estimating an integral we cannot compute]
Show that $\displaystyle 1 \le \int_0^1 \sqrt{1+x^3}\,dx \le \sqrt{2}$.

\emph{Solution.} For $0\le x\le 1$, we have $0\le x^3\le 1$, so $1 \le \sqrt{1+x^3} \le \sqrt{2}$. Here $b-a=1$, and the bounding property gives immediately
\[
1\cdot(1-0) \le \int_0^1 \sqrt{1+x^3}\,dx \le \sqrt{2}\,(1-0), \quad\text{i.e.}\quad 1 \le \int_0^1 \sqrt{1+x^3}\,dx \le 1.414 .
\]
(The true value is about $1.11$ — comfortably inside our bounds. There is no elementary primitive for $\sqrt{1+x^3}$, so such estimates, or the trapezium rule, are the only options.)
\end{exemplebox}

\courssub{Applications: Averages in Real Life}

The mean value formula turns ``average'' questions about continuously varying quantities into integrals.

\begin{exemplebox}[Average speed, temperature, flow]
\textbf{(a) Average speed.} A car's speed during the first $10$ seconds of motion is $v(t)=3t$ (in $\mathrm{m\,s^{-1}}$, $t$ in seconds). Its average speed over $[0,10]$ is
\[
\bar{v} = \frac{1}{10}\int_0^{10} 3t\,dt = \frac{1}{10}\left[\frac{3t^2}{2}\right]_0^{10} = \frac{1}{10}(150) = 15\ \mathrm{m\,s^{-1}} .
\]
(Here $v$ is linear, so this equals $\frac{v(0)+v(10)}{2} = \frac{0+30}{2} = 15$ — as the warning box predicted.)

\textbf{(b) Average temperature.} The temperature in a town on a given day is modelled by $T(t) = 24 + 6\sin\dfrac{\pi t}{12}$ degrees Celsius, where $t$ is the time in hours after midnight. The mean temperature over the day $[0,24]$ is
\[
\bar{T} = \frac{1}{24}\int_0^{24}\left(24+6\sin\frac{\pi t}{12}\right)dt = \frac{1}{24}\left[24t - \frac{72}{\pi}\cos\frac{\pi t}{12}\right]_0^{24} = \frac{1}{24}\left(576 - \frac{72}{\pi}(1-1)\right) = 24^\circ\mathrm{C}.
\]
The sine term contributes nothing over a full period: the average equals the central value $24$.

\textbf{(c) Average flow of a river.} The rate of flow of water through a pumping station in Bamenda during a 6-hour operation is $Q(t) = 50 + 20t - 2t^2$ cubic metres per hour. The average rate of flow is
\[
\bar{Q} = \frac{1}{6}\int_0^6 (50+20t-2t^2)\,dt = \frac{1}{6}\Big[50t+10t^2-\tfrac{2}{3}t^3\Big]_0^6 = \frac{1}{6}(300+360-144) = 86\ \mathrm{m^3\,h^{-1}} .
\]
The total volume pumped is $\bar{Q}\times 6 = 516\ \mathrm{m^3}$: mean value times duration, exactly as the geometric interpretation promises.
\end{exemplebox}

\courssub{Synthesis: a Typical GCE Long Question}

GCE questions love to tour the whole chapter with a single function: area, then volume of revolution, then mean value, then a numerical check. Here is a complete model.

\begin{exoresolu}[Area, volume and mean value of one function]
Let $f(x) = \sqrt{x}$ for $0 \le x \le 4$, and let $R$ be the region bounded by the curve $y=f(x)$, the $x$-axis and the line $x=4$.

\textbf{(a)} Find the area of $R$.

\textbf{(b)} Find the volume of the solid generated when $R$ is rotated through $360^\circ$ about the $x$-axis.

\textbf{(c)} Find the mean value of $f$ on $[0,4]$ and the value of $c\in[0,4]$ such that $f(c)$ equals this mean value.

\textbf{(d)} Use the trapezium rule with $4$ strips to estimate $\displaystyle\int_0^4 \sqrt{x}\,dx$, and comment.
\tcblower
\textbf{Solution.}
\textbf{(a)} $\displaystyle A = \int_0^4 x^{1/2}\,dx = \left[\frac{2}{3}x^{3/2}\right]_0^4 = \frac{2}{3}(8) - 0 = \frac{16}{3}$ square units.

\textbf{(b)} $\displaystyle V = \pi\int_0^4 \big(\sqrt{x}\big)^2 dx = \pi\int_0^4 x\,dx = \pi\left[\frac{x^2}{2}\right]_0^4 = 8\pi$ cubic units.

\textbf{(c)} $\bar{f} = \dfrac{1}{4-0}\cdot\dfrac{16}{3} = \dfrac{4}{3}$. Solve $\sqrt{c} = \dfrac{4}{3}$: $c = \dfrac{16}{9} \approx 1.78$, which lies in $[0,4]$. \checkmark

\textbf{(d)} With $n=4$ strips, $h = \dfrac{4-0}{4} = 1$. The ordinates are $y_0=\sqrt{0}=0,\ y_1=1,\ y_2=\sqrt{2}\approx 1.4142,\ y_3=\sqrt{3}\approx 1.7321,\ y_4=2$.
\[
T = \frac{h}{2}\big[y_0 + y_4 + 2(y_1+y_2+y_3)\big] = \frac{1}{2}\big[0 + 2 + 2(1 + 1.4142 + 1.7321)\big] = \frac{1}{2}(10.2926) \approx 5.146 .
\]
The exact value is $\frac{16}{3}\approx 5.333$. The trapezium rule \emph{underestimates} here because $\sqrt{x}$ is concave (curving downwards) on $[0,4]$: each trapezium lies below the curve.
\end{exoresolu}

\begin{exercice}[Practice]
\begin{enumerate}
\item Find the mean value of $f(x) = 3x^2 - 2x$ on $[1,4]$, and the value(s) of $c$ in $[1,4]$ for which $f(c)$ equals this mean value.
\item Prove that $\displaystyle \frac{\pi}{9} \le \int_0^{\pi/6}\frac{dx}{1+\sin x} \le \frac{\pi}{6}$. [\emph{Hint:} bound $\sin x$ on $[0,\pi/6]$, then bound the integrand.]
\item The cost rate of running a machine is $C(t) = 2000 + 300\sqrt{t}$ francs CFA per hour, for $0\le t\le 9$ hours. Find the average hourly cost rate over the 9-hour shift.
\item For $f(x)=x^2$ on $[0,3]$: find the area under the curve, the volume when the region is rotated about the $x$-axis, the mean value of $f$, and the value of $c$ given by the mean value theorem.
\end{enumerate}
\end{exercice}

\begin{corrige}[Exercise 1]
$\displaystyle\int_1^4 (3x^2-2x)\,dx = \big[x^3 - x^2\big]_1^4 = (64-16)-(1-1) = 48$. Mean value: $\bar{f} = \dfrac{48}{4-1} = 16$. Solve $3c^2 - 2c = 16$: $3c^2-2c-16=0$, so $c = \dfrac{2\pm\sqrt{4+192}}{6} = \dfrac{2\pm 14}{6}$, giving $c=\dfrac{8}{3}$ or $c=-\dfrac{7}{3}$. Only $c=\dfrac{8}{3}\in[1,4]$ is retained.
\end{corrige}

\begin{corrige}[Exercise 2]
On $[0,\pi/6]$, the sine function is increasing, so $0 \le \sin x \le \tfrac12$, hence $1 \le 1+\sin x \le \tfrac32$, and taking reciprocals (all quantities positive) reverses the inequalities: $\dfrac{2}{3} \le \dfrac{1}{1+\sin x} \le 1$. Integrating over an interval of length $\dfrac{\pi}{6}$:
\[
\frac{2}{3}\cdot\frac{\pi}{6} \le \int_0^{\pi/6}\frac{dx}{1+\sin x} \le 1\cdot\frac{\pi}{6}, \qquad\text{i.e.}\qquad \frac{\pi}{9} \le \int_0^{\pi/6}\frac{dx}{1+\sin x} \le \frac{\pi}{6}.
\]
\end{corrige}

\begin{corrige}[Exercise 3]
$\displaystyle\int_0^9 (2000+300\sqrt{t})\,dt = \Big[2000t + 300\cdot\tfrac{2}{3}t^{3/2}\Big]_0^9 = \Big[2000t + 200t^{3/2}\Big]_0^9 = 18000 + 200(27) = 23400$. Average rate: $\bar{C} = \dfrac{23400}{9} = 2600$ francs CFA per hour.
\end{corrige}

\begin{corrige}[Exercise 4]
Area: $\displaystyle\int_0^3 x^2\,dx = \left[\frac{x^3}{3}\right]_0^3 = 9$ square units. Volume: $\displaystyle\pi\int_0^3 x^4\,dx = \pi\left[\frac{x^5}{5}\right]_0^3 = \frac{243\pi}{5}$ cubic units. Mean value: $\bar{f} = \dfrac{9}{3} = 3$. Solve $c^2 = 3$: $c = \sqrt{3}\approx 1.73$, which lies in $[0,3]$ (the root $-\sqrt{3}$ is rejected).
\end{corrige}

\courssub{Revision Map and Examination Strategy}

\begin{popfigure}
\begin{center}
\begin{tikzpicture}[
  box/.style={draw=popblue, fill=popblueL, rounded corners=2pt, text width=3.4cm, align=center, font=\small, minimum height=1cm},
  arr/.style={->, thick, popdark}
]
\node[box, align=center] (riem) at (0,4.5){\textbf{Riemann sums}\\ $\int_a^b f = \lim \sum f(x_i)\,\Delta x$};
\node[box, align=center] (area) at (5,4.5){\textbf{Areas}\\ $A=\int_a^b f\,dx$; $\int_c^d x\,dy$ for $y$-axis};
\node[box, align=center] (vol) at (10,4.5){\textbf{Volumes}\\ $V_x=\pi\int y^2\,dx$; $V_y=\pi\int x^2\,dy$};
\node[box, align=center] (cent) at (0,1.5){\textbf{Centroids}\\ $\bar{x}=\frac{1}{A}\int x f\,dx$};
\node[box, align=center] (trap) at (5,1.5){\textbf{Trapezium rule}\\ $T=\frac{h}{2}[y_0+y_n+2\sum\text{middle}]$};
\node[box, align=center] (mean) at (10,1.5){\textbf{Mean value}\\ $\bar f=\frac{1}{b-a}\int_a^b f$; $f(c)=\bar f$};
\draw[arr] (riem) -- (area);
\draw[arr] (area) -- (vol);
\draw[arr] (riem) -- (cent);
\draw[arr] (area) -- (trap);
\draw[arr] (area) -- (mean);
\end{tikzpicture}
\legende{Revision map of the chapter Integration: every application flows from the definite integral.}
\end{center}
\end{popfigure}

\begin{aretenir}[Choosing the right tool in the examination]
\begin{itemize}
\item \textbf{Area under $y=f(x)$ above the $x$-axis:} $\displaystyle A=\int_a^b f(x)\,dx$. If the curve dips below the axis, split the integral at the zeros of $f$ and take absolute values of the pieces.
\item \textbf{Area between curve and $y$-axis:} express $x$ in terms of $y$ and use $\displaystyle A=\int_c^d x\,dy$.
\item \textbf{Volume about the $x$-axis:} $\displaystyle V=\pi\int_a^b y^2\,dx$; \textbf{about the $y$-axis:} $\displaystyle V=\pi\int_c^d x^2\,dy$. Do not forget the $\pi$ and do not forget to \emph{square}.
\item \textbf{Centroid of a region under $y=f(x)$:} $\bar{x} = \dfrac{1}{A}\displaystyle\int_a^b x\,f(x)\,dx$, $\bar{y} = \dfrac{1}{2A}\displaystyle\int_a^b [f(x)]^2\,dx$.
\item \textbf{No elementary primitive, or data given in a table:} trapezium rule, $T = \dfrac{h}{2}\big[y_0 + y_n + 2(\text{sum of middle ordinates})\big]$.
\item \textbf{``Average value'' or ``mean value'':} divide the integral by $(b-a)$. If asked for $c$, solve $f(c)=\bar{f}$ and check $c\in[a,b]$.
\item \textbf{``Show that the integral lies between \dots'':} bound $f$ by its minimum and maximum on the interval and multiply by $(b-a)$.
\end{itemize}
\end{aretenir}

\begin{attention}[Final checklist of traps]
\begin{itemize}
\item Mean value $\neq$ $\dfrac{f(a)+f(b)}{2}$ unless $f$ is linear.
\item In the mean value theorem, reject solutions of $f(c)=\bar{f}$ that fall outside $[a,b]$.
\item A negative integral is not a negative area: areas are obtained from $\displaystyle\int_a^b |f(x)|\,dx$.
\item Units: areas in square units, volumes in cubic units; volumes of revolution carry a factor $\pi$.
\item In the trapezium rule, $h = \dfrac{b-a}{n}$ where $n$ is the number of \emph{strips}, not the number of ordinates (there are $n+1$ ordinates).
\end{itemize}
\end{attention}

\begin{savaistu}[Why averages matter beyond the classroom]
Hydrologists studying the Sanaga river express its discharge as a mean rate of flow — exactly $\dfrac{1}{b-a}\displaystyle\int_a^b Q(t)\,dt$ — to design dams and predict electricity output at hydroelectric stations. Meteorological services define the mean daily temperature the same way. Every time you read ``average speed'', ``average cost'' or ``average power'' about a quantity that varies continuously, the definite integral is working quietly behind the scenes.
\end{savaistu}

\newpage

\coursec{Integration activity}

\coursec{Integrative Activity: Rainwater Harvesting Tank at Lyc\'ee Bilingue de Bafoussam}

\courssub{Learning Objectives}
\begin{aretenir}[Learning Objectives]
This integrative activity consolidates your mastery of the Integration chapter (PMS04) through a realistic engineering context. By the end, you should be able to:
\begin{itemize}
    \item Express a definite integral as the limit of a Riemann sum (Lesson 52).
    \item Compute the \cle{cross-sectional area} of a plane region bounded by a curve and the coordinate axes (Lessons 53, 54, 55).
    \item Calculate the \cle{volume of a solid of revolution} about the $x$-axis and the $y$-axis using the disc/washer method (Lessons 56, 57).
    \item Determine the \cle{centroid (centre of mass)} of a solid of revolution, applying the formulae for $\bar{x}$ (Lesson 58).
    \item Apply the \cle{trapezium rule} for numerical integration to approximate a definite integral from discrete data (Lesson 59).
    \item Use the \cle{Mean Value Theorem for Integrals} to find the average value of a function over an interval (Lesson 60).
    \item Communicate mathematical reasoning clearly, with labelled sketches and justifications for every formula used.
\end{itemize}
\end{aretenir}

\courssub{Theoretical Recall: Riemann Sums and Area Definitions}
\begin{propriete}[Definite Integral as Limit of Riemann Sums -- Lesson 52]
Let $f$ be a function defined on $[a,b]$. Partition $[a,b]$ into $n$ subintervals of width $\Delta x = \frac{b-a}{n}$. Choose sample points $x_i^* \in [x_{i-1}, x_i]$. The \cle{Riemann sum} is $\sum_{i=1}^n f(x_i^*)\Delta x$. The definite integral is the limit:
\[
\int_a^b f(x)\,\mathrm{d}x = \lim_{n\to\infty} \sum_{i=1}^n f(x_i^*)\Delta x.
\]
This limit exists if $f$ is continuous on $[a,b]$ (or has finitely many jump discontinuities). \textbf{Examinable:} Be able to write the Riemann sum for a given integral and recognise a limit of a sum as a definite integral.
\end{propriete}

\begin{propriete}[Areas Bounded by Curves and Axes -- Lessons 53, 54, 55]
\begin{itemize}
    \item Area between $y=f(x)\ge 0$, $x$-axis, $x=a$, $x=b$: $\displaystyle A = \int_a^b f(x)\,\mathrm{d}x$.
    \item Area between $x=g(y)\ge 0$, $y$-axis, $y=c$, $y=d$: $\displaystyle A = \int_c^d g(y)\,\mathrm{d}y$.
    \item If the curve crosses the axis, split the interval and take absolute values.
\end{itemize}
\end{propriete}

\begin{propriete}[Volumes of Revolution -- Lessons 56, 57]
\begin{itemize}
    \item About $x$-axis (disc method): $V = \pi \int_a^b [f(x)]^2\,\mathrm{d}x$ for $y=f(x)\ge 0$.
    \item About $y$-axis (disc method): $V = \pi \int_c^d [g(y)]^2\,\mathrm{d}y$ for $x=g(y)\ge 0$.
    \item Washer method for region between two curves.
\end{itemize}
\textbf{Proof (examinable):} Volume $\approx \sum \pi y_i^2 \Delta x \to \pi \int y^2\,\mathrm{d}x$.
\end{propriete}

\begin{propriete}[Centroid of a Solid of Revolution -- Lesson 58]
For a solid generated by rotating $y=f(x)\ge 0$ on $[a,b]$ about the $x$-axis, with uniform density:
\[
\bar{x} = \frac{\int_a^b x [f(x)]^2\,\mathrm{d}x}{\int_a^b [f(x)]^2\,\mathrm{d}x}, \qquad \bar{y}=0,\ \bar{z}=0 \text{ (by symmetry)}.
\]
For a thin shell (surface of revolution), the formula differs: $\bar{x} = \frac{\int x y \sqrt{1+(y')^2}\,\mathrm{d}x}{\int y \sqrt{1+(y')^2}\,\mathrm{d}x}$.
\end{propriete}

\begin{propriete}[Trapezium Rule -- Lesson 59]
For $n$ strips of width $h = \frac{b-a}{n}$, $x_i = a+ih$:
\[
\int_a^b f(x)\,\mathrm{d}x \approx \frac{h}{2}\left[ f(x_0) + f(x_n) + 2\sum_{i=1}^{n-1} f(x_i) \right].
\]
Error is proportional to $h^2$; exact for linear functions.
\end{propriete}

\begin{propriete}[Mean Value Theorem for Integrals -- Lesson 60]
If $f$ is continuous on $[a,b]$, $\exists\, c\in[a,b]$ such that
\[
f(c) = \frac{1}{b-a}\int_a^b f(x)\,\mathrm{d}x.
\]
The right-hand side is the \cle{mean (average) value} of $f$ on $[a,b]$.
\end{propriete}

\courssub{Context: Lyc\'ee Bilingue de Bafoussam -- Rainwater Harvesting Project}
\begin{experience}[Project Description]
The administration of \emph{Lyc\'ee Bilingue de Bafoussam} (West Region, Cameroon) launches a project to improve water autonomy during the dry season (November--March). The region receives heavy rainfall from April to October, making rainwater harvesting viable. The project committee designs a horizontal cylindrical tank whose \textbf{inner wall profile} (longitudinal cross-section through the axis) follows the curve
\[
\mathcal{C}:\quad y = 0.5\sqrt{x} = \frac{1}{2}x^{1/2}, \qquad 0 \le x \le 4
\]
where $x$ and $y$ are in \textbf{metres}. The tank is generated by rotating this profile completely about the $x$-axis (the horizontal axis). The tank rests on two concrete support cradles.

As members of the \emph{Young Engineers Club}, your group must produce a technical study before construction:
\begin{enumerate}
    \item Determine the cross-sectional area of the tank's longitudinal profile (area between $\mathcal{C}$ and the $x$-axis). This estimates the waterproofing membrane needed for the inner surface development.
    \item Calculate the exact \textbf{capacity} (volume) of the tank in $\mathrm{m}^3$ and litres ($1\,\mathrm{m}^3 = 1000\,\mathrm{L}$).
    \item Locate the \textbf{centre of mass of the empty tank} along its axis ($x$-coordinate $\bar{x}$), assuming uniform density and negligible wall thickness so that the centroid of the \emph{solid of revolution} approximates the centre of mass of the tank structure. This determines cradle placement for balanced loading.
    \item The construction team provides measured internal radii $y$ at 9 equally spaced stations (8 strips). Use the \textbf{trapezium rule} to approximate the volume from this discrete data and compare with the exact value.
    \item Hydrological data indicates that during the rainy season, the water depth at the tank's midpoint ($x=2$) varies with time $t$ (weeks) according to $h(t)$, whose \textbf{mean value} over the 12-week peak period is $0.65\,\mathrm{m}$. Use the Mean Value Theorem for Integrals to interpret this result and state the average depth.
    \item \textbf{(Extension)} Express the volume integral as a limit of Riemann sums. Also, set up the integral for the area bounded by the curve and the $y$-axis, and the volume if the profile were rotated about the $y$-axis.
\end{enumerate}
Your report must contain labelled sketches, all derivations, numerical results (3 significant figures), and a conclusion advising the committee.
\end{experience}

\begin{popfigure}
\begin{tikzpicture}[scale=0.9, >=stealth]
% Axes
\draw[->] (-0.5,0) -- (4.5,0) node[right] {$x\ (\mathrm{m})$};
\draw[->] (0,-1.2) -- (0,1.2) node[above] {$y\ (\mathrm{m})$};
% Curve y = 0.5*sqrt(x)
\draw[popblue, thick, domain=0:4, smooth, variable=\x] plot ({\x}, {0.5*sqrt(max(0,\x))});
% Shaded area (cross-section)
\fill[popblueL, opacity=0.3] (0,0) -- plot[domain=0:4, smooth, variable=\x] ({\x}, {0.5*sqrt(max(0,\x))}) -- (4,0) -- cycle;
% Radius at x=2
\draw[poporange, dashed] (2,0) -- (2,{0.5*sqrt(max(0,2))});
\draw[poporange, <->] (2.1,0) -- (2.1,{0.5*sqrt(max(0,2))}) node[midway, right] {$R(x)=0.5\sqrt{x}$};
% Trapezium strips (8 strips, h=0.5)
\foreach \i in {0,...,7} {
    \pgfmathsetmacro{\xa}{\i*0.5}
    \pgfmathsetmacro{\xb}{(\i+1)*0.5}
    \pgfmathsetmacro{\ya}{0.5*sqrt(\xa)}
    \pgfmathsetmacro{\yb}{0.5*sqrt(\xb)}
    \draw[popgreen!50!black, fill=popgreenL, opacity=0.2] (\xa,0) -- (\xa,\ya) -- (\xb,\yb) -- (\xb,0) -- cycle;
}
% Centroid marker
\draw[poppurple, thick] (2.667, -0.1) -- (2.667, 0.1) node[below] {$\bar{x}=8/3$};
% Labels
\node[popdark] at (2, -0.5) {\textbf{Longitudinal Profile (Cross-Section)}};
\node[popdark] at (4.2, 1.0) {$y=0.5\sqrt{x}$};
% 3D suggestion: rotation arrow
\draw[poppurple, thick, ->] (4.5, 0.5) arc (0:180:0.5 and 0.2);
\node[poppurple, right] at (4.5, 0.5) {\small Rotation about $x$-axis};
\end{tikzpicture}
\legende{Longitudinal profile of the tank (shaded blue area) rotated about the $x$-axis to generate the 3D volume. The 8 trapezium strips for numerical integration are shown in green. The centroid $\bar{x}=8/3\,\mathrm{m}$ is marked in purple.}
\end{popfigure}

\courssub{Tasks}

\begin{enumerate}
    \item \textbf{Cross-Sectional Area (Lessons 53, 54):} Sketch the region bounded by $y=0.5\sqrt{x}$, the $x$-axis, $x=0$ and $x=4$. Set up and evaluate the definite integral for the area $A$. State the units.
    \item \textbf{Area Bounded by the $y$-axis (Lesson 55):} Express $x$ in terms of $y$ and set up the integral for the area between the curve, the $y$-axis, $y=0$ and $y=1$. Evaluate it and verify it equals the area from Task 1.
    \item \textbf{Volume of Revolution about $x$-axis -- Capacity (Lesson 56):} Write the formula for the volume $V$ generated by rotating the region about the $x$-axis. Evaluate the integral exactly. Convert to litres.
    \item \textbf{Volume of Revolution about $y$-axis (Lesson 57):} Set up the integral for the volume if the same profile (for $0\le y\le 1$) were rotated about the $y$-axis. Evaluate it.
    \item \textbf{Riemann Sum Formulation (Lesson 52):} Write the Riemann sum for $\int_0^4 y^2\,\mathrm{d}x$ using a regular partition of $[0,4]$ into $n$ subintervals and right endpoints. State the limit that gives the exact volume.
    \item \textbf{Centre of Mass (Centroid) of the Solid (Lesson 58):} Write the formula for the $x$-coordinate of the centroid $\bar{x}$ of a solid of revolution about the $x$-axis. Compute $\bar{x}$ exactly and give a decimal approximation (3 s.f.). Explain why $\bar{y}=0$ and $\bar{z}=0$. Mark $\bar{x}$ on your sketch.
    \item \textbf{Numerical Verification -- Trapezium Rule (Lesson 59):} The table below gives measured radii $y$ at 9 stations.
    \begin{center}
    \begin{tabularx}{\linewidth}{|c|*{9}{X|}}\hline
    \rowcolor{popblueL}
    $x\ (\mathrm{m})$ & 0.0 & 0.5 & 1.0 & 1.5 & 2.0 & 2.5 & 3.0 & 3.5 & 4.0 \\ \hline
    $y\ (\mathrm{m})$ & 0.000 & 0.354 & 0.500 & 0.612 & 0.707 & 0.791 & 0.866 & 0.935 & 1.000 \\ \hline
    \end{tabularx}
    \end{center}
    Use the trapezium rule with $h=0.5$ to approximate $\int_0^4 y^2\,\mathrm{d}x$, then compute the approximate volume $V_{\text{trap}}$. Calculate the percentage error relative to the exact volume.
    \item \textbf{Mean Value Theorem Application (Lesson 60):} Let $h(t)$ be the water depth (metres) at the tank's midpoint ($x=2$) during the 12-week peak rainy season ($t\in[0,12]$). The Mean Value Theorem for Integrals guarantees a $c\in[0,12]$ such that $h(c) = \frac{1}{12}\int_0^{12} h(t)\,\mathrm{d}t$. If the average value (the right-hand side) is $0.65\,\mathrm{m}$, what is the \textbf{average depth of water} during the term? Explain the physical meaning for the water management committee.
\end{enumerate}

\courssub{Model Answer}

\begin{exoresolu}[Task 1: Cross-Sectional Area (Lessons 53, 54)]
The region is bounded above by $y = \frac{1}{2}x^{1/2}$, below by the $x$-axis ($y=0$), and laterally by $x=0$ and $x=4$. Since $f(x)\ge 0$ on $[0,4]$,
\[
A = \int_{0}^{4} y\,\mathrm{d}x = \int_{0}^{4} \frac{1}{2}x^{1/2}\,\mathrm{d}x.
\]
\[
A = \frac{1}{2} \left[ \frac{x^{3/2}}{3/2} \right]_{0}^{4} = \frac{1}{2} \cdot \frac{2}{3} \left[ x^{3/2} \right]_{0}^{4} = \frac{1}{3} \left( 4^{3/2} - 0 \right).
\]
Since $4^{3/2} = (\sqrt{4})^3 = 2^3 = 8$,
\[
A = \frac{8}{3}\ \mathrm{m}^2 \approx 2.67\ \mathrm{m}^2.
\]
\end{exoresolu}

\begin{aretenir}[Area under curve]
The area between a curve $y=f(x)\ge 0$ and the $x$-axis from $x=a$ to $x=b$ is $\int_a^b f(x)\,\mathrm{d}x$. Always sketch first to confirm limits and that the curve does not cross the axis.
\end{aretenir}

\begin{exoresolu}[Task 2: Area Bounded by the $y$-axis (Lesson 55)]
The curve is $y = \frac{1}{2}\sqrt{x} \implies \sqrt{x} = 2y \implies x = 4y^2$ for $0\le y\le 1$.
Area between the curve, $y$-axis ($x=0$), $y=0$, $y=1$:
\[
A = \int_{0}^{1} x\,\mathrm{d}y = \int_{0}^{1} 4y^2\,\mathrm{d}y = 4\left[ \frac{y^3}{3} \right]_0^1 = \frac{4}{3}\ \mathrm{m}^2.
\]
\textbf{Wait!} This is \textbf{not} the same region. The region in Task 1 is bounded by $x=4$ (vertical line), while this region is bounded by $y=1$ (horizontal line). They are different regions! The area in Task 1 is $\frac{8}{3}\,\mathrm{m}^2$; the area bounded by the $y$-axis and $y=1$ is $\frac{4}{3}\,\mathrm{m}^2$. The total rectangle $0\le x\le 4$, $0\le y\le 1$ has area $4\,\mathrm{m}^2$. Indeed $\frac{8}{3} + \frac{4}{3} = 4$.
\end{exoresolu}

\begin{attention}[Area bounded by $y$-axis vs $x$-axis]
``Area bounded by the $y$-axis'' means the region is bounded on the left by $x=0$ (the $y$-axis). The limits are on $y$. The region in Task 1 is bounded by the $x$-axis ($y=0$) and the vertical line $x=4$. They are complementary regions in the rectangle $[0,4]\times[0,1]$. Read the wording carefully in exams.
\end{attention}

\begin{exoresolu}[Task 3: Volume of Revolution about $x$-axis -- Capacity (Lesson 56)]
Rotating the region about the $x$-axis generates a solid. Using the \cle{disc method}, a thin slice at $x$ of thickness $\mathrm{d}x$ has volume $\mathrm{d}V = \pi y^2\,\mathrm{d}x$.
\[
V = \pi \int_{0}^{4} y^2\,\mathrm{d}x = \pi \int_{0}^{4} \left( \frac{1}{2}x^{1/2} \right)^2\,\mathrm{d}x = \pi \int_{0}^{4} \frac{1}{4}x\,\mathrm{d}x.
\]
\[
V = \frac{\pi}{4} \left[ \frac{x^2}{2} \right]_{0}^{4} = \frac{\pi}{4} \cdot \frac{16}{2} = \frac{\pi}{4} \cdot 8 = 2\pi\ \mathrm{m}^3.
\]
Exact capacity: $\boldsymbol{2\pi\ \mathrm{m}^3}$.
Decimal: $V \approx 6.283\ \mathrm{m}^3$ (3 s.f.).
In litres: $V = 2\pi \times 1000 \approx \boldsymbol{6280\ \mathrm{L}}$ (3 s.f.).
\end{exoresolu}

\begin{propriete}[Volume of Revolution about $x$-axis -- Disc Method]
If a region bounded by $y=f(x)\ge 0$, the $x$-axis, $x=a$, $x=b$ is rotated about the $x$-axis, the volume is $V = \pi \int_a^b [f(x)]^2\,\mathrm{d}x$. \textbf{Examinable proof:} Limit of sum of cylindrical discs $\sum \pi y_i^2 \Delta x$.
\end{propriete}

\begin{exoresolu}[Task 4: Volume of Revolution about $y$-axis (Lesson 57)]
The profile for $0\le y\le 1$ is $x = 4y^2$. Rotating this about the $y$-axis:
\[
V_y = \pi \int_{0}^{1} x^2\,\mathrm{d}y = \pi \int_{0}^{1} (4y^2)^2\,\mathrm{d}y = \pi \int_{0}^{1} 16y^4\,\mathrm{d}y = 16\pi \left[ \frac{y^5}{5} \right]_0^1 = \frac{16\pi}{5}\ \mathrm{m}^3 \approx 10.1\ \mathrm{m}^3.
\]
This is a different solid (a ``bowl'' shape) with larger volume.
\end{exoresolu}

\begin{exoresolu}[Task 5: Riemann Sum Formulation (Lesson 52)]
We want $\int_0^4 y^2\,\mathrm{d}x = \int_0^4 \frac{1}{4}x\,\mathrm{d}x$. Partition $[0,4]$ into $n$ equal subintervals: $\Delta x = \frac{4}{n}$. Right endpoints: $x_i = i\Delta x = \frac{4i}{n}$ for $i=1,\dots,n$.
Riemann sum:
\[
S_n = \sum_{i=1}^n f(x_i)\Delta x = \sum_{i=1}^n \frac{1}{4}\left(\frac{4i}{n}\right) \cdot \frac{4}{n} = \sum_{i=1}^n \frac{4i}{n^2} = \frac{4}{n^2}\sum_{i=1}^n i = \frac{4}{n^2}\cdot\frac{n(n+1)}{2} = \frac{2(n+1)}{n}.
\]
The definite integral is the limit:
\[
\int_0^4 \frac{1}{4}x\,\mathrm{d}x = \lim_{n\to\infty} S_n = \lim_{n\to\infty} \frac{2(n+1)}{n} = 2.
\]
Then $V = \pi \times 2 = 2\pi\ \mathrm{m}^3$, confirming Task 3.
\end{exoresolu}

\begin{exoresolu}[Task 6: Centre of Mass (Centroid) of the Solid (Lesson 58)]
For a solid of revolution about the $x$-axis with uniform density, the centroid lies on the axis of symmetry, so $\bar{y}=0$, $\bar{z}=0$. The $x$-coordinate is the first moment of volume divided by the total volume:
\[
\bar{x} = \frac{\int_0^4 x\,(\pi y^2)\,\mathrm{d}x}{\int_0^4 \pi y^2\,\mathrm{d}x} = \frac{\int_0^4 x y^2\,\mathrm{d}x}{\int_0^4 y^2\,\mathrm{d}x}.
\]
We have $y^2 = \frac{1}{4}x$. Numerator:
\[
\int_0^4 x \cdot \frac{1}{4}x\,\mathrm{d}x = \frac{1}{4} \int_0^4 x^2\,\mathrm{d}x = \frac{1}{4} \left[ \frac{x^3}{3} \right]_0^4 = \frac{1}{4} \cdot \frac{64}{3} = \frac{16}{3}.
\]
Denominator (from Task 3): $\int_0^4 y^2\,\mathrm{d}x = \frac{V}{\pi} = 2$.
\[
\bar{x} = \frac{16/3}{2} = \frac{8}{3}\ \mathrm{m} \approx 2.67\ \mathrm{m}.
\]
The support cradles should be placed symmetrically around $x = \frac{8}{3}\,\mathrm{m} \approx 2.67\,\mathrm{m}$ from the left end ($x=0$) to balance the empty tank.
\end{exoresolu}

\begin{attention}[Centroid vs Centre of Mass]
For a homogeneous solid, Centre of Mass = Centroid of Volume. The formula $\bar{x} = \frac{\int x y^2\,\mathrm{d}x}{\int y^2\,\mathrm{d}x}$ applies to the \textbf{solid volume}. If the tank were a thin shell, we would use the surface of revolution formula $\bar{x} = \frac{\int x y \sqrt{1+(y')^2}\,\mathrm{d}x}{\int y \sqrt{1+(y')^2}\,\mathrm{d}x}$. The problem statement allows the solid approximation.
\end{attention}

\begin{exoresolu}[Task 7: Trapezium Rule Verification (Lesson 59)]
We approximate $I = \int_0^4 y^2\,\mathrm{d}x$. The table gives $y$; we compute $y^2$ to 4 decimal places:
\begin{center}
\begin{tabularx}{\linewidth}{|c|*{9}{X|}}\hline
\rowcolor{popblueL}
$x$ & 0.0 & 0.5 & 1.0 & 1.5 & 2.0 & 2.5 & 3.0 & 3.5 & 4.0 \\ \hline
$y$ & 0.000 & 0.354 & 0.500 & 0.612 & 0.707 & 0.791 & 0.866 & 0.935 & 1.000 \\ \hline
$y^2$ & 0.0000 & 0.1253 & 0.2500 & 0.3745 & 0.4998 & 0.6257 & 0.7500 & 0.8742 & 1.0000 \\ \hline
\end{tabularx}
\end{center}
$h = 0.5$, $n=8$ strips. Trapezium Rule:
\[
I \approx \frac{h}{2} \left[ y_0^2 + y_8^2 + 2\sum_{i=1}^{7} y_i^2 \right].
\]
Sum of interior ordinates ($y_1^2$ to $y_7^2$):
$0.1253 + 0.2500 + 0.3745 + 0.4998 + 0.6257 + 0.7500 + 0.8742 = 3.4995$.
\[
I \approx \frac{0.5}{2} \left[ 0.0000 + 1.0000 + 2(3.4995) \right] = 0.25 \left[ 1.0000 + 6.9990 \right] = 0.25 \times 7.9990 = 1.99975.
\]
Exact value of $\int_0^4 y^2\,\mathrm{d}x = 2$ (from Task 3).
Approximate Volume:
\[
V_{\text{trap}} = \pi \times 1.99975 \approx 6.2823\ \mathrm{m}^3.
\]
Percentage Error:
\[
\%\text{ Error} = \left| \frac{V_{\text{trap}} - V_{\text{exact}}}{V_{\text{exact}}} \right| \times 100\% = \left| \frac{6.2823 - 6.2832}{6.2832} \right| \times 100\% \approx 0.014\%.
\]
\end{exoresolu}

\begin{methode}[Trapezium Rule for $\int_a^b f(x)\,\mathrm{d}x$ with $n$ strips]
\begin{enumerate}
    \item Calculate strip width $h = \frac{b-a}{n}$.
    \item List $x$-values: $x_0=a, x_1=a+h, \dots, x_n=b$.
    \item Evaluate $f(x_i)$ (here $f(x)=y^2$).
    \item Apply formula: $\int_a^b f(x)\,\mathrm{d}x \approx \frac{h}{2}[f(x_0)+f(x_n)+2\sum_{i=1}^{n-1}f(x_i)]$.
    \item Multiply by $\pi$ for volume of revolution.
\end{enumerate}
The small error ($<0.02\%$) confirms the exact calculus result and validates the measured radii.
\end{methode}

\begin{exoresolu}[Task 8: Mean Value Theorem for Integrals (Lesson 60)]
The Mean Value Theorem for Integrals states: If $h$ is continuous on $[a,b]$, there exists at least one $c \in [a,b]$ such that
\[
h(c) = \frac{1}{b-a} \int_a^b h(t)\,\mathrm{d}t.
\]
The quantity on the right is the \cle{mean value (average value)} of the function on the interval.
Here, $a=0$, $b=12$ (weeks), and the mean value is given as $0.65\,\mathrm{m}$.
\[
\frac{1}{12-0} \int_0^{12} h(t)\,\mathrm{d}t = 0.65\ \mathrm{m}.
\]
\textbf{Interpretation:} The average depth of water at the midpoint of the tank over the 12-week peak rainy season is exactly \textbf{0.65 m}.
This means that if the water level were held constant at $0.65\,\mathrm{m}$ for the entire 12 weeks, the total volume of water that would have passed through (or been stored at) that midpoint is the same as the actual varying volume. For the management committee: this average depth corresponds to an average cross-sectional area of water, allowing estimation of the average stored volume during the term for planning usage (e.g., irrigation, sanitation).
\end{exoresolu}

\begin{savaistu}[Cameroon Context]
In Bafoussam, the rainy season lasts about 7--8 months. A tank of $\approx 6300\,\mathrm{L}$ can supply a school of 500 students with $\approx 12\,\mathrm{L}$/day/student for about 10 days without rain. Knowing the \emph{average} depth helps size the pump and schedule rationing during dry spells within the season.
\end{savaistu}

\courssub{Extension Exercise}
\begin{exercice}[External Painting and Surface Area]
The tank is to be painted externally. The outer profile is $y = 0.5\sqrt{x} + 0.02$ (wall thickness $2\,\mathrm{cm}$) for $0 \le x \le 4$, rotated about the $x$-axis.
\begin{enumerate}
    \item Set up the integral for the \textbf{outer curved surface area} (area of revolution about $x$-axis). Do not evaluate.
    \item If one litre of paint covers $10\,\mathrm{m}^2$, estimate the paint required using the trapezium rule with 4 strips on the surface area integral. Comment on any difficulty.
\end{enumerate}
\end{exercice}

\begin{corrige}[Extension Exercise]
\begin{enumerate}
    \item Surface Area $S = 2\pi \int_0^4 y \sqrt{1 + \left(\frac{\mathrm{d}y}{\mathrm{d}x}\right)^2}\,\mathrm{d}x$.
    Here $y = 0.5\sqrt{x} + 0.02 = 0.5 x^{1/2} + 0.02$.
    $\frac{\mathrm{d}y}{\mathrm{d}x} = 0.25 x^{-1/2}$.
    $S = 2\pi \int_0^4 (0.5\sqrt{x} + 0.02) \sqrt{1 + \frac{0.0625}{x}}\,\mathrm{d}x$.
    \item Trapezium rule with $n=4$, $h=1$. $x$-values: $0, 1, 2, 3, 4$.
    Integrand $f(x) = (0.5\sqrt{x} + 0.02) \sqrt{1 + 0.0625/x}$.
    At $x=0$, $f(x) \to \infty$ (improper integral). A real tank would have a hemispherical cap at the pointed end; the model breaks down at the vertex. Start at $x=0.01$ or use a modified profile. This highlights the importance of checking the domain of the integrand in practical problems.
\end{enumerate}
\end{corrige}

\newpage

\coursec{Methods and worked examples}

\coursec{Integration}

\courssub{The Definite Integral as a Limit of Riemann Sums}

\begin{definition}[Partition and Riemann Sum]
Let $f$ be a function defined on a closed interval $[a,b]$. A \cle{partition} $P$ of $[a,b]$ is a finite set of points $\{x_0, x_1, \dots, x_n\}$ such that $a = x_0 < x_1 < \dots < x_n = b$. The \cle{norm} of the partition is $\|P\| = \max_{1 \le i \le n} (x_i - x_{i-1})$. Choosing a sample point $x_i^* \in [x_{i-1}, x_i]$ for each subinterval, the \cle{Riemann sum} is
\[
S(P) = \sum_{i=1}^{n} f(x_i^*)\, \Delta x_i, \qquad \Delta x_i = x_i - x_{i-1}.
\]
If $f$ is integrable on $[a,b]$, the definite integral is the limit of Riemann sums as the norm tends to zero:
\[
\int_a^b f(x)\,dx = \lim_{\|P\|\to 0} \sum_{i=1}^{n} f(x_i^*)\, \Delta x_i.
\]
\end{definition}

\begin{propriete}[Regular Partition and Right Endpoints]
For a regular partition into $n$ equal subintervals, $\Delta x = \dfrac{b-a}{n}$ and the right endpoints are $x_i = a + i\Delta x$ ($i=1,\dots,n$). Then
\[
\int_a^b f(x)\,dx = \lim_{n\to\infty} \sum_{i=1}^{n} f(a + i\Delta x)\,\Delta x.
\]
Useful summation formulas:
\[
\sum_{i=1}^{n} i = \frac{n(n+1)}{2},\quad
\sum_{i=1}^{n} i^2 = \frac{n(n+1)(2n+1)}{6},\quad
\sum_{i=1}^{n} i^3 = \left[\frac{n(n+1)}{2}\right]^2.
\]
\end{propriete}

\begin{methode}[Evaluating a Definite Integral as a Limit of Riemann Sums]
\begin{enumerate}
\item Write $\Delta x = \dfrac{b-a}{n}$ and $x_i = a + i\Delta x$ (right endpoints) or $x_i = a + (i-1)\Delta x$ (left endpoints).
\item Form the Riemann sum $S_n = \displaystyle\sum_{i=1}^{n} f(x_i)\,\Delta x$.
\item Simplify $S_n$ using the summation formulas above; express everything in powers of $n$.
\item Take the limit $n\to\infty$: terms containing $\frac{1}{n}, \frac{1}{n^2}, \dots$ vanish.
\end{enumerate}
\textbf{Reflex:} $\displaystyle\lim_{n\to\infty}\frac{1}{n}\sum_{i=1}^n f\!\left(\frac{i}{n}\right) = \int_0^1 f(x)\,dx$. More generally, $\displaystyle\lim_{n\to\infty}\frac{b-a}{n}\sum_{i=1}^n f\!\left(a+\frac{i(b-a)}{n}\right) = \int_a^b f(x)\,dx$.
\end{methode}

\begin{exemplebox}[Riemann sum for a polynomial]
Using Riemann sums, show that $\displaystyle\int_0^2 (x^2 + 2x)\,dx = \frac{20}{3}$.
\tcblower
\textbf{Step 1: partition.} $a=0$, $b=2$, $\Delta x = \frac{2}{n}$, $x_i = \frac{2i}{n}$.

\textbf{Step 2: form the sum.}
\begin{align*}
S_n &= \sum_{i=1}^{n} \left[ \left(\frac{2i}{n}\right)^2 + 2\left(\frac{2i}{n}\right) \right] \frac{2}{n} \\
&= \sum_{i=1}^{n} \left( \frac{4i^2}{n^2} + \frac{4i}{n} \right) \frac{2}{n} \\
&= \frac{8}{n^3}\sum_{i=1}^{n} i^2 + \frac{8}{n^2}\sum_{i=1}^{n} i \\
&= \frac{8}{n^3}\cdot\frac{n(n+1)(2n+1)}{6} + \frac{8}{n^2}\cdot\frac{n(n+1)}{2}.
\end{align*}

\textbf{Step 3: take the limit.}
\begin{align*}
S_n &= \frac{8}{6}\left(1+\frac{1}{n}\right)\left(2+\frac{1}{n}\right) + 4\left(1+\frac{1}{n}\right) \\
&\xrightarrow[n\to\infty]{} \frac{4}{3}\cdot 1 \cdot 2 + 4\cdot 1 = \frac{8}{3} + 4 = \frac{20}{3}.
\end{align*}
\textbf{Verification:} $\displaystyle\int_0^2 (x^2+2x)\,dx = \left[\frac{x^3}{3}+x^2\right]_0^2 = \frac{8}{3}+4 = \frac{20}{3}$. \checkmark
\end{exemplebox}

\begin{exemplebox}[Recognising a limit as a definite integral]
Express $\displaystyle\lim_{n\to\infty} \frac{1}{n}\sum_{i=1}^{n} \sqrt{1 + \frac{i}{n}}$ as a definite integral and evaluate it.
\tcblower
The sum has the form $\frac{1}{n}\sum f(i/n)$ with $f(x) = \sqrt{1+x}$ on $[0,1]$. Hence
\[
\lim_{n\to\infty} \frac{1}{n}\sum_{i=1}^{n} \sqrt{1 + \frac{i}{n}} = \int_0^1 \sqrt{1+x}\,dx.
\]
Evaluate: let $u = 1+x$, $du = dx$; limits $1\to 2$.
\[
\int_1^2 \sqrt{u}\,du = \left[\frac{2}{3}u^{3/2}\right]_1^2 = \frac{2}{3}(2\sqrt{2} - 1).
\]
\end{exemplebox}

\begin{popfigure}
\begin{tikzpicture}[scale=0.8]
\draw[->] (-0.5,0) -- (5,0) node[right] {$x$};
\draw[->] (0,-0.5) -- (0,3) node[above] {$y$};
\draw[domain=0:4, smooth, variable=\x, popblue, thick] plot ({\x}, {0.2*\x*\x + 0.5});
\foreach \i in {1,...,8} {
  \draw[poporange, fill=poporange!30] ({0.5*\i}, 0) rectangle ({0.5*\i-0.5}, {0.2*(0.5*\i)^2 + 0.5});
  \draw[poporange, dashed] ({0.5*\i}, 0) -- ({0.5*\i}, {0.2*(0.5*\i)^2 + 0.5});
}
\foreach \x in {0,1,2,3,4} \draw (\x, -0.1) -- (\x, 0.1) node[below] {$\x$};
\foreach \y in {1,2} \draw (-0.1, \y) -- (0.1, \y) node[left] {$\y$};
\node at (2.5, 2.5) {$f(x) = 0.2x^2 + 0.5$};
\node[poporange] at (2, -0.8) {Right-endpoint rectangles ($n=8$)};
\end{tikzpicture}
\legende{Right-endpoint Riemann sum approximating $\int_0^4 f(x)\,dx$}
\end{popfigure}

\courssub{Areas Bounded by Curves and the Coordinate Axes}

\begin{definition}[Area between a Curve and the $x$-axis]
Let $f$ be continuous on $[a,b]$. The \cle{area} of the region bounded by $y=f(x)$, the $x$-axis, and the lines $x=a$, $x=b$ is
\[
A = \int_a^b |f(x)|\,dx.
\]
If $f(x) \ge 0$ on $[a,b]$, $A = \int_a^b f(x)\,dx$. If $f(x) \le 0$, $A = -\int_a^b f(x)\,dx$. If $f$ changes sign, split the interval at the zeros of $f$ and add the absolute values of the integrals on each subinterval.
\end{definition}

\begin{methode}[Area between $y=f(x)$ and the $x$-axis]
\begin{enumerate}
\item Sketch the curve on $[a,b]$ and solve $f(x)=0$ to find all zeros inside the interval.
\item On each subinterval where $f$ does not change sign, compute $\int f(x)\,dx$.
\item Take the absolute value of each piece and sum them: $A = \sum \left|\int f(x)\,dx\right|$.
\item For the area between two curves $y=f(x)$ (upper) and $y=g(x)$ (lower) on $[a,b]$: $A = \int_a^b \big(f(x)-g(x)\big)\,dx$.
\end{enumerate}
\textbf{Reflex:} an integral can be negative; an \emph{area} never is. Always check the sign of $f$ before integrating.
\end{methode}

\begin{exemplebox}[Area with a sign change]
Find the total area enclosed between the curve $y = x^2 - 2x$ and the $x$-axis from $x = 0$ to $x = 3$.
\tcblower
\textbf{Step 1: zeros of $f$.} $x^2 - 2x = x(x-2) = 0 \Rightarrow x = 0$ or $x = 2$. On $[0,2]$, $f(x)\le 0$ (e.g.\ $f(1)=-1$); on $[2,3]$, $f(x)\ge 0$ (e.g.\ $f(3)=3$).

\textbf{Step 2: split the integral.} An antiderivative is $F(x) = \frac{x^3}{3} - x^2$.
\[
\int_0^2 (x^2-2x)\,dx = \left[\frac{x^3}{3}-x^2\right]_0^2 = \frac{8}{3}-4 = -\frac{4}{3},
\]
\[
\int_2^3 (x^2-2x)\,dx = \left[\frac{x^3}{3}-x^2\right]_2^3 = (9-9)-\left(\frac{8}{3}-4\right) = \frac{4}{3}.
\]

\textbf{Step 3: add absolute values.}
\[
A = \left|-\frac{4}{3}\right| + \frac{4}{3} = \frac{8}{3}\ \text{square units}.
\]
\begin{attention}
Had we computed $\int_0^3 (x^2-2x)\,dx = 0$, we would wrongly conclude the area is zero: the two parts cancel. Splitting at the zero $x=2$ is essential.
\end{attention}
\end{exemplebox}

\begin{exemplebox}[Area between two curves]
Find the area of the region bounded by $y = x^2$ and $y = 2x - x^2$.
\tcblower
\textbf{Step 1: intersection points.} $x^2 = 2x - x^2 \Rightarrow 2x^2 - 2x = 0 \Rightarrow 2x(x-1)=0 \Rightarrow x=0, 1$.

\textbf{Step 2: determine upper curve.} At $x=0.5$, $x^2 = 0.25$, $2x-x^2 = 0.75$. So $y = 2x-x^2$ is above $y=x^2$ on $[0,1]$.

\textbf{Step 3: integrate.}
\[
A = \int_0^1 \big[(2x-x^2) - x^2\big]\,dx = \int_0^1 (2x - 2x^2)\,dx = \left[x^2 - \frac{2}{3}x^3\right]_0^1 = 1 - \frac{2}{3} = \frac{1}{3}.
\]
\end{exemplebox}

\begin{popfigure}
\begin{tikzpicture}[scale=1.2]
\draw[->] (-0.5,0) -- (2,0) node[right] {$x$};
\draw[->] (0,-0.5) -- (0,1.5) node[above] {$y$};
\draw[domain=0:1.5, smooth, variable=\x, popblue, thick] plot ({\x}, {\x*\x}) node[right] {$y=x^2$};
\draw[domain=0:1.5, smooth, variable=\x, poporange, thick] plot ({\x}, {2*\x - \x*\x}) node[right] {$y=2x-x^2$};
\fill[popgreen!30] (0,0) -- plot[domain=0:1, smooth, variable=\x] ({\x}, {2*\x - \x*\x}) -- plot[domain=1:0, smooth, variable=\x] ({\x}, {\x*\x}) -- cycle;
\foreach \x in {0,1} \draw (\x, -0.05) -- (\x, 0.05) node[below] {$\x$};
\draw (0.5, 0.3) node[popgreen] {Area};
\end{tikzpicture}
\legende{Area between $y=x^2$ and $y=2x-x^2$}
\end{popfigure}

\courssub{Area Bounded by a Curve and the $y$-axis}

\begin{definition}[Area between a Curve and the $y$-axis]
Let $x = g(y)$ be continuous on $[c,d]$. The area of the region bounded by $x=g(y)$, the $y$-axis, and the lines $y=c$, $y=d$ is
\[
A = \int_c^d |g(y)|\,dy.
\]
If the curve is given as $y = f(x)$, solve for $x = g(y)$ (choosing the appropriate branch) and integrate with respect to $y$.
\end{definition}

\begin{methode}[Area between $x=g(y)$ and the $y$-axis]
\begin{enumerate}
\item Express $x$ as a function of $y$: solve $y = f(x)$ for $x = g(y)$.
\item Identify the limits as $y$-values, $y = c$ and $y = d$.
\item The area is $A = \displaystyle\int_c^d g(y)\,dy$ if the curve lies to the right of the $y$-axis ($g(y)\ge 0$); take $|g(y)|$ otherwise.
\item Alternative (rectangle method): compute the area of a bounding rectangle and subtract the area under $y=f(x)$.
\end{enumerate}
\textbf{Reflex:} integrate with respect to $y$; the integrand must be a function of $y$ alone.
\end{methode}

\begin{exemplebox}[Area bounded by the $y$-axis]
Find the area of the region bounded by the curve $y = x^2$, the $y$-axis and the lines $y = 1$ and $y = 4$.
\tcblower
\textbf{Step 1: solve for $x$.} $y = x^2 \Rightarrow x = \sqrt{y}$ (taking the branch $x\ge 0$).

\textbf{Step 2: integrate with respect to $y$ from $1$ to $4$.}
\[
A = \int_1^4 \sqrt{y}\,dy = \int_1^4 y^{1/2}\,dy = \left[\frac{2}{3}y^{3/2}\right]_1^4 = \frac{2}{3}\left(4^{3/2}-1^{3/2}\right) = \frac{2}{3}(8-1) = \frac{14}{3}.
\]

\textbf{Check by subtraction:} the rectangle $0\le x\le 2$, $1\le y\le 4$ has area $2\times 3 = 6$. The area under $y=x^2$ from $0$ to $2$ is $\int_0^2 x^2 dx = \frac{8}{3}$, and the rectangle $0\le x\le 2,\ 0\le y\le 1$ has area $2$. The desired area is the area of the large rectangle minus the area under the curve from $0$ to $2$ plus the area of the small rectangle? Wait, better: the region is the part of the rectangle $0\le x\le 2$, $1\le y\le 4$ that lies to the right of the curve. That equals area of rectangle minus area under curve from $x=0$ to $x=2$ for $y$ from $1$ to $4$? Actually, the region is bounded by $x=0$, $y=1$, $y=4$, and $x=\sqrt{y}$. The rectangle method: consider the rectangle $0\le x\le 2$, $0\le y\le 4$ (area $8$). Subtract the area under $y=x^2$ from $0$ to $2$ ($\frac{8}{3}$) and subtract the rectangle $0\le x\le 1$, $0\le y\le 1$ (area $1$)? No, that's not the region. Let's stick to the direct integration which is correct. The check: the region is exactly $\int_1^4 \sqrt{y}\,dy = 14/3 \approx 4.667$. The rectangle $0\le x\le 2$, $1\le y\le 4$ has area $6$. The area to the left of the curve in that rectangle is $\int_1^4 \sqrt{y}\,dy$? No, the curve is the right boundary. The area is exactly the integral. The alternative check: area of rectangle $0\le x\le 2$, $1\le y\le 4$ is $6$. The area between the curve and the line $x=2$ from $y=1$ to $4$ is $\int_1^4 (2-\sqrt{y})\,dy = [2y - \frac{2}{3}y^{3/2}]_1^4 = (8 - \frac{16}{3}) - (2 - \frac{2}{3}) = \frac{8}{3} - \frac{4}{3} = \frac{4}{3}$. Then the desired area = rectangle - that area = $6 - \frac{4}{3} = \frac{14}{3}$. \checkmark

Hence $A = \dfrac{14}{3}$ square units.
\end{exemplebox}

\begin{exemplebox}[Area with curve crossing the $y$-axis]
Find the area of the region bounded by $y = x^3$, the $y$-axis, and the lines $y = -1$ and $y = 1$.
\tcblower
\textbf{Step 1: solve for $x$.} $x = \sqrt[3]{y} = y^{1/3}$. This is defined for all $y$ and is negative for $y<0$, positive for $y>0$.

\textbf{Step 2: split at $y=0$.}
\[
A = \int_{-1}^0 |y^{1/3}|\,dy + \int_0^1 |y^{1/3}|\,dy = \int_{-1}^0 (-y^{1/3})\,dy + \int_0^1 y^{1/3}\,dy.
\]
\[
= -\left[\frac{3}{4}y^{4/3}\right]_{-1}^0 + \left[\frac{3}{4}y^{4/3}\right]_0^1 = -\left(0 - \frac{3}{4}(1)\right) + \frac{3}{4} = \frac{3}{4} + \frac{3}{4} = \frac{3}{2}.
\]
\end{exemplebox}

\courssub{Volumes of Solids of Revolution}

\begin{definition}[Solid of Revolution]
When a plane region is rotated through $360^\circ$ about a line (the \cle{axis of revolution}) in its plane, it generates a \cle{solid of revolution}. The volume can be computed by the \cle{disc/washer method} (slicing perpendicular to the axis) or the \cle{shell method} (slicing parallel to the axis). The syllabus requires the disc/washer method.
\end{definition}

\begin{propriete}[Disc Method about the $x$-axis]
If the region bounded by $y=f(x)\ge 0$, the $x$-axis, $x=a$, $x=b$ is rotated about the $x$-axis, the volume is
\[
V = \pi \int_a^b [f(x)]^2\,dx.
\]
For a region between two curves $y=f(x)$ (upper) and $y=g(x)$ (lower, $g(x)\ge 0$) rotated about the $x$-axis (washer method):
\[
V = \pi \int_a^b \left([f(x)]^2 - [g(x)]^2\right)\,dx.
\]
\end{propriete}

\begin{methode}[Disc method about the $x$-axis]
\begin{enumerate}
\item Sketch the region and identify the axis of rotation (the $x$-axis).
\item The cross-section perpendicular to the $x$-axis is a \cle{disc} of radius $r = f(x)$, area $\pi[f(x)]^2$.
\item Apply $V = \pi\displaystyle\int_a^b [f(x)]^2\,dx$.
\item For a region between two curves rotated about the $x$-axis (washer method): $V = \pi\displaystyle\int_a^b \left([f(x)]^2 - [g(x)]^2\right)dx$.
\end{enumerate}
\textbf{Reflex:} it is $y^2$ that is integrated, not $y$. Do not forget the factor $\pi$.
\end{methode}

\begin{exemplebox}[Volume about the $x$-axis]
The region bounded by the curve $y = \sqrt{x}$, the $x$-axis and the lines $x = 1$ and $x = 4$ is rotated through $360^{\circ}$ about the $x$-axis. Find the volume of the solid generated.
\tcblower
\textbf{Step 1: set up the disc formula.} Here $f(x) = \sqrt{x}$, so $[f(x)]^2 = x$, and $a=1$, $b=4$.

\textbf{Step 2: integrate.}
\[
V = \pi\int_1^4 (\sqrt{x})^2\,dx = \pi\int_1^4 x\,dx = \pi\left[\frac{x^2}{2}\right]_1^4 = \pi\left(\frac{16}{2}-\frac{1}{2}\right) = \frac{15\pi}{2}.
\]

\textbf{Conclusion:} $V = \dfrac{15\pi}{2}$ cubic units $\approx 23.56$ cubic units.

\textbf{Remark:} the solid is a paraboloid segment; the squaring of $\sqrt{x}$ turned the integral into a trivial polynomial one — always simplify $[f(x)]^2$ before integrating.
\end{exemplebox}

\begin{exemplebox}[Washer method about the $x$-axis]
The region bounded by $y = x^2$ and $y = x$ is rotated about the $x$-axis. Find the volume.
\tcblower
\textbf{Step 1: intersection points.} $x^2 = x \Rightarrow x(x-1)=0 \Rightarrow x=0, 1$. On $[0,1]$, $y=x$ is above $y=x^2$.

\textbf{Step 2: washer formula.}
\[
V = \pi\int_0^1 \left( x^2 - (x^2)^2 \right)\,dx = \pi\int_0^1 (x^2 - x^4)\,dx = \pi\left[\frac{x^3}{3} - \frac{x^5}{5}\right]_0^1 = \pi\left(\frac{1}{3} - \frac{1}{5}\right) = \frac{2\pi}{15}.
\]
\end{exemplebox}

\begin{popfigure}
\begin{tikzpicture}[scale=1]
\draw[->] (-0.5,0) -- (2,0) node[right] {$x$};
\draw[->] (0,-0.5) -- (0,1.5) node[above] {$y$};
\draw[domain=0:1.2, smooth, variable=\x, popblue, thick] plot ({\x}, {\x}) node[right] {$y=x$};
\draw[domain=0:1.2, smooth, variable=\x, poporange, thick] plot ({\x}, {\x*\x}) node[right] {$y=x^2$};
\fill[popblue!20] (0,0) -- plot[domain=0:1, smooth, variable=\x] ({\x}, {\x}) -- plot[domain=1:0, smooth, variable=\x] ({\x}, {\x*\x}) -- cycle;
\draw[poporange, dashed] (1,0) -- (1,1);
\draw[poporange, dashed] (0.5,0) -- (0.5,0.5);
\draw[poporange, dashed] (0.5,0.5) -- (0.5,0.25);
\node at (0.5, 0.7) {Region};
\draw[<->, popgreen] (0.5, 0.25) -- (0.5, 0.5) node[midway, right] {$R_{\text{outer}} - R_{\text{inner}}$};
\end{tikzpicture}
\legende{Washer cross-section for rotation about $x$-axis}
\end{popfigure}

\begin{propriete}[Disc Method about the $y$-axis]
If the region bounded by $x=g(y)\ge 0$, the $y$-axis, $y=c$, $y=d$ is rotated about the $y$-axis, the volume is
\[
V = \pi \int_c^d [g(y)]^2\,dy = \pi \int_c^d x^2\,dy.
\]
For a region between two curves $x=f(y)$ (right) and $x=g(y)$ (left, $g(y)\ge 0$) rotated about the $y$-axis:
\[
V = \pi \int_c^d \left([f(y)]^2 - [g(y)]^2\right)\,dy.
\]
\end{propriete}

\begin{methode}[Disc method about the $y$-axis]
\begin{enumerate}
\item Rewrite the curve as $x = g(y)$ and determine the $y$-limits $c$ and $d$.
\item The disc radius is now $x = g(y)$, so $V = \pi\displaystyle\int_c^d [g(y)]^2\,dy = \pi\int_c^d x^2\,dy$.
\item Alternatively, when rotating about the $y$-axis a region under $y=f(x)$, use the rectangle-minus-volume technique or the shell idea: $V = \pi b^2 d - \pi a^2 c - \pi\int_c^d x^2\,dy$ as appropriate to the sketch.
\end{enumerate}
\textbf{Reflex:} about the $y$-axis $\Rightarrow$ integrate $x^2$ with respect to $y$.
\end{methode}

\begin{exemplebox}[Volume about the $y$-axis]
The region bounded by the curve $y = x^2$, the $y$-axis and the line $y = 4$ (with $x\ge 0$) is rotated about the $y$-axis. Find the volume generated.
\tcblower
\textbf{Step 1: express $x$ in terms of $y$.} $x = \sqrt{y}$, so $x^2 = y$. The limits are $y = 0$ (at $x=0$) and $y = 4$.

\textbf{Step 2: apply the formula.}
\[
V = \pi\int_0^4 x^2\,dy = \pi\int_0^4 y\,dy = \pi\left[\frac{y^2}{2}\right]_0^4 = \pi\cdot\frac{16}{2} = 8\pi.
\]

\textbf{Check (complement method):} rotating the rectangle $0\le x\le 2$, $0\le y\le 4$ about the $y$-axis gives a cylinder of volume $\pi r^2 h = \pi(2^2)(4) = 16\pi$. The region under $y=x^2$ from $0$ to $2$ rotated about the $y$-axis leaves the volume we want as the complement within the cylinder: indeed the paraboloid inside the cylinder occupies half of it, $8\pi$, consistent with the known result that a paraboloid of revolution fills half its circumscribing cylinder. \checkmark

Hence $V = 8\pi$ cubic units.
\end{exemplebox}

\begin{exemplebox}[Washer method about the $y$-axis]
The region bounded by $y = x^2$, $y = 4$, and the $y$-axis ($x\ge 0$) is rotated about the $y$-axis. Find the volume. (Same as above, but let's do it as a washer: outer radius $x=2$ from $y=0$ to $4$, inner radius $x=\sqrt{y}$.)
\tcblower
Actually, the region is bounded by $x=0$, $y=4$, and $y=x^2$. Rotating about $y$-axis: outer boundary is $x=2$ (from $y=0$ to $4$)? No, the region is $0\le x\le \sqrt{y}$, $0\le y\le 4$. The outer radius is $\sqrt{y}$, inner is $0$. So it's a disc, not a washer. For a washer example: region between $y=x^2$ and $y=2x$ rotated about $y$-axis. Intersections: $x^2=2x \Rightarrow x=0,2$. In terms of $y$: $x=\sqrt{y}$ and $x=y/2$. For $y\in[0,4]$, $y/2 \le \sqrt{y}$? At $y=1$, $0.5 \le 1$, yes. So outer radius is $\sqrt{y}$, inner is $y/2$.
\[
V = \pi\int_0^4 \left( (\sqrt{y})^2 - (y/2)^2 \right)\,dy = \pi\int_0^4 \left( y - \frac{y^2}{4} \right)\,dy = \pi\left[ \frac{y^2}{2} - \frac{y^3}{12} \right]_0^4 = \pi\left( 8 - \frac{64}{12} \right) = \pi\left( 8 - \frac{16}{3} \right) = \frac{8\pi}{3}.
\]
\end{exemplebox}

\begin{popfigure}
\begin{tikzpicture}[scale=0.8]
\draw[->] (-0.5,0) -- (3,0) node[right] {$x$};
\draw[->] (0,-0.5) -- (0,4.5) node[above] {$y$};
\draw[domain=0:2, smooth, variable=\x, popblue, thick] plot ({\x}, {\x*\x}) node[right] {$y=x^2$};
\draw[poporange, thick] (0,0) -- (0,4);
\draw[poporange, thick] (0,4) -- (2,4);
\fill[popgreen!30] (0,0) -- plot[domain=0:2, smooth, variable=\x] ({\x}, {\x*\x}) -- (2,4) -- (0,4) -- cycle;
\draw[dashed] (1,0) node[below] {$1$} -- (1,1);
\draw[dashed] (1,1) -- (0,1) node[left] {$1$};
\node at (1, 2) {Region};
\draw[<->, popgreen] (0, 2) -- (1.414, 2) node[midway, above] {$x=\sqrt{y}$};
\end{tikzpicture}
\legende{Region rotated about $y$-axis to form a paraboloid}
\end{popfigure}

\courssub{Centroids and Centres of Mass by Integration}

\begin{definition}[Centroid of a Plane Lamina]
The \cle{centroid} $(\bar{x}, \bar{y})$ of a uniform plane lamina (thin plate of constant density) is the point where the lamina would balance. For a region bounded by $y=f(x)\ge 0$, the $x$-axis, $x=a$, $x=b$, with area $A = \int_a^b f(x)\,dx$:
\[
\bar{x} = \frac{1}{A} \int_a^b x\,f(x)\,dx, \qquad
\bar{y} = \frac{1}{A} \cdot \frac{1}{2} \int_a^b [f(x)]^2\,dx.
\]
The factor $\frac{1}{2}$ in $\bar{y}$ arises because the centroid of a thin vertical strip of height $f(x)$ is at its midpoint, $y = f(x)/2$.
\end{definition}

\begin{propriete}[Centroid Formulas]
\begin{itemize}
\item \textbf{Area:} $A = \int_a^b f(x)\,dx$ (or $\int_c^d g(y)\,dy$ for $x=g(y)$).
\item \textbf{Moment about $y$-axis:} $M_y = \int_a^b x f(x)\,dx \quad\Rightarrow\quad \bar{x} = M_y/A$.
\item \textbf{Moment about $x$-axis:} $M_x = \frac{1}{2}\int_a^b [f(x)]^2\,dx \quad\Rightarrow\quad \bar{y} = M_x/A$.
\item \textbf{For a uniform rod/wire along $x$-axis with density $\rho(x)$:} $\bar{x} = \dfrac{\int_a^b x\rho(x)\,dx}{\int_a^b \rho(x)\,dx}$.
\item \textbf{For a solid of revolution about $x$-axis:} the centroid lies on the axis, $\bar{y}=0$, and $\bar{x} = \dfrac{\int_a^b x [f(x)]^2\,dx}{\int_a^b [f(x)]^2\,dx}$.
\end{itemize}
\end{propriete}

\begin{methode}[Centroid of a plane region under $y=f(x)$]
\begin{enumerate}
\item Find the limits $a,b$ (intersections with $x$-axis or given lines) and compute the area $A = \int_a^b f(x)\,dx$.
\item Compute the moment about the $y$-axis: $M_y = \displaystyle\int_a^b x\,f(x)\,dx$, then $\bar{x} = \dfrac{M_y}{A}$.
\item Compute the moment about the $x$-axis: $M_x = \dfrac{1}{2}\displaystyle\int_a^b [f(x)]^2\,dx$, then $\bar{y} = \dfrac{M_x}{A}$.
\item Use symmetry whenever possible: if the region is symmetric about the $y$-axis, $\bar{x}=0$; if symmetric about the $x$-axis, $\bar{y}=0$.
\end{enumerate}
\textbf{Reflex:} centroid = (first moment) $\div$ (total area or mass). The factor $\tfrac12$ in $M_x$ comes from the centroid of each vertical strip being at height $f(x)/2$.
\end{methode}

\begin{exemplebox}[Centroid of a region under a parabola]
Find the coordinates of the centroid of the region bounded by $y = 4 - x^2$ and the $x$-axis.
\tcblower
\textbf{Step 1: limits and area.} $4 - x^2 = 0 \Rightarrow x = \pm 2$.
\[
A = \int_{-2}^{2}(4-x^2)\,dx = \left[4x - \frac{x^3}{3}\right]_{-2}^{2} = \left(8-\frac{8}{3}\right) - \left(-8+\frac{8}{3}\right) = \frac{32}{3}.
\]

\textbf{Step 2: $\bar{x}$.} By symmetry of the region about the $y$-axis, $\bar{x} = 0$. (Check: $\int_{-2}^{2} x(4-x^2)\,dx = 0$ since the integrand is odd.)

\textbf{Step 3: $\bar{y}$.}
\[
M_x = \frac{1}{2}\int_{-2}^{2}(4-x^2)^2\,dx = \frac{1}{2}\int_{-2}^{2}\left(16 - 8x^2 + x^4\right)dx = \frac{1}{2}\left[16x - \frac{8x^3}{3} + \frac{x^5}{5}\right]_{-2}^{2}.
\]
Evaluating at $2$: $16(2) - \frac{8(8)}{3} + \frac{32}{5} = 32 - \frac{64}{3} + \frac{32}{5} = \frac{480 - 320 + 96}{15} = \frac{256}{15}$.
At $-2$: $-32 + \frac{64}{3} - \frac{32}{5} = -\frac{256}{15}$.
Difference: $\frac{256}{15} - (-\frac{256}{15}) = \frac{512}{15}$. Then $M_x = \frac{1}{2} \cdot \frac{512}{15} = \frac{256}{15}$.
\[
\bar{y} = \frac{M_x}{A} = \frac{256/15}{32/3} = \frac{256}{15}\times\frac{3}{32} = \frac{8}{5}.
\]

\textbf{Conclusion:} the centroid is $\left(0,\ \dfrac{8}{5}\right)$. Note $\bar{y} = 1.6 < 2 = $ half the maximum height, which is sensible since the lamina is wider near the base.
\end{exemplebox}

\begin{exemplebox}[Centroid of a region bounded by two curves]
Find the centroid of the region bounded by $y = x^2$ and $y = 2x$.
\tcblower
\textbf{Step 1: intersections.} $x^2 = 2x \Rightarrow x=0, 2$. Upper curve: $y=2x$, lower: $y=x^2$.

\textbf{Step 2: area.}
\[
A = \int_0^2 (2x - x^2)\,dx = \left[x^2 - \frac{x^3}{3}\right]_0^2 = 4 - \frac{8}{3} = \frac{4}{3}.
\]

\textbf{Step 3: $\bar{x}$.}
\[
M_y = \int_0^2 x(2x - x^2)\,dx = \int_0^2 (2x^2 - x^3)\,dx = \left[\frac{2x^3}{3} - \frac{x^4}{4}\right]_0^2 = \frac{16}{3} - 4 = \frac{4}{3}.
\]
\[
\bar{x} = \frac{M_y}{A} = \frac{4/3}{4/3} = 1.
\]

\textbf{Step 4: $\bar{y}$.}
\[
M_x = \frac{1}{2}\int_0^2 \left( (2x)^2 - (x^2)^2 \right)\,dx = \frac{1}{2}\int_0^2 (4x^2 - x^4)\,dx = \frac{1}{2}\left[ \frac{4x^3}{3} - \frac{x^5}{5} \right]_0^2 = \frac{1}{2}\left( \frac{32}{3} - \frac{32}{5} \right) = \frac{1}{2}\cdot\frac{160-96}{15} = \frac{32}{15}.
\]
\[
\bar{y} = \frac{M_x}{A} = \frac{32/15}{4/3} = \frac{32}{15}\cdot\frac{3}{4} = \frac{8}{5} = 1.6.
\]
\textbf{Centroid:} $(1, 1.6)$.
\end{exemplebox}

\begin{popfigure}
\begin{tikzpicture}[scale=1]
\draw[->] (-0.5,0) -- (3,0) node[right] {$x$};
\draw[->] (0,-0.5) -- (0,4.5) node[above] {$y$};
\draw[domain=0:2.2, smooth, variable=\x, popblue, thick] plot ({\x}, {2*\x}) node[right] {$y=2x$};
\draw[domain=0:2.2, smooth, variable=\x, poporange, thick] plot ({\x}, {\x*\x}) node[right] {$y=x^2$};
\fill[popgreen!30] (0,0) -- plot[domain=0:2, smooth, variable=\x] ({\x}, {2*\x}) -- plot[domain=2:0, smooth, variable=\x] ({\x}, {\x*\x}) -- cycle;
\draw[popgreen, thick] (1, 1.6) circle (0.08) node[above right] {$(\bar{x},\bar{y})=(1,1.6)$};
\foreach \x in {0,1,2} \draw (\x, -0.1) -- (\x, 0.1) node[below] {$\x$};
\foreach \y in {1,2,3,4} \draw (-0.1, \y) -- (0.1, \y) node[left] {$\y$};
\end{tikzpicture}
\legende{Centroid of region between $y=x^2$ and $y=2x$}
\end{popfigure}

\courssub{Numerical Integration — the Trapezium Rule}

\begin{definition}[Trapezium Rule]
The \cle{trapezium rule} (or trapezoidal rule) approximates $\int_a^b f(x)\,dx$ by dividing $[a,b]$ into $n$ strips of equal width $h = \frac{b-a}{n}$ and approximating the area under the curve by trapezia. If $y_i = f(x_i)$ with $x_i = a + ih$, then
\[
\int_a^b f(x)\,dx \approx \frac{h}{2}\left[ y_0 + y_n + 2\sum_{i=1}^{n-1} y_i \right].
\]
The error is proportional to $h^2$ (or $1/n^2$) and depends on $f''(\xi)$ for some $\xi\in(a,b)$.
\end{definition}

\begin{propriete}[Error Estimate]
If $f''$ is continuous on $[a,b]$, the error $E$ in the trapezium rule satisfies
\[
|E| \le \frac{(b-a)^3}{12n^2} \max_{x\in[a,b]} |f''(x)|.
\]
Doubling $n$ roughly quarters the error.
\end{propriete}

\begin{methode}[Applying the trapezium rule]
\begin{enumerate}
\item Divide $[a,b]$ into $n$ strips of equal width $h = \dfrac{b-a}{n}$; the ordinates are $y_0, y_1, \dots, y_n$ with $y_i = f(a+ih)$.
\item Tabulate the values of $f$ at $x_0, x_1, \dots, x_n$ (work to at least 4 decimal places).
\item Apply $\displaystyle\int_a^b f(x)\,dx \approx \frac{h}{2}\left[y_0 + y_n + 2\left(y_1 + y_2 + \dots + y_{n-1}\right)\right]$.
\item More strips $\Rightarrow$ better accuracy. The rule \emph{over}estimates for a concave-up curve ($f''>0$, chords lie above the curve) and \emph{under}estimates for a concave-down curve ($f''<0$).
\end{enumerate}
\textbf{Reflex:} first and last ordinates counted once; all interior ordinates counted twice.
\end{methode}

\begin{exemplebox}[Trapezium rule with 4 strips]
Use the trapezium rule with 4 strips to estimate $\displaystyle\int_0^2 \frac{1}{1+x^2}\,dx$, giving your answer to 3 decimal places. Compare with the exact value $\arctan 2 \approx 1.1071$.
\tcblower
\textbf{Step 1: strip width.} $h = \dfrac{2-0}{4} = 0.5$.

\textbf{Step 2: table of ordinates.}
\begin{center}
\begin{tabular}{|c|c|c|c|c|c|}
\hline
\rowcolor{popblueL} $x$ & $0$ & $0.5$ & $1.0$ & $1.5$ & $2.0$ \\
\hline
$y = \frac{1}{1+x^2}$ & $1.0000$ & $0.8000$ & $0.5000$ & $0.3077$ & $0.2000$ \\
\hline
\end{tabular}
\end{center}
(For instance $y_3 = \frac{1}{1+2.25} = \frac{1}{3.25} = 0.3077$.)

\textbf{Step 3: apply the formula.}
\begin{align*}
\int_0^2 \frac{dx}{1+x^2} &\approx \frac{0.5}{2}\Big[y_0 + y_4 + 2(y_1 + y_2 + y_3)\Big] \\
&= 0.25\Big[1.0000 + 0.2000 + 2(0.8000 + 0.5000 + 0.3077)\Big] \\
&= 0.25\big[1.2000 + 2(1.6077)\big] = 0.25 \times 4.4154 = 1.10385.
\end{align*}

\textbf{Conclusion:} $\displaystyle\int_0^2 \frac{dx}{1+x^2} \approx 1.104$ (3 d.p.). The exact value is $\arctan 2 \approx 1.1071$; the error is about $0.003$, i.e.\ roughly $0.3\%$ — an excellent estimate with only 4 strips. Since $f''(x) = \frac{6x^2-2}{(1+x^2)^3}$, on $[0,2]$ it changes sign; near $x=0$ it's negative (concave down) so the rule underestimates there, but overall the estimate is slightly below the true value. Increasing $n$ to 8 would roughly quarter the error.
\end{exemplebox}

\begin{exemplebox}[Trapezium rule with a table of values]
The table below gives values of a function $f(x)$. Use the trapezium rule with 6 strips to estimate $\int_0^3 f(x)\,dx$.

\begin{center}
\begin{tabular}{|c|c|c|c|c|c|c|c|}
\hline
\rowcolor{popblueL} $x$ & $0$ & $0.5$ & $1.0$ & $1.5$ & $2.0$ & $2.5$ & $3.0$ \\
\hline
$f(x)$ & $2.0$ & $2.5$ & $3.2$ & $4.1$ & $5.3$ & $6.8$ & $8.5$ \\
\hline
\end{tabular}
\end{center}
\tcblower
$h = \frac{3-0}{6} = 0.5$.
\[
\int_0^3 f(x)\,dx \approx \frac{0.5}{2}\left[ 2.0 + 8.5 + 2(2.5 + 3.2 + 4.1 + 5.3 + 6.8) \right] = 0.25\left[ 10.5 + 2(21.9) \right] = 0.25 \times 54.3 = 13.575.
\]
\end{exemplebox}

\begin{popfigure}
\begin{tikzpicture}[scale=1]
\draw[->] (-0.5,0) -- (4,0) node[right] {$x$};
\draw[->] (0,-0.5) -- (0,3) node[above] {$y$};
\draw[domain=0:3.5, smooth, variable=\x, popblue, thick] plot ({\x}, {0.2*\x*\x + 0.5});
\foreach \i in {0,...,6} {
  \pgfmathsetmacro{\x}{0.5*\i}
  \pgfmathsetmacro{\y}{0.2*\x*\x + 0.5}
  \draw[poporange, fill=poporange!20] (\x, 0) -- (\x, \y);
  \ifnum\i<6
    \pgfmathsetmacro{\xnext}{0.5*\i+0.5}
    \pgfmathsetmacro{\ynext}{0.2*\xnext*\xnext + 0.5}
    \draw[poporange, thick] (\x, \y) -- (\xnext, \ynext);
  \fi
}
\foreach \x in {0,0.5,1,1.5,2,2.5,3} \draw (\x, -0.1) -- (\x, 0.1) node[below] {$\x$};
\node at (2, 2.5) {Trapezium rule approximation};
\end{tikzpicture}
\legende{Trapezium rule with $n=6$ strips}
\end{popfigure}

\courssub{Mean Value Theorem for Integrals and Applications}

\begin{definition}[Mean Value of a Function]
The \cle{mean value} (or average value) of a continuous function $f$ on $[a,b]$ is
\[
\bar{f} = \frac{1}{b-a} \int_a^b f(x)\,dx.
\]
Geometrically, $\bar{f}$ is the height of the rectangle on base $[a,b]$ having the same area as the region under the curve $y=f(x)$.
\end{definition}

\begin{propriete}[Mean Value Theorem for Integrals]
If $f$ is continuous on $[a,b]$, then there exists at least one number $c \in (a,b)$ such that
\[
f(c) = \bar{f} = \frac{1}{b-a} \int_a^b f(x)\,dx,
\]
or equivalently,
\[
\int_a^b f(x)\,dx = f(c)(b-a).
\]
\end{propriete}

\begin{methode}[Finding the mean value and the point $c$]
\begin{enumerate}
\item Compute the definite integral $I = \int_a^b f(x)\,dx$.
\item The mean value is $\bar{f} = \dfrac{I}{b-a}$.
\item To find $c$ guaranteed by the Mean Value Theorem: solve $f(c) = \bar{f}$ and select the solution(s) lying in the open interval $(a,b)$.
\end{enumerate}
\textbf{Reflex:} mean value $=$ (area under curve) $\div$ (length of interval).
\end{methode}

\begin{exemplebox}[Mean value and the point $c$]
Find the mean value of $f(x) = x^2 + 1$ on $[0,3]$, and determine the value(s) of $c$ guaranteed by the Mean Value Theorem for integrals.
\tcblower
\textbf{Step 1: compute the integral.}
\[
\int_0^3 (x^2+1)\,dx = \left[\frac{x^3}{3} + x\right]_0^3 = 9 + 3 = 12.
\]

\textbf{Step 2: mean value.}
\[
\bar{f} = \frac{1}{3-0}\times 12 = 4.
\]
So the rectangle of base $3$ and height $4$ has the same area as the region under the parabola.

\textbf{Step 3: find $c$.} Solve $f(c) = \bar f$:
\[
c^2 + 1 = 4 \Rightarrow c^2 = 3 \Rightarrow c = \pm\sqrt{3}.
\]
Only $c = \sqrt{3} \approx 1.732$ lies in $(0,3)$.

\textbf{Conclusion:} the mean value is $4$, attained at $c = \sqrt{3}$, confirming the Mean Value Theorem: $\int_0^3 (x^2+1)\,dx = f(\sqrt{3})\times(3-0) = 4\times 3 = 12$. \checkmark
\end{exemplebox}

\begin{exemplebox}[Application: average velocity]
A particle moves along a line with velocity $v(t) = 3t^2 - 12t + 9$ m/s for $0 \le t \le 4$. Find the average velocity over this interval.
\tcblower
\[
\bar{v} = \frac{1}{4-0} \int_0^4 (3t^2 - 12t + 9)\,dt = \frac{1}{4} \left[ t^3 - 6t^2 + 9t \right]_0^4 = \frac{1}{4} \left( 64 - 96 + 36 \right) = \frac{4}{4} = 1\ \text{m/s}.
\]
The Mean Value Theorem guarantees a time $c \in (0,4)$ where $v(c) = 1$. Solve $3c^2 - 12c + 9 = 1 \Rightarrow 3c^2 - 12c + 8 = 0 \Rightarrow c = \frac{12 \pm \sqrt{144-96}}{6} = \frac{12 \pm \sqrt{48}}{6} = 2 \pm \frac{2\sqrt{3}}{3}$. Both $2 - \frac{2\sqrt{3}}{3} \approx 0.845$ and $2 + \frac{2\sqrt{3}}{3} \approx 3.155$ lie in $(0,4)$.
\end{exemplebox}

\begin{popfigure}
\begin{tikzpicture}[scale=1]
\draw[->] (-0.5,0) -- (4,0) node[right] {$x$};
\draw[->] (0,-0.5) -- (0,3) node[above] {$y$};
\draw[domain=0:3.5, smooth, variable=\x, popblue, thick] plot ({\x}, {0.3*\x*\x - 0.5*\x + 1});
\draw[poporange, dashed] (0, 1.5) -- (3.5, 1.5) node[right] {$\bar{f}=1.5$};
\fill[popgreen!20] (0,0) -- (0,1.5) -- (3.5,1.5) -- (3.5,0) -- cycle;
\draw[popgreen, thick] (1.5, 0) -- (1.5, 1.5) node[midway, right] {$\bar{f}$};
\node at (1.75, 2.5) {Rectangle area = Area under curve};
\end{tikzpicture}
\legende{Geometric interpretation of the mean value}
\end{popfigure}

\begin{aretenir}[Key Formulas for Integration Applications]
\begin{itemize}
\item \textbf{Riemann sum:} $\int_a^b f(x)\,dx = \lim_{n\to\infty} \sum_{i=1}^n f(a+i\Delta x)\Delta x$, $\Delta x = \frac{b-a}{n}$.
\item \textbf{Area ($x$-axis):} $A = \int_a^b |f(x)|\,dx$; split at zeros of $f$.
\item \textbf{Area ($y$-axis):} $A = \int_c^d |g(y)|\,dy$ where $x=g(y)$.
\item \textbf{Volume ($x$-axis):} $V = \pi\int_a^b [f(x)]^2\,dx$ (disc); $V = \pi\int_a^b ([f]^2-[g]^2)\,dx$ (washer).
\item \textbf{Volume ($y$-axis):} $V = \pi\int_c^d [g(y)]^2\,dy$ (disc); $V = \pi\int_c^d ([f]^2-[g]^2)\,dy$ (washer).
\item \textbf{Centroid:} $\bar{x} = \frac{1}{A}\int x f\,dx$, $\bar{y} = \frac{1}{2A}\int f^2\,dx$.
\item \textbf{Trapezium rule:} $\int_a^b f \approx \frac{h}{2}[y_0+y_n+2\sum_{i=1}^{n-1}y_i]$, $h=\frac{b-a}{n}$.
\item \textbf{Mean value:} $\bar{f} = \frac{1}{b-a}\int_a^b f\,dx$; $\exists c\in(a,b): f(c)=\bar{f}$.
\end{itemize}
\end{aretenir}

\begin{attention}[Frequent Examination Traps]
\begin{itemize}
\item Giving a \emph{negative} area: always integrate $|f|$ or split at the zeros of $f$.
\item Squaring incorrectly in volume problems: $V = \pi\int y^2\,dx$, never $\pi\left(\int y\,dx\right)^2$.
\item Mixing the axes: about the $x$-axis use $y^2\,dx$; about the $y$-axis use $x^2\,dy$.
\item In the trapezium rule, forgetting to double the interior ordinates or using the wrong strip width $h$.
\item In centroid questions, omitting the factor $\tfrac12$ in $M_x = \tfrac12\int y^2\,dx$.
\item Forgetting to check that the solution $c$ from the Mean Value Theorem lies in the \emph{open} interval $(a,b)$.
\item Confusing the mean value of a function with the mean value theorem for derivatives.
\end{itemize}
\end{attention}

\begin{savaistu}[Historical Note]
The method of exhaustion used by Archimedes (c. 287–212 BC) to find areas and volumes is the ancient precursor of Riemann sums. The trapezium rule was known to the Babylonians (c. 50 BC) for astronomical calculations. The Mean Value Theorem for integrals was formalised by Cauchy in the 19th century, building on the work of Lagrange and Rolle. In Cameroon, integration techniques are essential in engineering (hydroelectric dam design on the Sanaga River), physics (computing work done by variable forces), and economics (consumer surplus calculations).
\end{savaistu}

\begin{exercice}[Practice Exercises]
\begin{enumerate}
\item Express $\displaystyle\lim_{n\to\infty} \frac{1}{n}\sum_{i=1}^{n} \left(2 + \frac{3i}{n}\right)^2$ as a definite integral and evaluate it.
\item Find the area bounded by $y = x^3 - 4x$ and the $x$-axis.
\item Find the area of the region bounded by $x = y^2 - 4$ and the $y$-axis.
\item The region bounded by $y = \sin x$, the $x$-axis, $x=0$, $x=\pi$ is rotated about the $x$-axis. Find the volume.
\item The region bounded by $y = x^2$, $y = 0$, $x = 2$ is rotated about the $y$-axis. Find the volume.
\item Find the centroid of the region bounded by $y = \sqrt{x}$, $y = 0$, $x = 4$.
\item Use the trapezium rule with 5 strips to estimate $\int_1^2 \ln x\,dx$. (Use $\ln 1=0$, $\ln 1.2=0.1823$, $\ln 1.4=0.3365$, $\ln 1.6=0.4700$, $\ln 1.8=0.5878$, $\ln 2=0.6931$.)
\item Find the mean value of $f(x) = \cos x$ on $[0, \pi/2]$ and the corresponding $c$.
\end{enumerate}
\end{exercice}

\begin{corrige}[Exercise 1]
$\frac{1}{n}\sum (2+3i/n)^2 = \frac{1}{n}\sum f(i/n)$ with $f(x) = (2+3x)^2$ on $[0,1]$. Integral: $\int_0^1 (2+3x)^2\,dx = \int_0^1 (4+12x+9x^2)\,dx = [4x+6x^2+3x^3]_0^1 = 4+6+3 = 13$.
\end{corrige}

\begin{corrige}[Exercise 2]
Zeros: $x(x^2-4)=0 \Rightarrow x=-2,0,2$. Area = $\int_{-2}^0 (x^3-4x)\,dx + \int_0^2 -(x^3-4x)\,dx = [x^4/4-2x^2]_{-2}^0 - [x^4/4-2x^2]_0^2 = (0 - (4-8)) - ((4-8)-0) = 4 - (-4) = 8$ sq units.
\end{corrige}

\begin{corrige}[Exercise 3]
$x = y^2-4$ meets $y$-axis at $x=0 \Rightarrow y=\pm 2$. For $y\in[-2,2]$, $x\le 0$, so $|x| = 4-y^2$. $A = \int_{-2}^2 (4-y^2)\,dy = [4y - y^3/3]_{-2}^2 = (8-8/3) - (-8+8/3) = 32/3$ sq units.
\end{corrige}

\begin{corrige}[Exercise 4]
$V = \pi\int_0^\pi \sin^2 x\,dx = \pi\int_0^\pi \frac{1-\cos 2x}{2}\,dx = \frac{\pi}{2}[x - \frac{\sin 2x}{2}]_0^\pi = \frac{\pi^2}{2}$ cubic units.
\end{corrige}

\begin{corrige}[Exercise 5]
$x = \sqrt{y}$, $x^2 = y$, $y\in[0,4]$. $V = \pi\int_0^4 y\,dy = 8\pi$ cubic units. (Or shell method: $V = 2\pi\int_0^2 x(x^2)\,dx = 2\pi[x^4/4]_0^2 = 8\pi$.)
\end{corrige}

\begin{corrige}[Exercise 6]
$A = \int_0^4 \sqrt{x}\,dx = [2/3 x^{3/2}]_0^4 = 16/3$. $M_y = \int_0^4 x\sqrt{x}\,dx = \int_0^4 x^{3/2}\,dx = [2/5 x^{5/2}]_0^4 = 64/5$. $\bar{x} = (64/5)/(16/3) = 12/5 = 2.4$. $M_x = \frac12\int_0^4 x\,dx = \frac12[8] = 4$. $\bar{y} = 4/(16/3) = 3/4 = 0.75$. Centroid: $(2.4, 0.75)$.
\end{corrige}

\begin{corrige}[Exercise 7]
$h = (2-1)/5 = 0.2$. $y_0=0, y_1=0.1823, y_2=0.3365, y_3=0.4700, y_4=0.5878, y_5=0.6931$. Sum = $0 + 0.6931 + 2(0.1823+0.3365+0.4700+0.5878) = 0.6931 + 2(1.5766) = 3.8463$. Integral $\approx 0.2/2 \times 3.8463 = 0.3846$. Exact: $\int_1^2 \ln x\,dx = [x\ln x - x]_1^2 = 2\ln 2 - 2 - (0-1) = 2\ln 2 - 1 \approx 0.3863$. Error $\approx 0.0017$.
\end{corrige}

\begin{corrige}[Exercise 8]
$\bar{f} = \frac{1}{\pi/2}\int_0^{\pi/2} \cos x\,dx = \frac{2}{\pi}[\sin x]_0^{\pi/2} = \frac{2}{\pi}$. Solve $\cos c = 2/\pi \Rightarrow c = \arccos(2/\pi) \approx 0.881 \in (0,\pi/2)$.
\end{corrige}

\newpage

\coursec{Exercises}

\courssub{I check my knowledge}

\begin{exercice}[MCQ on Riemann sums and definite integrals]
Which of the following statements is \textbf{always} true for a function $f$ continuous on $[a,b]$?
\begin{enumerate}
\item The definite integral $\int_a^b f(x)\,dx$ equals the area between the curve $y=f(x)$ and the $x$-axis.
\item $\int_a^b f(x)\,dx = \lim_{n\to\infty} \sum_{i=1}^n f(x_i^*)\Delta x$ where $\Delta x = \frac{b-a}{n}$ and $x_i^*$ is any point in the $i$-th subinterval.
\item If $f(x) \ge 0$ on $[a,b]$, then $\int_a^b f(x)\,dx \ge 0$.
\item $\int_a^b f(x)\,dx = F(b)-F(a)$ where $F$ is any antiderivative of $f$.
\end{enumerate}
\end{exercice}

\begin{exercice}[True/False with justification]
Determine whether each statement is true or false. Justify your answer with a brief explanation or a counterexample.
\begin{enumerate}
\item If $\int_0^2 f(x)\,dx = 5$, then $\int_0^2 3f(x)\,dx = 15$.
\item $\int_{-1}^1 x^3\,dx = 0$ because $x^3$ is an odd function.
\item The area of the region bounded by $y=\sin x$ and the $x$-axis on $[0,2\pi]$ is $\int_0^{2\pi} \sin x\,dx$.
\item The trapezium rule always gives an overestimate for $\int_a^b f(x)\,dx$ when $f''(x) > 0$ on $[a,b]$.
\end{enumerate}
\end{exercice}

\begin{exercice}[Definitions and notations]
\begin{enumerate}
\item Write the definition of the definite integral $\int_a^b f(x)\,dx$ as a limit of Riemann sums.
\item State the Mean Value Theorem for Integrals for a continuous function on $[a,b]$.
\item Give the formula for the volume of the solid obtained by rotating the region under $y=f(x)\ge0$ on $[a,b]$ about the $x$-axis.
\item Write the coordinates $(\bar{x},\bar{y})$ of the centroid of a uniform lamina bounded by $y=f(x)$, $y=0$, $x=a$, $x=b$ (with $f(x)\ge0$).
\end{enumerate}
\end{exercice}

\begin{exercice}[Direct calculations]
Evaluate the following definite integrals exactly.
\begin{enumerate}
\item $\displaystyle\int_0^2 (3x^2 - 4x + 1)\,dx$
\item $\displaystyle\int_1^e \frac{\ln x}{x}\,dx$
\item $\displaystyle\int_0^{\pi/2} \sin^2 x\,dx$
\item $\displaystyle\int_0^1 \frac{dx}{1+x^2}$
\end{enumerate}
\end{exercice}

\courssub{I practise}

\begin{exercice}[Area between curve and $x$-axis]
Find the \textbf{total area} enclosed between the curve $y = x^3 - 4x$ and the $x$-axis. (Hint: the curve crosses the $x$-axis at three points; split the integral accordingly.)
\end{exercice}

\begin{exercice}[Area between two curves]
\begin{enumerate}
\item Find the area of the region bounded by the curves $y = x^2$ and $y = 2 - x^2$.
\item Find the area of the region bounded by the curve $x = y^2 - 2$ and the $y$-axis.
\end{enumerate}
\end{exercice}

\begin{exercice}[Volume of revolution about the $x$-axis]
The region under the curve $y = e^x$ from $x = 0$ to $x = 1$ is rotated through $360^\circ$ about the $x$-axis. Find the exact volume of the solid generated.
\end{exercice}

\begin{exercice}[Volume of revolution about the $y$-axis]
The same region as in the previous exercise (under $y = e^x$, $0 \le x \le 1$) is rotated through $360^\circ$ about the $y$-axis.
\begin{enumerate}
\item Set up the integral for the volume using the method of cylindrical shells.
\item Evaluate the integral exactly.
\end{enumerate}
\end{exercice}

\courssub{I prepare for the GCE}

\begin{exercice}[{Riemann sums, Mean Value Theorem and Trapezium rule from a table] \hfill [16 marks}]
\begin{enumerate}
\item By evaluating a limit of Riemann sums with right endpoints, show from first principles that
\[
\int_0^2 x^2\,dx = \frac{8}{3},
\]
using the formula $1^2+2^2+\cdots+n^2 = \frac{n(n+1)(2n+1)}{6}$. \hfill [6 marks]
\item Find the mean value of $f(x) = \sin x$ on the interval $[0,\pi]$. Determine the value(s) of $c \in [0,\pi]$ guaranteed by the Mean Value Theorem for Integrals. \hfill [5 marks]
\item A function $f$ is given only by the following table of values measured every $0.5$ units:
\[
\begin{array}{c|ccccc}
x & 0 & 0.5 & 1.0 & 1.5 & 2.0 \\ \hline
f(x) & 1.0 & 0.8 & 0.5 & 0.3 & 0.2
\end{array}
\]
Use the trapezium rule to estimate $\int_0^2 f(x)\,dx$. State, with a sketch, whether your answer is an overestimate or an underestimate. \hfill [5 marks]
\end{enumerate}
\end{exercice}

\begin{exercice}[{Trapezium rule for $\pi$ and volumes of revolution] \hfill [15 marks}]
\begin{enumerate}
\item Use the trapezium rule with $5$ ordinates (i.e. $4$ strips) to estimate
\[
\int_0^1 \frac{dx}{1+x^2}
\]
to $4$ decimal places. Given that the exact value is $\pi/4$, deduce an approximation for $\pi$ and comment on the accuracy. \hfill [7 marks]
\item The region under $y = e^x$ from $x = 0$ to $x = 1$ is rotated through $360^\circ$ about the $x$-axis. Find the exact volume generated. \hfill [4 marks]
\item The same region is rotated about the $y$-axis. Set up the integral for the volume using the shell method and evaluate it exactly. \hfill [4 marks]
\end{enumerate}
\end{exercice}

\begin{exercice}[{Synthesis: Area, volume, mean value, centroid and centre of mass] \hfill [20 marks}]
\begin{enumerate}
\item Consider the function $f(x) = 2x - x^2$ on the interval $[0,2]$.
\begin{enumerate}
\item Find the area of the region bounded by the curve $y = f(x)$ and the $x$-axis. \hfill [3 marks]
\item Find the volume of the solid generated when this region is rotated through $360^\circ$ about the $x$-axis. \hfill [4 marks]
\item Find the mean value of $f$ on $[0,2]$. \hfill [2 marks]
\item Find the coordinates of the centroid of this plane region. \hfill [4 marks]
\end{enumerate}
\item The region bounded by the curve $y = x^2$, the $x$-axis and the line $x = 2$ is rotated through $360^\circ$ about the $x$-axis.
\begin{enumerate}
\item Find the volume of the solid formed. \hfill [3 marks]
\item Find the $x$-coordinate of the centre of mass of this solid (assuming uniform density). \hfill [4 marks]
\end{enumerate}
\end{enumerate}
\end{exercice}

\newpage

\coursec{Solutions to the exercises}

\begin{corrige}[Exercise 1 — MCQ on Riemann sums and definite integrals]
\begin{enumerate}
\item \textbf{False.} The definite integral gives the \emph{signed} area: regions below the $x$-axis count negatively. The total geometric area is $\int_a^b |f(x)|\,dx$, which equals $\int_a^b f(x)\,dx$ only when $f(x)\ge 0$ everywhere on $[a,b]$.
\item \textbf{True.} This is precisely the definition of the definite integral for a continuous (hence integrable) function: the limit exists and is independent of the choice of the sample points $x_i^*$.
\item \textbf{True.} Every Riemann sum $\sum f(x_i^*)\Delta x$ is a sum of non-negative terms, so its limit is non-negative. (Equivalently, integration preserves inequalities: $f \ge 0 \Rightarrow \int_a^b f \ge 0$.)
\item \textbf{True.} This is the Fundamental Theorem of Calculus: since $f$ is continuous, any antiderivative $F$ satisfies $\int_a^b f(x)\,dx = F(b)-F(a)$.
\end{enumerate}
\end{corrige}

\begin{corrige}[Exercise 2 — True/False with justification]
\begin{enumerate}
\item \textbf{True.} By linearity of the integral, $\int_0^2 3f(x)\,dx = 3\int_0^2 f(x)\,dx = 3\times 5 = 15$.
\item \textbf{True.} For an odd function, $\int_{-a}^{a} f(x)\,dx = 0$. Check directly: $\int_{-1}^1 x^3\,dx = \left[\frac{x^4}{4}\right]_{-1}^1 = \frac14 - \frac14 = 0$. Geometrically, the area on $[-1,0]$ (negative) cancels the area on $[0,1]$ (positive).
\item \textbf{False.} $\sin x \le 0$ on $[\pi, 2\pi]$, so $\int_0^{2\pi}\sin x\,dx = 0$, which is the \emph{signed} area, not the total area. The correct total area is
\[
\int_0^{2\pi}|\sin x|\,dx = \int_0^{\pi}\sin x\,dx - \int_{\pi}^{2\pi}\sin x\,dx = 2+2 = 4.
\]
\item \textbf{True.} If $f''(x) > 0$, then $f$ is strictly convex, so each chord lies strictly above the curve. Each trapezium therefore has area greater than the area under the arc, and the trapezium rule overestimates the integral.
\end{enumerate}
\end{corrige}

\begin{corrige}[Exercise 3 — Definitions and notations]
\begin{enumerate}
\item Partition $[a,b]$ into $n$ subintervals of equal width $\Delta x = \dfrac{b-a}{n}$, with endpoints $x_i = a + i\Delta x$, and choose any sample point $x_i^* \in [x_{i-1}, x_i]$. Then
\[
\int_a^b f(x)\,dx = \lim_{n\to\infty}\sum_{i=1}^{n} f(x_i^*)\,\Delta x .
\]
\item If $f$ is continuous on $[a,b]$, there exists at least one $c\in[a,b]$ such that
\[
\int_a^b f(x)\,dx = f(c)\,(b-a), \qquad\text{i.e. } f(c) = \frac{1}{b-a}\int_a^b f(x)\,dx .
\]
The number $\dfrac{1}{b-a}\displaystyle\int_a^b f(x)\,dx$ is the \cle{mean value} of $f$ on $[a,b]$.
\item Disc method: $\displaystyle V = \pi\int_a^b [f(x)]^2\,dx$.
\item With $A = \displaystyle\int_a^b f(x)\,dx$:
\[
\bar{x} = \frac{1}{A}\int_a^b x\,f(x)\,dx, \qquad \bar{y} = \frac{1}{A}\int_a^b \tfrac12 [f(x)]^2\,dx .
\]
\end{enumerate}
\end{corrige}

\begin{corrige}[Exercise 4 — Direct calculations]
\begin{enumerate}
\item
\[
\int_0^2 (3x^2-4x+1)\,dx = \Big[x^3 - 2x^2 + x\Big]_0^2 = (8 - 8 + 2) - 0 = \boxed{2}.
\]
\item Substitute $u = \ln x$, $du = \frac{dx}{x}$; when $x=1$, $u=0$; when $x=e$, $u=1$:
\[
\int_1^e \frac{\ln x}{x}\,dx = \int_0^1 u\,du = \left[\frac{u^2}{2}\right]_0^1 = \boxed{\tfrac12}.
\]
\item Use $\sin^2 x = \frac{1-\cos 2x}{2}$:
\[
\int_0^{\pi/2}\sin^2 x\,dx = \left[\frac{x}{2} - \frac{\sin 2x}{4}\right]_0^{\pi/2} = \frac{\pi}{4} - 0 = \boxed{\frac{\pi}{4}}.
\]
\item
\[
\int_0^1 \frac{dx}{1+x^2} = \Big[\arctan x\Big]_0^1 = \arctan 1 - \arctan 0 = \boxed{\frac{\pi}{4}}.
\]
\end{enumerate}
\end{corrige}

\begin{corrige}[Exercise 5 — Area between curve and $x$-axis]
The curve $y = x^3 - 4x = x(x-2)(x+2)$ crosses the $x$-axis at $x = -2, 0, 2$.

\textbf{Sign analysis.} On $[-2,0]$: test $x=-1$: $y = -1+4 = 3 > 0$. On $[0,2]$: test $x=1$: $y = 1 - 4 = -3 < 0$.

Hence the total area is
\[
A = \int_{-2}^{0}(x^3-4x)\,dx - \int_{0}^{2}(x^3-4x)\,dx .
\]
An antiderivative is $F(x) = \dfrac{x^4}{4} - 2x^2$.
\[
\int_{-2}^{0}(x^3-4x)\,dx = F(0) - F(-2) = 0 - (4 - 8) = 4,
\]
\[
\int_{0}^{2}(x^3-4x)\,dx = F(2) - F(0) = (4-8) - 0 = -4.
\]
Therefore
\[
A = 4 - (-4) = \boxed{8 \text{ square units}}.
\]

\textbf{Warning (common trap).} Writing simply $\int_{-2}^{2}(x^3-4x)\,dx = 0$ gives the \emph{signed} area (the function is odd), not the total area. Always split the integral at the zeros of $f$ and take the absolute value of each piece.
\end{corrige}

\begin{corrige}[Exercise 6 — Area between two curves]
\begin{enumerate}
\item \textbf{Intersection points:} $x^2 = 2 - x^2 \Rightarrow 2x^2 = 2 \Rightarrow x = \pm 1$.

On $[-1,1]$, test $x=0$: $2 - x^2 = 2 > 0 = x^2$, so $y = 2-x^2$ is the upper curve.
\[
A = \int_{-1}^{1}\big[(2-x^2) - x^2\big]\,dx = \int_{-1}^{1}(2 - 2x^2)\,dx = \left[2x - \frac{2x^3}{3}\right]_{-1}^{1}.
\]
\[
A = \left(2 - \frac{2}{3}\right) - \left(-2 + \frac{2}{3}\right) = \frac{4}{3} + \frac{4}{3} = \boxed{\frac{8}{3} \text{ square units}}.
\]

\item \textbf{Intersection with the $y$-axis:} $y^2 - 2 = 0 \Rightarrow y = \pm\sqrt{2}$.

Here we integrate with respect to $y$. The curve $x = y^2 - 2$ lies to the \emph{left} of the $y$-axis ($x \le 0$) for $-\sqrt2 \le y \le \sqrt2$, so the area is
\[
A = \int_{-\sqrt2}^{\sqrt2}\big[0 - (y^2 - 2)\big]\,dy = \int_{-\sqrt2}^{\sqrt2}(2 - y^2)\,dy = \left[2y - \frac{y^3}{3}\right]_{-\sqrt2}^{\sqrt2}.
\]
\[
A = \left(2\sqrt2 - \frac{2\sqrt2}{3}\right) - \left(-2\sqrt2 + \frac{2\sqrt2}{3}\right) = \frac{4\sqrt2}{3} + \frac{4\sqrt2}{3} = \boxed{\frac{8\sqrt2}{3} \text{ square units}}.
\]
\end{enumerate}
\end{corrige}

\begin{corrige}[Exercise 7 — Volume of revolution about the $x$-axis]
By the disc method with $y = e^x$ on $[0,1]$:
\[
V = \pi\int_0^1 (e^x)^2\,dx = \pi\int_0^1 e^{2x}\,dx = \pi\left[\frac{e^{2x}}{2}\right]_0^1 = \pi\left(\frac{e^2}{2} - \frac{1}{2}\right).
\]
\[
\boxed{V = \frac{\pi}{2}\left(e^2 - 1\right) \text{ cubic units}} \approx 10.04 \text{ cubic units}.
\]
\end{corrige}

\begin{corrige}[Exercise 8 — Volume of revolution about the $y$-axis]
\begin{enumerate}
\item \textbf{Shell method.} A thin vertical strip at position $x$, of width $dx$ and height $e^x$, rotated about the $y$-axis generates a cylindrical shell of radius $x$, height $e^x$ and thickness $dx$, hence of volume $dV = 2\pi x\, e^x\, dx$. Therefore
\[
V = 2\pi\int_0^1 x\,e^x\,dx .
\]
\item \textbf{Evaluation by parts.} Let $u = x$, $dv = e^x dx$, so $du = dx$, $v = e^x$:
\[
\int_0^1 x e^x\,dx = \Big[x e^x\Big]_0^1 - \int_0^1 e^x\,dx = (e - 0) - \Big[e^x\Big]_0^1 = e - (e - 1) = 1.
\]
Hence
\[
\boxed{V = 2\pi \text{ cubic units}}.
\]
\end{enumerate}

\textbf{Warning (common trap).} Do \emph{not} write $V = \pi\int_0^1 x^2\,dy$ blindly: rotation about the $y$-axis with shells uses $2\pi\int x\,y\,dx$; with discs (washers) one would need to solve $x = \ln y$ and subtract the inner radius, giving $V = \pi\int_1^e \big[1 - (\ln y)^2\big]\,dy$, which also yields $2\pi$ after more work.
\end{corrige}

\begin{corrige}[Exercise 9 — Riemann sums, Mean Value Theorem and Trapezium rule]
\begin{enumerate}
\item \textbf{From first principles.} Partition $[0,2]$ into $n$ equal strips: $\Delta x = \frac{2}{n}$, right endpoints $x_i = \frac{2i}{n}$, $i = 1, \dots, n$. The Riemann sum is
\[
S_n = \sum_{i=1}^{n} f(x_i)\,\Delta x = \sum_{i=1}^{n}\left(\frac{2i}{n}\right)^2 \cdot \frac{2}{n} = \frac{8}{n^3}\sum_{i=1}^{n} i^2 = \frac{8}{n^3}\cdot\frac{n(n+1)(2n+1)}{6}.
\]
Simplify:
\[
S_n = \frac{4}{3}\cdot\frac{(n+1)(2n+1)}{n^2} = \frac{4}{3}\left(1+\frac1n\right)\left(2+\frac1n\right).
\]
Taking the limit:
\[
\int_0^2 x^2\,dx = \lim_{n\to\infty} S_n = \frac{4}{3}(1)(2) = \frac{8}{3}. \qquad\blacksquare
\]
\textit{Mark scheme:} correct $\Delta x$ and $x_i$ \textbf{[1]}; Riemann sum set up \textbf{[1]}; use of $\sum i^2$ formula \textbf{[1]}; algebraic simplification \textbf{[1]}; limit taken correctly \textbf{[1]}; conclusion $\frac83$ \textbf{[1]}.

\item \textbf{Mean value of $\sin x$ on $[0,\pi]$.}
\[
\bar{f} = \frac{1}{\pi - 0}\int_0^{\pi}\sin x\,dx = \frac{1}{\pi}\Big[-\cos x\Big]_0^{\pi} = \frac{1}{\pi}(1+1) = \frac{2}{\pi}.
\]
By the Mean Value Theorem for Integrals, there exists $c\in[0,\pi]$ with $\sin c = \frac{2}{\pi}$. Since $\frac{2}{\pi}\approx 0.6366 < 1$, solutions exist:
\[
c = \arcsin\frac{2}{\pi} \approx 0.6901 \quad\text{or}\quad c = \pi - \arcsin\frac{2}{\pi} \approx 2.4515.
\]
\textit{Mark scheme:} mean value formula \textbf{[1]}; integral $=2$ \textbf{[1]}; $\bar f = 2/\pi$ \textbf{[1]}; equation $\sin c = 2/\pi$ \textbf{[1]}; both values of $c$ \textbf{[1]}.

\item \textbf{Trapezium rule} with $h = 0.5$, ordinates $y_0 = 1.0,\ y_1 = 0.8,\ y_2 = 0.5,\ y_3 = 0.3,\ y_4 = 0.2$:
\[
T = \frac{h}{2}\Big[y_0 + y_4 + 2(y_1 + y_2 + y_3)\Big] = \frac{0.5}{2}\big[1.0 + 0.2 + 2(0.8+0.5+0.3)\big].
\]
\[
T = 0.25\,[1.2 + 2(1.6)] = 0.25 \times 4.4 = \boxed{1.1}.
\]
\textbf{Over/underestimate:} the tabulated values decrease and appear \emph{convex} (the successive drops $0.2, 0.3, 0.2, 0.1$ suggest the curve bends upwards, like a decaying exponential). For a convex curve each chord lies above the arc, so the trapezia cover more than the true area: the estimate $1.1$ is an \textbf{overestimate}.

\begin{popfigure}
\begin{center}
\begin{tikzpicture}[scale=2.2]
\draw[->] (-0.2,0) -- (2.4,0) node[right] {$x$};
\draw[->] (0,-0.1) -- (0,1.3) node[above] {$y$};
\draw[popblue, thick, domain=0:2, samples=60] plot (\x, {1/(1+1.2*\x)});
\draw[poporange, thick] (0,1) -- (0.5,0.8) -- (1,0.5) -- (1.5,0.3) -- (2,0.2);
\foreach \x/\y in {0/1, 0.5/0.8, 1/0.5, 1.5/0.3, 2/0.2} {
  \draw[dashed, gray] (\x,0) -- (\x,\y);
  \fill[popdark] (\x,\y) circle (0.015);
}
\node[popblue] at (1.5,0.85) {curve};
\node[poporange] at (0.55,1.12) {chords};
\end{tikzpicture}
\end{center}
\legende{For a convex curve the chords (orange) lie above the curve (blue): the trapezium rule overestimates.}
\end{popfigure}
\textit{Mark scheme:} formula with $h=0.5$ \textbf{[1]}; correct ordinate sum \textbf{[1]}; answer $1.1$ \textbf{[1]}; sketch with chords above curve \textbf{[1]}; conclusion "overestimate" justified by convexity \textbf{[1]}.
\end{enumerate}
\end{corrige}

\begin{corrige}[Exercise 10 — Trapezium rule for $\pi$ and volumes of revolution]
\begin{enumerate}
\item With $5$ ordinates on $[0,1]$: $h = \frac{1-0}{4} = 0.25$. Compute $f(x) = \dfrac{1}{1+x^2}$:
\begin{center}
\begin{tabular}{|c|ccccc|}
\hline
\rowcolor{popblueL} $x$ & $0$ & $0.25$ & $0.5$ & $0.75$ & $1$ \\
\hline
$f(x)$ & $1$ & $0.9412$ & $0.8$ & $0.64$ & $0.5$ \\
\hline
\end{tabular}
\end{center}
(Values: $\frac{1}{1.0625} = 0.94118$, $\frac{1}{1.25} = 0.8$, $\frac{1}{1.5625} = 0.64$.)

Trapezium rule:
\[
T = \frac{0.25}{2}\Big[1 + 0.5 + 2(0.9412 + 0.8 + 0.64)\Big] = 0.125\,[1.5 + 2(2.3812)].
\]
\[
T = 0.125 \times 6.2624 = 0.7828.
\]
So $\displaystyle\int_0^1\frac{dx}{1+x^2} \approx 0.7828$, and since the exact value is $\frac{\pi}{4}$:
\[
\pi \approx 4 \times 0.7828 = \boxed{3.1312}.
\]
\textbf{Comment on accuracy.} The true value is $\pi \approx 3.1416$, so the error is about $0.0104$, i.e. roughly $0.33\%$: the approximation is correct to about $2$ decimal places. Here $T = 0.7828 < \frac{\pi}{4} = 0.7854$, so the trapezium rule has \emph{underestimated} the integral. This is consistent with the shape of the curve: $f''(x) = \dfrac{6x^2 - 2}{(1+x^2)^3}$, so $f$ is concave on $\left[0, \frac{1}{\sqrt3}\right)$ (where $f'' < 0$, chords below the curve) and convex only on $\left(\frac{1}{\sqrt3}, 1\right]$; the concave part dominates and the chords lie, on balance, below the arc. Finally, since the error of the trapezium rule behaves like $h^2$, doubling the number of strips would roughly \emph{quarter} the error.

\textit{Mark scheme:} $h = 0.25$ \textbf{[1]}; table of ordinates correct \textbf{[2]}; trapezium formula applied \textbf{[1]}; $T = 0.7828$ \textbf{[1]}; $\pi \approx 3.1312$ \textbf{[1]}; sensible comment on accuracy/error \textbf{[1]}.

\item \textbf{Volume about the $x$-axis} (disc method):
\[
V = \pi\int_0^1 e^{2x}\,dx = \pi\left[\frac{e^{2x}}{2}\right]_0^1 = \boxed{\frac{\pi}{2}\left(e^2 - 1\right)}.
\]
\textit{Mark scheme:} correct formula $\pi\int y^2 dx$ \textbf{[1]}; $e^{2x}$ integrated \textbf{[1]}; limits substituted \textbf{[1]}; exact answer \textbf{[1]}.

\item \textbf{Volume about the $y$-axis} (shell method): $V = 2\pi\displaystyle\int_0^1 x e^x\,dx$. Integration by parts:
\[
\int_0^1 x e^x\,dx = \Big[xe^x\Big]_0^1 - \Big[e^x\Big]_0^1 = e - (e-1) = 1,
\]
so $\boxed{V = 2\pi}$ cubic units.

\textit{Mark scheme:} shell formula $2\pi\int xy\,dx$ \textbf{[1]}; by parts set up \textbf{[1]}; integral $=1$ \textbf{[1]}; $V = 2\pi$ \textbf{[1]}.
\end{enumerate}
\end{corrige}

\begin{corrige}[Exercise 11 — Synthesis: area, volume, mean value, centroid and centre of mass]
\textbf{Part 1.} $f(x) = 2x - x^2 = x(2-x) \ge 0$ on $[0,2]$ (a downward parabola with zeros at $0$ and $2$, vertex $(1,1)$).

\begin{enumerate}
\item \textbf{Area.}
\[
A = \int_0^2 (2x - x^2)\,dx = \left[x^2 - \frac{x^3}{3}\right]_0^2 = 4 - \frac{8}{3} = \boxed{\frac{4}{3} \text{ sq. units}}.
\]
\textit{[3 marks: integral set up \textbf{[1]}, antiderivative \textbf{[1]}, answer \textbf{[1]}.]}

\item \textbf{Volume about the $x$-axis.}
\[
V = \pi\int_0^2 (2x - x^2)^2\,dx = \pi\int_0^2 (4x^2 - 4x^3 + x^4)\,dx
\]
\[
= \pi\left[\frac{4x^3}{3} - x^4 + \frac{x^5}{5}\right]_0^2 = \pi\left(\frac{32}{3} - 16 + \frac{32}{5}\right).
\]
Common denominator $15$: $\dfrac{160}{15} - \dfrac{240}{15} + \dfrac{96}{15} = \dfrac{16}{15}$. Hence
\[
\boxed{V = \frac{16\pi}{15} \text{ cubic units}}.
\]
\textit{[4 marks: $\pi\int y^2 dx$ \textbf{[1]}, expansion \textbf{[1]}, integration \textbf{[1]}, exact answer \textbf{[1]}.]}

\item \textbf{Mean value.}
\[
\bar{f} = \frac{1}{2-0}\int_0^2 (2x-x^2)\,dx = \frac{1}{2}\cdot\frac{4}{3} = \boxed{\frac{2}{3}}.
\]
\textit{[2 marks: formula \textbf{[1]}, answer \textbf{[1]}.]}

\item \textbf{Centroid.} By symmetry of the parabola about $x = 1$, $\bar{x} = 1$. Verify:
\[
\int_0^2 x(2x - x^2)\,dx = \int_0^2 (2x^2 - x^3)\,dx = \left[\frac{2x^3}{3} - \frac{x^4}{4}\right]_0^2 = \frac{16}{3} - 4 = \frac{4}{3},
\]
so $\bar{x} = \dfrac{4/3}{4/3} = 1$. \checkmark

For $\bar{y}$:
\[
\int_0^2 \tfrac12(2x-x^2)^2\,dx = \frac{1}{2}\cdot\frac{16}{15} = \frac{8}{15} \quad (\text{from part (b), without the } \pi),
\]
\[
\bar{y} = \frac{1}{A}\cdot\frac{8}{15} = \frac{3}{4}\cdot\frac{8}{15} = \frac{2}{5}.
\]
\[
\boxed{(\bar{x}, \bar{y}) = \left(1,\ \frac{2}{5}\right)}.
\]
\textit{[4 marks: $\bar x$ formula or symmetry \textbf{[1]}, $\bar x = 1$ \textbf{[1]}, $\bar y$ formula with $\tfrac12 y^2$ \textbf{[1]}, $\bar y = \frac25$ \textbf{[1]}.]}
\end{enumerate}

\textbf{Part 2.} Region under $y = x^2$, $0 \le x \le 2$, rotated about the $x$-axis.
\begin{enumerate}
\item \textbf{Volume.}
\[
V = \pi\int_0^2 (x^2)^2\,dx = \pi\int_0^2 x^4\,dx = \pi\left[\frac{x^5}{5}\right]_0^2 = \boxed{\frac{32\pi}{5} \text{ cubic units}}.
\]
\textit{[3 marks: formula \textbf{[1]}, integration \textbf{[1]}, answer \textbf{[1]}.]}

\item \textbf{Centre of mass of the solid.} By symmetry the centre of mass lies on the $x$-axis. For a solid of revolution of uniform density,
\[
\bar{x} = \frac{\displaystyle\pi\int_0^2 x\,y^2\,dx}{\displaystyle\pi\int_0^2 y^2\,dx} = \frac{\displaystyle\int_0^2 x\cdot x^4\,dx}{32/5}.
\]
Numerator:
\[
\int_0^2 x^5\,dx = \left[\frac{x^6}{6}\right]_0^2 = \frac{64}{6} = \frac{32}{3}.
\]
Therefore
\[
\bar{x} = \frac{32/3}{32/5} = \frac{5}{3}.
\]
\[
\boxed{\bar{x} = \frac{5}{3}} \quad\text{(the centre of mass is at } \left(\tfrac53, 0\right)\text{).}
\]
This is sensible: the solid is a "trumpet" widening towards $x = 2$, so more mass lies to the right, and indeed $\frac53 > 1$.

\textit{[4 marks: correct formula $\bar x = \pi\int xy^2 dx / V$ \textbf{[1]}, numerator $\frac{32}{3}$ \textbf{[1]}, division by $V$ \textbf{[1]}, $\bar x = \frac53$ \textbf{[1]}.]}
\end{enumerate}

\textbf{Examiner's expectations for this synthesis question.}
\begin{itemize}
\item Quote each formula \emph{before} substituting: $A = \int y\,dx$, $V = \pi\int y^2 dx$, $\bar f = \frac{1}{b-a}\int f$, $\bar y = \frac{1}{A}\int \frac12 y^2 dx$, $\bar x_{\text{solid}} = \dfrac{\pi\int x y^2 dx}{V}$.
\item Keep answers \textbf{exact} (fractions and $\pi$) unless a decimal is requested.
\item A common trap: forgetting the factor $\tfrac12$ in the $\bar y$ formula, or using $\int x y^2 dx$ with $y$ instead of $y^2$.
\end{itemize}
\end{corrige}

\newpage

\coursec{Summary sheet}

\begin{popfigure}
\begin{center}\tcbox[colback=popink,colframe=popink,coltext=white,arc=9pt,boxsep=4pt,left=6pt,right=6pt]{\bfseries\small Integration (PMS04)}\end{center}
\vspace{-2pt}
\begin{tcbraster}[raster columns=3,raster equal height=rows,raster column skip=6pt,raster row skip=6pt,enhanced,blankest]
\begin{tcolorbox}[enhanced,colback=white,colframe=popblue,boxrule=0.9pt,arc=3pt,title={Riemann sums},fonttitle=\bfseries\scriptsize,coltitle=white,colbacktitle=popblue,left=4pt,right=4pt,top=2pt,bottom=2pt,boxsep=1pt,toptitle=1.5pt,bottomtitle=1.5pt,halign=left]
\begin{itemize}[nosep,leftmargin=*,itemsep=1pt,font=\scriptsize]
\item $\sum f(x_i)\Delta x$
\item $\int_a^b f\,dx$
\end{itemize}
\end{tcolorbox}
\begin{tcolorbox}[enhanced,colback=white,colframe=popgreen,boxrule=0.9pt,arc=3pt,title={Areas},fonttitle=\bfseries\scriptsize,coltitle=white,colbacktitle=popgreen,left=4pt,right=4pt,top=2pt,bottom=2pt,boxsep=1pt,toptitle=1.5pt,bottomtitle=1.5pt,halign=left]
\begin{itemize}[nosep,leftmargin=*,itemsep=1pt,font=\scriptsize]
\item $A=\int_a^b f\,dx$
\item $A=\int_c^d x\,dy$
\item Between curves
\end{itemize}
\end{tcolorbox}
\begin{tcolorbox}[enhanced,colback=white,colframe=poporange,boxrule=0.9pt,arc=3pt,title={Volumes},fonttitle=\bfseries\scriptsize,coltitle=white,colbacktitle=poporange,left=4pt,right=4pt,top=2pt,bottom=2pt,boxsep=1pt,toptitle=1.5pt,bottomtitle=1.5pt,halign=left]
\begin{itemize}[nosep,leftmargin=*,itemsep=1pt,font=\scriptsize]
\item $V=\pi\int y^2 dx$
\item $V=\pi\int x^2 dy$
\end{itemize}
\end{tcolorbox}
\begin{tcolorbox}[enhanced,colback=white,colframe=popgold,boxrule=0.9pt,arc=3pt,title={Centroids},fonttitle=\bfseries\scriptsize,coltitle=white,colbacktitle=popgold,left=4pt,right=4pt,top=2pt,bottom=2pt,boxsep=1pt,toptitle=1.5pt,bottomtitle=1.5pt,halign=left]
\begin{itemize}[nosep,leftmargin=*,itemsep=1pt,font=\scriptsize]
\item $\bar{x}=\frac{\int x f\,dx}{\int f\,dx}$
\item Pappus
\end{itemize}
\end{tcolorbox}
\begin{tcolorbox}[enhanced,colback=white,colframe=poppurple,boxrule=0.9pt,arc=3pt,title={Trapezium rule},fonttitle=\bfseries\scriptsize,coltitle=white,colbacktitle=poppurple,left=4pt,right=4pt,top=2pt,bottom=2pt,boxsep=1pt,toptitle=1.5pt,bottomtitle=1.5pt,halign=left]
\begin{itemize}[nosep,leftmargin=*,itemsep=1pt,font=\scriptsize]
\item $\frac{h}{2}[y_0+2(\cdots)+y_n]$
\end{itemize}
\end{tcolorbox}
\begin{tcolorbox}[enhanced,colback=white,colframe=poppink,boxrule=0.9pt,arc=3pt,title={Mean value},fonttitle=\bfseries\scriptsize,coltitle=white,colbacktitle=poppink,left=4pt,right=4pt,top=2pt,bottom=2pt,boxsep=1pt,toptitle=1.5pt,bottomtitle=1.5pt,halign=left]
\begin{itemize}[nosep,leftmargin=*,itemsep=1pt,font=\scriptsize]
\item $f(c)=\frac{1}{b-a}\int_a^b f$
\end{itemize}
\end{tcolorbox}
\end{tcbraster}

\legende{Mind map of the chapter: from Riemann sums to applications of the definite integral.}
\end{popfigure}

\begin{bilan}
\textbf{Key definitions}
\begin{itemize}
\item \cle{Definite integral} as a limit of Riemann sums: if $f$ is continuous on $[a,b]$, then
\[ \int_a^b f(x)\,dx=\lim_{n\to\infty}\sum_{i=1}^{n} f(x_i^{*})\,\Delta x,\qquad \Delta x=\frac{b-a}{n}. \]
\item \cle{Fundamental Theorem of Calculus}: if $F'=f$, then $\displaystyle\int_a^b f(x)\,dx=F(b)-F(a)=\bigl[F(x)\bigr]_a^b$.
\item \cle{Mean value} of $f$ on $[a,b]$: $\displaystyle \bar{f}=\frac{1}{b-a}\int_a^b f(x)\,dx$.
\end{itemize}

\textbf{Formulas to know by heart}
\begin{itemize}
\item Area under a curve (above the $x$-axis): $\displaystyle A=\int_a^b f(x)\,dx$; if $f\le 0$, $A=-\displaystyle\int_a^b f(x)\,dx=\left|\int_a^b f(x)\,dx\right|$.
\item Area between two curves ($f\ge g$ on $[a,b]$): $\displaystyle A=\int_a^b\bigl[f(x)-g(x)\bigr]dx$.
\item Area bounded by the $y$-axis: $\displaystyle A=\int_c^d x\,dy=\int_c^d g(y)\,dy$ where $x=g(y)$.
\item Volume about the $x$-axis: $\displaystyle V=\pi\int_a^b y^2\,dx$; about the $y$-axis: $\displaystyle V=\pi\int_c^d x^2\,dy$ (disc method). Washer: $V=\pi\int_a^b(y_1^2-y_2^2)\,dx$.
\item Centroid of the region under $y=f(x)$, $a\le x\le b$:
\[ \bar{x}=\frac{\displaystyle\int_a^b x f(x)\,dx}{\displaystyle\int_a^b f(x)\,dx},\qquad
\bar{y}=\frac{\displaystyle\tfrac12\int_a^b [f(x)]^2\,dx}{\displaystyle\int_a^b f(x)\,dx}. \]
\item Centre of mass of a lamina with density $\rho$: replace area by mass $M=\int \rho\,dA$; moments $M_y=\int x\,dm$, $M_x=\int y\,dm$, then $\bar{x}=M_y/M$, $\bar{y}=M_x/M$.
\item \cle{Trapezium rule} with $n$ strips of width $h=\dfrac{b-a}{n}$:
\[ \int_a^b f(x)\,dx\approx \frac{h}{2}\Bigl[y_0+2(y_1+y_2+\cdots+y_{n-1})+y_n\Bigr]. \]
\item \cle{Mean Value Theorem for integrals}: if $f$ is continuous on $[a,b]$, there exists $c\in[a,b]$ with
\[ \int_a^b f(x)\,dx=(b-a)\,f(c). \]
\end{itemize}

\textbf{Essential methods}
\begin{enumerate}
\item \textbf{Area:} sketch the curve, find intersections with the axes (solve $f(x)=0$), split the integral where $f$ changes sign, integrate each part and add the \emph{absolute values}.
\item \textbf{Volume:} identify the axis of rotation, square the correct radius ($y$ for the $x$-axis, $x$ for the $y$-axis), multiply by $\pi$, integrate.
\item \textbf{Centroid:} compute the area (denominator) first, then the two moments; check that $(\bar{x},\bar{y})$ lies inside the region.
\item \textbf{Trapezium rule:} build a table of $x_i$ and $y_i=f(x_i)$ with equal step $h$; keep at least 4 decimal places; the first and last ordinates are counted once, all others twice.
\end{enumerate}

\textbf{Common mistakes to avoid}
\begin{itemize}
\item Taking $\displaystyle\int_a^b f(x)\,dx$ as an area when $f$ is negative: an area is always \emph{positive}.
\item Forgetting the factor $\pi$ in volumes of revolution, or squaring $x$ instead of $y$.
\item Mixing the axes: for rotation about the $y$-axis, express $x$ in terms of $y$ and integrate with respect to $y$.
\item In the trapezium rule, using the wrong $h$ (it is $\frac{b-a}{n}$, not $\frac{b-a}{n-1}$) or doubling $y_0$ and $y_n$.
\item Dropping the denominator (total area or mass) in centroid calculations.
\item Confusing the mean value $\bar f$ with $f\!\left(\tfrac{a+b}{2}\right)$: they are equal only in special cases.
\end{itemize}

\textbf{GCE tips}
\begin{itemize}
\item Always \emph{sketch} the region first: most marks are lost through wrong limits, not wrong integration.
\item State units: areas in square units, volumes in cubic units; leave answers exact (in terms of $\pi$) unless a decimal is requested.
\item For the trapezium rule, show the ordinate table clearly --- method marks are awarded for it.
\item Check plausibility: compare a trapezium estimate with a quick sketch; a volume should be positive and of reasonable size.
\item When a question says ``hence'', use the previous result (e.g.\ use a computed area to find a centroid or a mean value).
\end{itemize}
\end{bilan}

\findecours

\end{document}
